% Dérivation — Mathématiques 1re Tech — Cours signé OnBuch+
% Programme officiel MINESEC. Compiler avec : tectonic N14-derivation.tex
\newcommand{\DOCMATIERE}{Mathématiques}
\newcommand{\DOCNIVEAU}{1re Tech}
\newcommand{\DOCMODULE}{Mathématiques – classe de Première technique}
\newcommand{\DOCLECON}{N14}
\newcommand{\DOCTITRE}{Dérivation}
% OnBuch+ — préambule « cours signés » ENRICHI (compilé avec tectonic / XeLaTeX)
% Système visuel « Pop » d'OnBuch+ : fond crème (jamais blanc), bordures encre
% épaisses, ombres « dures » décalées, titres Archivo Black en capitales.
% Version MATHÉMATIQUES : graphes (pgfplots/tikz), boîtes pédagogiques
% (définition, propriété, méthode, attention, le savais-tu, exercice résolu,
% exercice, TP, bilan), page de garde et signature OnBuch+ ancrée en bas.
%
% Macros attendues dans main.tex AVANT \input{preamble} :
%   \DOCMATIERE  \DOCNIVEAU  \DOCMODULE  \DOCLECON (numéro)  \DOCTITRE
\documentclass[11pt]{article}

\usepackage{fontspec}
\usepackage[french]{babel}
\usepackage[a4paper,margin=2cm,top=3.4cm,bottom=2.8cm,headheight=1.4cm,footskip=1.3cm]{geometry}
\usepackage{amsmath,amssymb,amsfonts}
\usepackage{mathtools} % \xmapsto, \coloneqq... souvent utilisés par les modèles
\usepackage{cancel} % simplifications barrées
% \abs{z} (module, valeur absolue) et \norm{u} (norme), utilisés par les modèles
\providecommand{\abs}{}\let\abs\relax\DeclarePairedDelimiter{\abs}{\lvert}{\rvert}
\providecommand{\norm}{}\let\norm\relax\DeclarePairedDelimiter{\norm}{\lVert}{\rVert}
\usepackage[table]{xcolor} % [table] : fournit aussi \rowcolors (alternance de lignes)
\usepackage{pagecolor}
\usepackage{tikz}
\usetikzlibrary{arrows.meta,positioning,calc,shapes.geometric,shapes.misc,shapes.arrows,decorations.pathmorphing,decorations.pathreplacing,decorations.markings,patterns,shadows,mindmap,trees,fit,backgrounds,matrix,intersections,angles,quotes}
\usepackage{pgfplots}
\usepackage{pgf-pie} % diagrammes circulaires \pie (statistiques)
\pgfplotsset{compat=1.18}
\usepgfplotslibrary{fillbetween,groupplots} % aires entre courbes (calcul intégral)
\usepackage{enumitem}
\usepackage{fancyhdr}
% [most] charge toutes les bibliothèques tcolorbox (listings, minted, raster,
% documentation...) et gonfle inutilement la mémoire TeX sur un document long
% (~300 boîtes) : on ne charge que ce qui est réellement utilisé.
\usepackage[skins,breakable]{tcolorbox}
\usepackage{graphicx}
\usepackage{colortbl}
\usepackage{booktabs}
\usepackage{tabularx}
\usepackage{diagbox} % case d'en-tête diagonale des tableaux à double entrée
% Colonnes extensibles centrée / gauche / droite (C, L, R), souvent utilisées par les modèles
\newcolumntype{C}{>{\centering\arraybackslash}X}
\newcolumntype{L}{>{\raggedright\arraybackslash}X}
\newcolumntype{R}{>{\raggedleft\arraybackslash}X}
\usepackage{array}
\usepackage{multicol}
\usepackage{siunitx}
% parse-numbers=false : accepte aussi \SI{2,5 \times 10^{-3}}{...}
\sisetup{locale=FR,output-decimal-marker={,},group-minimum-digits=4,parse-numbers=false}
\DeclareSIUnit\ohm{\ensuremath{\symup{\Omega}}} % U+2126 absent de la police
\usepackage{qrcode}
% \degree : commande fréquemment utilisée par les modèles pour un angle ou une
% température en dehors de \SI{}{} (non fournie par les paquets chargés).
\providecommand{\degree}{^{\circ}}
% \qed (fin de démonstration), utilisé par les modèles sans amsthm
\providecommand{\qed}{\hfill$\square$}
% \barre : double barre des valeurs interdites dans un tableau de variations (tkz-tab)
\providecommand{\barre}{\ensuremath{\|}}
% environnement proof (amsthm non chargé) utilisé par les modèles
\ifdefined\proof\else\newenvironment{proof}[1][Démonstration]{\par\noindent\textit{#1.}\ }{\qed\par}\fi
\usepackage{lastpage}
\usepackage{hyperref}
\hypersetup{hidelinks}

% ============================================================
% Palette « Pop » OnBuch+ — mêmes hex que la classe P (plus_ui.dart)
% ============================================================
\definecolor{poporange}{HTML}{F59321}
\definecolor{popdark}{HTML}{E07A0C}
\definecolor{popcream}{HTML}{FFF6E8}
\definecolor{popink}{HTML}{1C1714}
\definecolor{poppurple}{HTML}{7A5AE0}
\definecolor{popgreen}{HTML}{2E9E6A}
\definecolor{poppink}{HTML}{E0457B}
\definecolor{popblue}{HTML}{2F7DE1}
\definecolor{popgold}{HTML}{C98A12}
\definecolor{popmuted}{HTML}{8A7F76}
\definecolor{popline}{HTML}{EADFD2}
\definecolor{popgray}{HTML}{8A7F76}
\colorlet{popgrey}{popgray}
% Alias tolérants (le modèle nomme parfois les couleurs différemment)
\colorlet{popwhite}{white}
\colorlet{popblack}{popink}
\colorlet{popred}{poppink}
\colorlet{popyellow}{popgold}
% Teintes claires pour fonds de tableaux / zones
\colorlet{popblueL}{popblue!10!white}
\colorlet{poporangeL}{poporange!14!white}
\colorlet{popgreenL}{popgreen!12!white}
\colorlet{poppurpleL}{poppurple!12!white}
\colorlet{poppinkL}{poppink!12!white}
\colorlet{popgoldL}{popgold!14!white}

\pagecolor{popcream}

% ============================================================
% Typographie
% ============================================================
\defaultfontfeatures{Ligatures=TeX}
\setmainfont{PlusJakartaSans-Medium.ttf}[Path=../../../fonts/,BoldFont=PlusJakartaSans-Bold.ttf]
% Maths en sans-serif (Fira Math) avec chiffres et lettres droites de Plus
% Jakarta, pour que les formules se fondent dans le texte.
\usepackage[math-style=ISO,bold-style=ISO]{unicode-math}
\setmathfont{FiraMath-Regular.otf}[Path=../../../fonts/]
\setmathfont{PlusJakartaSans-Medium.ttf}[Path=../../../fonts/,range={up/num,up/{Latin,latin},\mathup}]
\usepackage{icomma} % virgule décimale sans espace en mode math
% Caractères mathématiques Unicode absents des polices de texte (Plus Jakarta,
% Archivo Black) : rendus par leur équivalent mathématique, en texte comme en
% formule — sinon ils s'affichent en carré vide (ex. « Arithmétique dans ℤ »).
\usepackage{newunicodechar}
\newunicodechar{ℝ}{\ensuremath{\mathbb{R}}}
\newunicodechar{ℤ}{\ensuremath{\mathbb{Z}}}
\newunicodechar{ℂ}{\ensuremath{\mathbb{C}}}
\newunicodechar{ℕ}{\ensuremath{\mathbb{N}}}
\newunicodechar{ℚ}{\ensuremath{\mathbb{Q}}}
\newunicodechar{𝔻}{\ensuremath{\mathbb{D}}}
\newunicodechar{α}{\ensuremath{\alpha}}
\newunicodechar{β}{\ensuremath{\beta}}
\newunicodechar{γ}{\ensuremath{\gamma}}
\newunicodechar{δ}{\ensuremath{\delta}}
\newunicodechar{ε}{\ensuremath{\varepsilon}}
\newunicodechar{θ}{\ensuremath{\theta}}
\newunicodechar{λ}{\ensuremath{\lambda}}
\newunicodechar{μ}{\ensuremath{\mu}}
\newunicodechar{σ}{\ensuremath{\sigma}}
\newunicodechar{τ}{\ensuremath{\tau}}
\newunicodechar{φ}{\ensuremath{\varphi}}
\newunicodechar{ψ}{\ensuremath{\psi}}
\newunicodechar{ω}{\ensuremath{\omega}}
\newunicodechar{Σ}{\ensuremath{\Sigma}}
% majuscules produites par \MakeUppercase dans les titres (α-aminés -> Α)
\newunicodechar{Α}{\ensuremath{\alpha}}
\newunicodechar{Β}{\ensuremath{\beta}}
\newunicodechar{Γ}{\ensuremath{\Gamma}}
\newunicodechar{Δ}{\ensuremath{\Delta}}
\newunicodechar{Ε}{\ensuremath{\varepsilon}}
\newunicodechar{Θ}{\ensuremath{\Theta}}
\newunicodechar{Λ}{\ensuremath{\Lambda}}
\newunicodechar{Μ}{\ensuremath{\mu}}
\newunicodechar{Τ}{\ensuremath{\tau}}
\newunicodechar{Φ}{\ensuremath{\Phi}}
\newunicodechar{Ψ}{\ensuremath{\Psi}}
\newunicodechar{Ω}{\ensuremath{\Omega}}
\newunicodechar{↦}{\ensuremath{\mapsto}}
\newunicodechar{∈}{\ensuremath{\in}}
\newunicodechar{∉}{\ensuremath{\notin}}
\newunicodechar{∧}{\ensuremath{\wedge}}
\newunicodechar{⇔}{\ensuremath{\Leftrightarrow}}
\newunicodechar{⟺}{\ensuremath{\Longleftrightarrow}}
\newunicodechar{⇒}{\ensuremath{\Rightarrow}}
\newunicodechar{⟹}{\ensuremath{\Longrightarrow}}
\newunicodechar{∘}{\ensuremath{\circ}}
\newunicodechar{∪}{\ensuremath{\cup}}
\newunicodechar{∩}{\ensuremath{\cap}}
\newunicodechar{≡}{\ensuremath{\equiv}}
\newunicodechar{⊂}{\ensuremath{\subset}}
\newunicodechar{⊥}{\ensuremath{\perp}}
\newunicodechar{∃}{\ensuremath{\exists}}
\newunicodechar{∀}{\ensuremath{\forall}}
\newunicodechar{∛}{\ensuremath{\sqrt[3]{\,}}}
\newunicodechar{∜}{\ensuremath{\sqrt[4]{\,}}}
\newunicodechar{⃗}{\ensuremath{{}^{\to}}}
\newunicodechar{ⁿ}{\ifmmode^{n}\else\textsuperscript{\ensuremath{n}}\fi}
\newunicodechar{ˣ}{\ifmmode^{x}\else\textsuperscript{\ensuremath{x}}\fi}
\newunicodechar{ᵇ}{\ifmmode^{b}\else\textsuperscript{\ensuremath{b}}\fi}
\newunicodechar{ʳ}{\ifmmode^{r}\else\textsuperscript{\ensuremath{r}}\fi}
\newunicodechar{ᵘ}{\ifmmode^{u}\else\textsuperscript{\ensuremath{u}}\fi}
\newunicodechar{ʸ}{\ifmmode^{y}\else\textsuperscript{\ensuremath{y}}\fi}
\newunicodechar{ᵃ}{\ifmmode^{a}\else\textsuperscript{\ensuremath{a}}\fi}
\newunicodechar{ᵉ}{\ifmmode^{e}\else\textsuperscript{\ensuremath{e}}\fi}
\newunicodechar{ᵏ}{\ifmmode^{k}\else\textsuperscript{\ensuremath{k}}\fi}
\newunicodechar{ᵖ}{\ifmmode^{p}\else\textsuperscript{\ensuremath{p}}\fi}
\newunicodechar{ᴺ}{\ifmmode^{N}\else\textsuperscript{\ensuremath{N}}\fi}
\newunicodechar{ᶜ}{\ifmmode^{c}\else\textsuperscript{\ensuremath{c}}\fi}
\newunicodechar{ʰ}{\ifmmode^{h}\else\textsuperscript{\ensuremath{h}}\fi}
\newunicodechar{ᵠ}{\ifmmode^{q}\else\textsuperscript{\ensuremath{q}}\fi}
\newunicodechar{ₙ}{\ifmmode_{n}\else\textsubscript{\ensuremath{n}}\fi}
\newunicodechar{ₐ}{\ifmmode_{a}\else\textsubscript{\ensuremath{a}}\fi}
\newunicodechar{ᵢ}{\ifmmode_{i}\else\textsubscript{\ensuremath{i}}\fi}
\newunicodechar{ᵣ}{\ifmmode_{r}\else\textsubscript{\ensuremath{r}}\fi}
\newunicodechar{ₖ}{\ifmmode_{k}\else\textsubscript{\ensuremath{k}}\fi}
\newunicodechar{ᵦ}{\ifmmode_{\beta}\else\textsubscript{\ensuremath{\beta}}\fi}
\newunicodechar{ₓ}{\ifmmode_{x}\else\textsubscript{\ensuremath{x}}\fi}
\newfontfamily\popdisplay{ArchivoBlack-Regular.ttf}[Path=../../../fonts/]
\newfontfamily\poplogo{SpaceGrotesk-Bold.ttf}[Path=../../../fonts/]
\newfontfamily\popsemi{PlusJakartaSans-SemiBold.ttf}[Path=../../../fonts/]
\newfontfamily\popxbold{PlusJakartaSans-ExtraBold.ttf}[Path=../../../fonts/]
\newfontfamily\popxboldls{PlusJakartaSans-ExtraBold.ttf}[Path=../../../fonts/,LetterSpace=3.0]

\setlist[enumerate]{leftmargin=1.7em,itemsep=5pt,topsep=4pt}
\setlist[itemize]{leftmargin=1.7em,itemsep=4pt,topsep=4pt,label={\color{poporange}\textbullet}}
\setlength{\parindent}{0pt}
\setlength{\parskip}{5pt}
\renewcommand{\arraystretch}{1.25}

% pgfplots : style Pop par défaut
\pgfplotsset{
  every axis/.append style={axis line style={popink,line width=1pt},tick style={popink},
    grid style={popline,line width=0.5pt},label style={font=\small},tick label style={font=\footnotesize}},
  popcurve/.style={poporange,line width=1.8pt,no markers,smooth},
}
% Même style disponible dans un \draw TikZ simple (hors pgfplots)
\tikzset{popcurve/.style={poporange,line width=1.8pt}}

% ============================================================
% Pill / squiggle
% ============================================================
\newcommand{\poppill}[3]{%
  \tikz[baseline=-0.55ex]\node[draw=popink,line width=1.2pt,rounded corners=2.6pt,%
    fill=#2,inner xsep=6.5pt,inner ysep=3pt]{%
    \popxboldls\fontsize{7}{7}\selectfont\color{#3}\MakeUppercase{#1}};%
}
\newcommand{\popsquiggle}[1]{%
  \tikz[baseline=-0.4ex]{
    \draw[#1,line width=2.6pt,line cap=round]
      plot [smooth] coordinates {(0,0.09) (0.34,0.01) (0.68,0.21) (1.02,0.01) (1.36,0.10)};
  }%
}
% Mot-clé surligné (marqueur orange sous le texte)
\newcommand{\cle}[1]{{\popxbold\color{popdark}#1}}

% ============================================================
% En-tête / pied
% ============================================================
\pagestyle{fancy}
\fancyhf{}
\renewcommand{\headrulewidth}{0pt}
\renewcommand{\footrulewidth}{0pt}
\lhead{\raisebox{-0.15ex}{\poplogo\fontsize{13}{13}\selectfont\color{popink}OnBuch\color{poporange}+}}
\rhead{\poppill{\DOCMATIERE\ \textbullet\ \DOCNIVEAU\ \textbullet\ Leçon \DOCLECON}{white}{popink}}
\lfoot{\popxbold\fontsize{7.6}{7.6}\selectfont\color{popmuted}\MakeUppercase{Cours signé OnBuch+}\ \textbullet\ {\popsemi\DOCTITRE}}
\rfoot{\tikz[baseline=-0.6ex]\node[rounded corners=8.5pt,fill=popink,minimum height=17pt,minimum width=17pt,inner xsep=5pt,inner sep=0]{\popxbold\fontsize{8}{8}\selectfont\color{popcream}\thepage\,/\,\pageref*{LastPage}};}

% ============================================================
% Titres
% ============================================================
\newcommand{\chaptitre}[2]{%
  \begin{tcolorbox}[enhanced,colback=poporange,colframe=popink,boxrule=2.6pt,arc=11pt,
      drop shadow=popink,left=15pt,right=15pt,top=12pt,bottom=11pt,before skip=3pt,after skip=18pt]
    \poppill{#2}{popink}{popcream}\par\vspace{9pt}
    {\raggedright\hyphenpenalty=10000\exhyphenpenalty=10000\popdisplay\fontsize{22}{24}\selectfont\color{popcream}\MakeUppercase{#1}\par}
  \end{tcolorbox}
}
% Section numérotée : pastille orange avec le numéro + titre Archivo
\newcounter{popsec}
\newcommand{\coursec}[1]{%
  \stepcounter{popsec}\setcounter{popsub}{0}%
  \par\vspace{16pt}\noindent
  \begin{minipage}{\linewidth}
  \tikz[baseline=-0.7ex]\node[draw=popink,line width=1.6pt,fill=poporange,rounded corners=4pt,minimum size=22pt,inner sep=0]{\popdisplay\fontsize{12}{12}\selectfont\color{popcream}\thepopsec};%
  \hspace{8pt}{\popdisplay\fontsize{15}{17}\selectfont\color{popink}\MakeUppercase{#1}}\par
  \vspace{4pt}\noindent{\color{poporange}\rule{36pt}{3.6pt}}
  \end{minipage}\par\nopagebreak\vspace{8pt}
}
\newcounter{popsub}
\newcommand{\courssub}[1]{%
  \stepcounter{popsub}%
  \par\vspace{9pt}\noindent{\popxbold\fontsize{11.5}{13}\selectfont\color{popink}{\color{poporange}\thepopsec.\thepopsub}\hspace{6pt}#1}\par\nopagebreak\vspace{4pt}
}

% ============================================================
% Boîtes pédagogiques — même recette : fond blanc, bordure épaisse colorée,
% ombre dure encre, badge plein. Titre optionnel [#1].
% ============================================================
\tcbset{popbase/.style={breakable,enhanced,colback=white,boxrule=2.3pt,arc=9pt,
  shadow={3pt}{-3pt}{0pt}{popink},left=14pt,right=14pt,top=11pt,bottom=11pt,before skip=11pt,after skip=15pt,
  fontupper=\popsemi\fontsize{10.6}{14.5}\selectfont\color{popink},
  fontlower=\popsemi\fontsize{10.6}{14.5}\selectfont\color{popink}}}
% Coche dessinée (le glyphe ✓ n'existe pas dans les polices du document)
\newcommand{\popcheck}[1]{\tikz[baseline=-0.5ex]\draw[#1,line width=1.6pt,line cap=round,line join=round](-0.13,0.0)--(-0.04,-0.09)--(0.13,0.1);}
\newcommand{\popboxhead}[3]{\poppill{#1}{#2}{white}\ifx\relax#3\relax\else\hspace{6pt}{\popxbold\fontsize{10.6}{12}\selectfont\color{popink}#3}\fi\par\vspace{7pt}}

\newtcolorbox{aretenir}[1][]{popbase,colframe=popgreen,before upper={\popboxhead{À retenir}{popgreen}{#1}}}
\newtcolorbox{exemplebox}[1][]{popbase,colframe=poporange,before upper={\popboxhead{Exemple}{poporange}{#1}}}
\newtcolorbox{definition}[1][]{popbase,colframe=popblue,before upper={\popboxhead{Définition}{popblue}{#1}}}
\newtcolorbox{propriete}[1][]{popbase,colframe=poppurple,before upper={\popboxhead{Propriété}{poppurple}{#1}}}
\newtcolorbox{methode}[1][]{popbase,colframe=popgold,before upper={\popboxhead{Méthode}{popgold}{#1}}}
\newtcolorbox{attention}[1][]{popbase,colframe=poppink,before upper={\popboxhead{Attention}{poppink}{#1}}}
\newtcolorbox{experience}[1][]{popbase,colframe=popblue,colback=popblueL,before upper={\popboxhead{Expérience}{popblue}{#1}}}
% « Le savais-tu ? » : carte violette pleine (contexte, culture, Cameroun)
\newtcolorbox{savaistu}[1][]{popbase,colback=poppurple,colframe=popink,
  fontupper=\popsemi\fontsize{10.4}{14.2}\selectfont\color{white},
  before upper={\poppill{Le savais-tu\,?}{white}{poppurple}\ifx\relax#1\relax\else\hspace{6pt}{\popxbold\color{white}#1}\fi\par\vspace{7pt}}}
% Exercice résolu : énoncé en haut, \tcblower puis la solution sur fond crème
\newcounter{popexo}\newcounter{popexr}
\newtcolorbox{exoresolu}[1][]{popbase,colframe=poporange,colbacklower=popcream,
  segmentation style={popink,line width=1.2pt,dashed},
  before upper={\stepcounter{popexr}\popboxhead{Exercice résolu \thepopexr}{poporange}{#1}},
  before lower={\poppill{Solution}{popink}{popcream}\par\vspace{6pt}}}
% Exercice (non résolu en place) — la correction est dans la partie Corrigés
\newtcolorbox{exercice}[1][]{popbase,colframe=popink,
  before upper={\stepcounter{popexo}\popboxhead{Exercice \thepopexo}{popink}{#1}}}
% Corrigé
\newtcolorbox{corrige}[1][]{popbase,colframe=popgreen,colback=popgreenL,
  before upper={\popboxhead{Corrigé}{popgreen}{#1}}}
% Objectifs / prérequis / bilan
\newtcolorbox{objectifs}{popbase,colframe=popink,colback=poporangeL,
  before upper={\popboxhead{Objectifs de la leçon}{popink}{}}}
\newtcolorbox{prerequis}{popbase,colframe=popink,colback=popblueL,
  before upper={\popboxhead{Prérequis}{popblue}{}}}
\newtcolorbox{bilan}{popbase,colframe=popink,colback=white,boxrule=2.8pt,
  before upper={\popboxhead{Fiche bilan}{poporange}{L'essentiel à connaître pour l'examen}}}
% Figure encadrée avec légende
\newtcolorbox{popfigure}[1][]{enhanced,breakable=false,colback=white,colframe=popline,boxrule=1.4pt,arc=8pt,
  left=10pt,right=10pt,top=10pt,bottom=8pt,before skip=10pt,after skip=12pt,halign=center,
  lower separated=false,halign lower=center,
  fontlower=\popsemi\fontsize{9}{11.5}\selectfont\color{popmuted},
  #1}
\newcommand{\legende}[1]{\par\vspace{6pt}{\popsemi\fontsize{9}{11.5}\selectfont\color{popmuted}#1}}

% ============================================================
% Signature OnBuch+
% ============================================================
\newcommand{\poplogoinline}[1]{{\poplogo\fontsize{#1}{#1}\selectfont\color{popink}OnBuch\color{poporange}+}}
\newcommand{\signatureonbuch}{%
  \begin{tcolorbox}[enhanced,colback=popink,colframe=popink,boxrule=2.2pt,arc=10pt,
      drop shadow=poporange,left=14pt,right=14pt,top=10pt,bottom=10pt,before skip=0pt,after skip=0pt]
    \tikz[baseline=-0.7ex]\node[circle,fill=popgreen,draw=popcream,line width=1.3pt,minimum size=22pt,inner sep=0]{\popcheck{white}};%
    \hspace{9pt}\begin{minipage}[c]{0.62\linewidth}
      {\popxbold\fontsize{12}{13}\selectfont\color{popcream}Signé\ \textbullet\ {\poplogo OnBuch\color{poporange}+}}\par\vspace{2pt}
      {\raggedright\popsemi\fontsize{8}{10.5}\selectfont\color{popline}\MakeUppercase{Équipe pédagogique OnBuch+}\\\MakeUppercase{Conforme au programme officiel MINESEC}\par}
    \end{minipage}\hfill
    {\poplogo\fontsize{22}{22}\selectfont\color{popcream}OnBuch\color{poporange}+}
  \end{tcolorbox}%
}

% ============================================================
% Page de garde
% #1 titre · #2 matière · #3 niveau · #4 module · #5 n° leçon · #6 durée ·
% #7 description / objectifs (texte ou itemize)
% ============================================================
\newcommand{\pagegarde}[7]{%
  \thispagestyle{empty}
  \newgeometry{a4paper,margin=1.7cm,top=1.6cm,bottom=1.4cm}%
  \begin{tikzpicture}[remember picture,overlay]
    \fill[poporange!18] ([xshift=-3.2cm,yshift=-3.6cm]current page.north east) circle (4.6cm);
    \fill[poppurple!14] ([xshift=2.4cm,yshift=7.6cm]current page.south west) circle (3.8cm);
  \end{tikzpicture}%
  {\poplogo\fontsize{30}{30}\selectfont\color{popink}OnBuch\color{poporange}+}\hfill
  \raisebox{0.6ex}{\poppill{Cours signé \textbullet\ Premium}{popink}{popcream}}\par
  \vspace{1pt}\popsquiggle{poppurple}\par
  \vspace{8pt}
  \begin{tcolorbox}[enhanced,colback=poporange,colframe=popink,boxrule=2.8pt,arc=13pt,
      drop shadow=popink,left=18pt,right=18pt,top=12pt,bottom=13pt,after skip=10pt]
    \poppill{#2 \textbullet\ #3}{popink}{popcream}\hspace{5pt}\poppill{#4}{white}{popink}\par\vspace{8pt}
    {\popdisplay\fontsize{14}{14}\selectfont\color{popink}LEÇON #5}\par\vspace{4pt}
    {\raggedright\hyphenpenalty=10000\exhyphenpenalty=10000\popdisplay\fontsize{25}{27}\selectfont\color{popcream}\MakeUppercase{#1}\par}
    \vspace{8pt}
    {\popxbold\fontsize{9.5}{9.5}\selectfont\color{popink}\MakeUppercase{Durée conseillée : #6}}
  \end{tcolorbox}
  \begin{tcolorbox}[enhanced,colback=white,colframe=popink,boxrule=2.2pt,arc=10pt,
      drop shadow=popink,left=16pt,right=16pt,top=10pt,bottom=10pt,after skip=8pt]
    \poppill{Dans cette leçon}{poppurple}{white}\par\vspace{5pt}
    \popsemi\fontsize{9.3}{12.6}\selectfont\color{popink}#7
  \end{tcolorbox}
  \noindent\begin{minipage}[t]{0.487\linewidth}
    \begin{tcolorbox}[enhanced,colback=popgreen,colframe=popink,boxrule=2.2pt,arc=10pt,
        drop shadow=popink,left=11pt,right=11pt,top=10pt,bottom=10pt,height=3.1cm,valign=center]
      \tcbox[colback=white,colframe=popink,boxrule=1.3pt,arc=5pt,boxsep=0pt,nobeforeafter,
        left=3pt,right=3pt,top=3pt,bottom=3pt,valign=center]{\qrcode[height=1.9cm]{https://play.google.com/store/apps/details?id=com.luvvix.onbuch&hl=fr}}%
      \hspace{8pt}\begin{minipage}[c]{0.5\linewidth}
        {\popdisplay\fontsize{11}{12.5}\selectfont\color{white}\MakeUppercase{Télécharge OnBuch}}\par\vspace{4pt}
        {\raggedright\popsemi\fontsize{8.4}{11}\selectfont\color{white}Gratuit \textbullet\ Google Play\\Tuteur IA, annales, concours, résultats\par}
      \end{minipage}
    \end{tcolorbox}
  \end{minipage}\hfill
  \begin{minipage}[t]{0.487\linewidth}
    \begin{tcolorbox}[enhanced,colback=poppurple,colframe=popink,boxrule=2.2pt,arc=10pt,
        drop shadow=popink,left=11pt,right=11pt,top=10pt,bottom=10pt,height=3.1cm,valign=center]
      \tikz[baseline=-0.7ex]\node[circle,fill=white,draw=popink,line width=1.3pt,minimum size=1.6cm,inner sep=0]{\popdisplay\fontsize{24}{24}\selectfont\color{poppurple}+};%
      \hspace{8pt}\begin{minipage}[c]{0.6\linewidth}
        {\popdisplay\fontsize{11}{12.5}\selectfont\color{white}\MakeUppercase{Passe à OnBuch+}}\par\vspace{4pt}
        {\raggedright\popsemi\fontsize{8.4}{11}\selectfont\color{white}Tous les cours signés, Tuteur IA illimité, fiches PDF — onglet OnBuch+ de l'app\par}
      \end{minipage}
    \end{tcolorbox}
  \end{minipage}
  \vspace{8pt}
  \popcontenu
  \vfill
  \signatureonbuch
  \restoregeometry
  \newpage
}

% Rangée « Ce cours contient » (page de garde)
\newcommand{\poptile}[3]{% #1 couleur #2 gros texte #3 légende
  \begin{tcolorbox}[enhanced,colback=#1,colframe=popink,boxrule=2pt,arc=9pt,shadow={2.5pt}{-2.5pt}{0pt}{popink},
      left=6pt,right=6pt,top=9pt,bottom=9pt,halign=center,valign=center,height=1.75cm,width=\linewidth,nobeforeafter]
    {\popdisplay\fontsize{12.5}{13}\selectfont\color{white}\MakeUppercase{#2}}\par\vspace{3pt}
    {\centering\popsemi\fontsize{7.8}{9.5}\selectfont\color{white}#3\par}
  \end{tcolorbox}}
\newcommand{\popcontenu}{%
  \par\noindent{\popxboldls\fontsize{7.5}{7.5}\selectfont\color{popmuted}CE COURS SIGNÉ CONTIENT}\par\vspace{7pt}
  \noindent\begin{minipage}[t]{0.235\linewidth}\poptile{popblue}{Cours}{complet, illustré, pas à pas}\end{minipage}\hfill
  \begin{minipage}[t]{0.235\linewidth}\poptile{poporange}{Méthodes}{et exercices résolus}\end{minipage}\hfill
  \begin{minipage}[t]{0.235\linewidth}\poptile{poppink}{Exercices}{type examen + corrigés}\end{minipage}\hfill
  \begin{minipage}[t]{0.235\linewidth}\poptile{popgreen}{Fiche bilan}{l'essentiel à retenir}\end{minipage}\par}

% Sommaire
\newtcolorbox{sommaire}{enhanced,colback=white,colframe=popink,boxrule=2.2pt,arc=10pt,
  drop shadow=popink,left=18pt,right=18pt,top=13pt,bottom=13pt,before skip=6pt,after skip=20pt,
  before upper={\poppill{Sommaire}{poppurple}{white}\par\vspace{9pt}\popsemi\fontsize{10.6}{14.5}\selectfont\color{popink}}}

% Signature de fin de cours — poussée tout en bas de la dernière page
\newcommand{\findecours}{%
  \par\vfill\nopagebreak
  \begin{center}\popsquiggle{poporange}\end{center}\vspace{4pt}
  \signatureonbuch
}
\AtBeginDocument{\def\llbracket{\mathopen{[\mkern-3.5mu[}}\def\rrbracket{\mathclose{]\mkern-3.5mu]}}} % intervalles d'entiers (glyphe ⟦ absent de la police)
% Glyphes absents de Fira Math : redéfinis à partir de glyphes disponibles
\AtBeginDocument{%
  \def\setminus{\mathbin{\backslash}}\let\smallsetminus\setminus
  \def\vdots{\mathord{\vbox{\baselineskip3.2pt\lineskiplimit0pt\kern4pt\hbox{.}\hbox{.}\hbox{.}}}}%
  \def\ddots{\mathinner{\mkern1mu\raise6.5pt\hbox{.}\mkern2mu\raise3.6pt\hbox{.}\mkern2mu\raise0.7pt\hbox{.}\mkern1mu}}%
  \def\triangle{\mathord{\scalebox{0.9}{$\symup{\Delta}$}}}%
  \def\nabla{\mathord{\raisebox{\depth}{\scalebox{1}[-1]{$\symup{\Delta}$}}}}%
  \def\checkmark{\popcheck{popdark}}%
  \def\star{\mathbin{\ast}}\def\bigstar{\mathord{\ast}}%
  \def\diamond{\mathbin{\scalebox{0.6}{$\Diamond$}}}%
  \def\bigcup{\mathop{\mathchoice{\raisebox{-0.3em}{\scalebox{1.7}{$\cup$}}}{\scalebox{1.25}{$\cup$}}{\cup}{\cup}}\displaylimits}%
  \def\bigcap{\mathop{\mathchoice{\raisebox{-0.3em}{\scalebox{1.7}{$\cap$}}}{\scalebox{1.25}{$\cap$}}{\cap}{\cap}}\displaylimits}%
  \def\lesseqqgtr{\mathrel{\lessgtr}}\def\bigoplus{\mathop{\oplus}\displaylimits}%
}
\providecommand{\ssi}{si et seulement si} % abréviation fréquente
\AtBeginDocument{\providecommand{\wideparen}{\overparen}} % arc de cercle

\begin{document}

\pagegarde{Dérivation}{Mathématiques}{1re Tech}{Mathématiques – classe de Première technique}{N14}{6 séances (fiche de progression harmonisée)}{Cette leçon introduit la notion de nombre dérivé d'une fonction en un point, la fonction dérivée sur un intervalle et ses opérations. Elle établit le lien fondamental entre le signe de la dérivée et le sens de variation, et applique ces outils à la tangente, à l'approximation affine et à des problèmes concrets d'optimisation.
\vspace{4pt}
\begin{multicols}{2}\small
\begin{itemize}
\item À la fin de cette leçon, je sais calculer un taux d'accroissement et déterminer le nombre dérivé d'une fonction en un réel par la définition.
\item À la fin de cette leçon, je sais étudier la dérivabilité à gauche et à droite en un réel et conclure sur la dérivabilité.
\item À la fin de cette leçon, je sais déterminer la fonction dérivée des fonctions usuelles et des sommes, produits, quotients et composées par une fonction affine.
\item À la fin de cette leçon, je sais déterminer l'équation de la tangente à une courbe en un de ses points.
\item À la fin de cette leçon, je sais utiliser le signe de la dérivée pour dresser le tableau de variations d'une fonction.
\item À la fin de cette leçon, je sais donner l'approximation affine d'une fonction autour d'un réel et l'utiliser pour estimer une valeur.
\end{itemize}\end{multicols}}

\begin{sommaire}
\begin{multicols}{2}
\begin{enumerate}
\item Nombre dérivé en un réel
\item Fonction dérivée sur un intervalle
\item Opérations sur les fonctions dérivables
\item Tangente et approximation affine
\item Dérivée et sens de variation
\item Applications à des situations concrètes
\item Activité d'intégration
\item Méthodes et exercices résolus
\item Exercices
\item Corrigés des exercices
\item Fiche bilan
\end{enumerate}
\end{multicols}
\end{sommaire}

\begin{objectifs}
\begin{itemize}
\item À la fin de cette leçon, je sais calculer un taux d'accroissement et déterminer le nombre dérivé d'une fonction en un réel par la définition.
\item À la fin de cette leçon, je sais étudier la dérivabilité à gauche et à droite en un réel et conclure sur la dérivabilité.
\item À la fin de cette leçon, je sais déterminer la fonction dérivée des fonctions usuelles et des sommes, produits, quotients et composées par une fonction affine.
\item À la fin de cette leçon, je sais déterminer l'équation de la tangente à une courbe en un de ses points.
\item À la fin de cette leçon, je sais utiliser le signe de la dérivée pour dresser le tableau de variations d'une fonction.
\item À la fin de cette leçon, je sais donner l'approximation affine d'une fonction autour d'un réel et l'utiliser pour estimer une valeur.
\item À la fin de cette leçon, je sais modéliser une situation concrète de la vie courante (coût, recette, vitesse) par une fonction et l'optimiser grâce à la dérivation.
\item À la fin de cette leçon, je sais rédiger et justifier une démarche complète : hypothèses, calculs, conclusion.
\end{itemize}
\end{objectifs}

\begin{prerequis}
\begin{itemize}
\item Notion de fonction, ensemble de définition, lecture graphique (Seconde)
\item Calcul algébrique : développement, factorisation, fractions (Seconde)
\item Notion de limite d'une fonction en un point (début d'année de Première)
\item Équations de droites : pente (coefficient directeur) et point de passage (Seconde)
\item Sens de variation d'une fonction et tableau de variations (Seconde)
\end{itemize}
\end{prerequis}

\newpage

\coursec{Nombre dérivé en un réel}

\textbf{Situation d'accroche.} Au marché Mokolo de Yaoundé, Mama Odile vend chaque matin ses beignets. Elle a remarqué quelque chose de curieux : lorsqu'elle augmente le prix du beignet de 25 FCFA, sa recette quotidienne \emph{augmente} d'abord... puis, passé un certain prix, elle \emph{diminue}, car les clients deviennent trop rares. Il existe donc un prix « idéal » qui rend la recette maximale. Mais comment le trouver autrement que par tâtonnements ?

Tout repose sur une question simple : \emph{quand une grandeur varie, à quelle vitesse varie-t-elle ?} Si l'on sait mesurer la « vitesse de variation » de la recette en fonction du prix, on peut détecter le moment précis où cette vitesse s'annule : c'est là que la recette est maximale. Cet outil de mesure de la variation instantanée s'appelle la \cle{dérivation}. C'est l'un des concepts les plus puissants des mathématiques : il sert à optimiser un bénéfice, à calculer la vitesse d'un véhicule, l'intensité d'un courant ou la pente d'une toiture.

\textbf{Problématique.} Comment définir et calculer rigoureusement la « vitesse de variation instantanée » d'une fonction en un point ?

\courssub{Le taux d'accroissement : mesurer une variation moyenne}

\textbf{Intuition.} Sur la route Yaoundé--Douala, un chauffeur parcourt 240 km en 3 heures. Sa vitesse \emph{moyenne} est $\dfrac{240}{3} = 80$ km/h : c'est le rapport entre la variation de la position et la variation du temps. De la même façon, pour une fonction $f$, on mesure la variation moyenne entre deux points en divisant la variation des images par la variation des antécédents.

\begin{definition}[Taux d'accroissement]
Soit $f$ une fonction définie sur un intervalle contenant le réel $a$, et $h$ un réel non nul tel que $a+h$ appartienne à cet intervalle. Le \cle{taux d'accroissement} (ou taux de variation) de $f$ entre $a$ et $a+h$ est le réel :
\[ \tau(h) = \frac{f(a+h) - f(a)}{h}. \]
\end{definition}

\begin{exemplebox}[Un taux d'accroissement concret]
Soit $f(x) = x^2$ et $a = 1$. Calculons $\tau(h)$ pour quelques valeurs de $h$ :
\begin{align*}
\tau(1) &= \frac{f(2)-f(1)}{1} = \frac{4-1}{1} = 3 ;\\
\tau(0{,}5) &= \frac{f(1{,}5)-f(1)}{0{,}5} = \frac{2{,}25-1}{0{,}5} = 2{,}5 ;\\
\tau(0{,}1) &= \frac{f(1{,}1)-f(1)}{0{,}1} = \frac{1{,}21-1}{0{,}1} = 2{,}1.
\end{align*}
On observe que lorsque $h$ se rapproche de $0$, le taux $\tau(h)$ se rapproche de $2$. Ce nombre $2$ sera le \emph{nombre dérivé} de $f$ en $1$.
\end{exemplebox}

\textbf{Interprétation géométrique.} Sur la courbe représentative de $f$, notons $A$ le point de coordonnées $(a\,;\,f(a))$ et $M$ le point de coordonnées $(a+h\,;\,f(a+h))$. Le taux d'accroissement $\tau(h)$ est exactement le \cle{coefficient directeur} (la pente) de la droite $(AM)$, appelée \cle{sécante} à la courbe. Quand $h$ tend vers $0$, le point $M$ glisse le long de la courbe vers le point $A$, et la sécante $(AM)$ « pivote » autour de $A$ vers une position limite : la \cle{tangente}.

\begin{popfigure}
\begin{center}
\begin{tikzpicture}
\begin{axis}[
  width=0.85\linewidth, height=7cm,
  axis lines=middle,
  xlabel={$x$}, ylabel={$y$},
  xmin=-0.5, xmax=3.2, ymin=-0.5, ymax=5.5,
  xtick={1}, ytick={1},
  legend style={at={(0.02,0.98)}, anchor=north west, font=\small}
]
\addplot[popcurve,domain=-0.3:2.3,samples=200]{x^2};
\addlegendentry{$\mathcal{C}_f : y=x^2$}
\addplot[poporange, thick, domain=-0.2:2.6]{2*x-1};
\addlegendentry{tangente en $A$}
\addplot[popgreen, dashed, domain=0.4:2.6]{3*x-2};
\addplot[poppurple, dashed, domain=0.4:2.6]{2.5*x-1.5};
\addplot[popblue, dashed, domain=0.4:2.6]{2.1*x-1.1};
\addplot[only marks, mark=*, popdark] coordinates {(1,1) (2,4) (1.5,2.25) (1.1,1.21)};
\node[popdark, below right] at (axis cs:1,1) {$A(1\,;\,1)$};
\node[popgreen, above] at (axis cs:2,4) {$M_1$};
\node[poppurple, right] at (axis cs:1.5,2.25) {$M_2$};
\node[popblue, right] at (axis cs:1.1,1.21) {$M_3$};
\end{axis}
\end{tikzpicture}
\end{center}
\legende{Quand $M$ se rapproche de $A$, la sécante $(AM)$ pivote vers la tangente en $A$, de pente $f'(1)=2$.}
\end{popfigure}

\courssub{Nombre dérivé d'une fonction en un réel}

\textbf{Intuition.} Le taux d'accroissement donne une vitesse \emph{moyenne} sur un intervalle. Pour obtenir la vitesse \emph{instantanée}, on resserre l'intervalle à l'extrême : on fait tendre $h$ vers $0$ et on regarde si le taux se stabilise vers un nombre fini.

\begin{definition}[Nombre dérivé en un réel]
Soit $f$ une fonction définie sur un intervalle ouvert contenant le réel $a$. On dit que $f$ est \cle{dérivable en $a$} si le taux d'accroissement $\tau(h) = \dfrac{f(a+h)-f(a)}{h}$ admet une \textbf{limite finie} lorsque $h$ tend vers $0$. Cette limite est alors appelée \cle{nombre dérivé} de $f$ en $a$ et se note $f'(a)$ :
\[ f'(a) = \lim_{h \to 0} \frac{f(a+h)-f(a)}{h}. \]
\end{definition}

\begin{aretenir}[Interprétations du nombre dérivé]
\begin{itemize}
\item \textbf{Géométrique} : $f'(a)$ est le \cle{coefficient directeur} (la pente) de la tangente à la courbe $\mathcal{C}_f$ au point $A(a\,;\,f(a))$.
\item \textbf{Cinématique} : si $x(t)$ est la position d'un mobile à l'instant $t$, alors $x'(t_0)$ est sa \cle{vitesse instantanée} à l'instant $t_0$.
\item \textbf{Économique} : si $R(p)$ est la recette de Mama Odile en fonction du prix $p$, alors $R'(p)$ mesure la variation instantanée de la recette quand le prix varie.
\end{itemize}
\end{aretenir}

\begin{exoresolu}[Calcul de $f'(2)$ par la définition]
\textbf{Énoncé.} Soit $f$ la fonction définie sur $\mathbb{R}$ par $f(x) = x^2 - 3x$. En utilisant la définition, montrer que $f$ est dérivable en $a = 2$ et déterminer $f'(2)$.
\tcblower
\textbf{Solution détaillée.} On calcule d'abord le taux d'accroissement entre $2$ et $2+h$, avec $h \neq 0$ :
\begin{align*}
\tau(h) &= \frac{f(2+h)-f(2)}{h} \\
&= \frac{(2+h)^2 - 3(2+h) - (4 - 6)}{h} \\
&= \frac{4 + 4h + h^2 - 6 - 3h + 2}{h} \\
&= \frac{h^2 + h}{h} = h + 1.
\end{align*}
Lorsque $h$ tend vers $0$, $\tau(h) = h+1$ tend vers $1$, qui est un nombre \emph{fini}. Donc $f$ est dérivable en $2$ et :
\[ f'(2) = 1. \]
\textbf{Interprétation} : la tangente à la courbe de $f$ au point d'abscisse $2$ a pour pente $1$.
\end{exoresolu}

\begin{methode}[Étudier la dérivabilité en un point par la définition]
Pour étudier la dérivabilité de $f$ en $a$ :
\begin{enumerate}
\item Vérifier que $f$ est définie en $a$ et au voisinage de $a$.
\item Former le taux d'accroissement $\tau(h) = \dfrac{f(a+h)-f(a)}{h}$.
\item \textbf{Simplifier} $\tau(h)$ en factorisant par $h$ au numérateur (c'est presque toujours possible quand $f$ est dérivable), puis simplifier par $h$.
\item Chercher la limite quand $h$ tend vers $0$ : si elle est \emph{finie}, $f$ est dérivable en $a$ et $f'(a)$ est cette limite ; sinon, $f$ n'est pas dérivable en $a$.
\end{enumerate}
\end{methode}

\courssub{Exemples guidés : fonctions de référence}

\begin{exemplebox}[Dérivabilité de $f(x) = x^2$ en tout réel $a$]
Soit $f(x) = x^2$ et $a$ un réel quelconque. Pour $h \neq 0$ :
\[ \tau(h) = \frac{(a+h)^2 - a^2}{h} = \frac{a^2 + 2ah + h^2 - a^2}{h} = \frac{h(2a+h)}{h} = 2a + h. \]
Quand $h$ tend vers $0$, $\tau(h)$ tend vers $2a$. Donc $f$ est dérivable en tout réel $a$ et $\boxed{f'(a) = 2a}$. Par exemple $f'(1) = 2$, ce qui confirme l'observation de la figure précédente.
\end{exemplebox}

\begin{exemplebox}[Dérivabilité de $f(x) = \dfrac{1}{x}$ en $a \neq 0$]
Soit $f(x) = \dfrac{1}{x}$, définie pour $x \neq 0$, et $a \neq 0$. Pour $h \neq 0$ :
\begin{align*}
\tau(h) &= \frac{\dfrac{1}{a+h} - \dfrac{1}{a}}{h} = \frac{\dfrac{a - (a+h)}{a(a+h)}}{h} = \frac{-h}{h \cdot a(a+h)} = \frac{-1}{a(a+h)}.
\end{align*}
Quand $h$ tend vers $0$, $a+h$ tend vers $a$, donc $\tau(h)$ tend vers $-\dfrac{1}{a^2}$, qui est fini. Donc $f$ est dérivable en tout $a \neq 0$ et $\boxed{f'(a) = -\dfrac{1}{a^2}}$. Par exemple $f'(2) = -\dfrac{1}{4}$ : la tangente en $2$ « descend ».
\end{exemplebox}

\begin{exemplebox}[Dérivabilité de $f(x) = \sqrt{x}$ en $a > 0$]
Soit $f(x) = \sqrt{x}$ et $a > 0$. Pour $h \neq 0$, on utilise l'\emph{expression conjuguée} pour faire disparaître les racines au numérateur :
\begin{align*}
\tau(h) &= \frac{\sqrt{a+h}-\sqrt{a}}{h} = \frac{(\sqrt{a+h}-\sqrt{a})(\sqrt{a+h}+\sqrt{a})}{h(\sqrt{a+h}+\sqrt{a})} \\
&= \frac{(a+h)-a}{h(\sqrt{a+h}+\sqrt{a})} = \frac{1}{\sqrt{a+h}+\sqrt{a}}.
\end{align*}
Quand $h$ tend vers $0$, le dénominateur tend vers $2\sqrt{a}$, donc $\tau(h)$ tend vers $\dfrac{1}{2\sqrt{a}}$. Donc $f$ est dérivable en tout $a > 0$ et $\boxed{f'(a) = \dfrac{1}{2\sqrt{a}}}$.
\end{exemplebox}

\begin{attention}[La fonction racine carrée n'est pas dérivable en 0]
Pour $f(x) = \sqrt{x}$ en $a = 0$ : $\tau(h) = \dfrac{\sqrt{h}}{h} = \dfrac{1}{\sqrt{h}}$ (pour $h > 0$). Quand $h$ tend vers $0$, $\tau(h)$ devient infiniment grand : la limite n'est \textbf{pas finie}. Donc $\sqrt{x}$ n'est \emph{pas} dérivable en $0$ : sa courbe admet en $0$ une tangente \emph{verticale} (pente infinie). La condition « limite finie » de la définition est essentielle.
\end{attention}

\courssub{Dérivabilité à gauche et à droite en un réel}

\textbf{Intuition.} Certaines courbes présentent un « point anguleux » : à gauche du point, la courbe arrive avec une certaine pente ; à droite, elle repart avec une autre pente. Il n'y a alors pas de tangente unique, mais deux \emph{demi-tangentes}. Pour décrire cette situation, on étudie le taux d'accroissement séparément de chaque côté.

\begin{definition}[Nombres dérivés à gauche et à droite]
Soit $f$ une fonction définie au voisinage de $a$.
\begin{itemize}
\item On dit que $f$ est \cle{dérivable à gauche en $a$} si $\tau(h)$ admet une limite finie quand $h$ tend vers $0$ par valeurs \emph{négatives} ($h < 0$). Cette limite est notée $f'_g(a)$ :
\[ f'_g(a) = \lim_{\substack{h \to 0 \\ h < 0}} \frac{f(a+h)-f(a)}{h}. \]
\item On dit que $f$ est \cle{dérivable à droite en $a$} si $\tau(h)$ admet une limite finie quand $h$ tend vers $0$ par valeurs \emph{positives} ($h > 0$). Cette limite est notée $f'_d(a)$ :
\[ f'_d(a) = \lim_{\substack{h \to 0 \\ h > 0}} \frac{f(a+h)-f(a)}{h}. \]
\end{itemize}
\end{definition}

\begin{propriete}[Condition de dérivabilité en un point]
La fonction $f$ est dérivable en $a$ \textbf{si et seulement si} elle est dérivable à gauche et à droite en $a$ et que ces deux nombres dérivés sont \emph{égaux} :
\[ f \text{ dérivable en } a \iff f'_g(a) = f'_d(a) \quad (\text{existants et finis}). \]
Dans ce cas, $f'(a) = f'_g(a) = f'_d(a)$.
\end{propriete}

\begin{exemplebox}[Contre-exemple fondamental : $f(x) = |x|$ en 0]
Soit $f(x) = |x|$. Étudions sa dérivabilité en $0$. On a $f(0) = 0$, donc :
\[ \tau(h) = \frac{|h| - 0}{h} = \frac{|h|}{h}. \]
\begin{itemize}
\item Si $h > 0$ : $|h| = h$, donc $\tau(h) = \dfrac{h}{h} = 1$. Ainsi $f'_d(0) = 1$.
\item Si $h < 0$ : $|h| = -h$, donc $\tau(h) = \dfrac{-h}{h} = -1$. Ainsi $f'_g(0) = -1$.
\end{itemize}
Comme $f'_g(0) = -1 \neq 1 = f'_d(0)$, la fonction valeur absolue \textbf{n'est pas dérivable en 0}. Sa courbe présente en $0$ un \cle{point anguleux} avec deux demi-tangentes de pentes $-1$ et $1$.
\end{exemplebox}

\begin{popfigure}
\begin{center}
\begin{tikzpicture}
\begin{axis}[
  width=0.8\linewidth, height=6.5cm,
  axis lines=middle,
  xlabel={$x$}, ylabel={$y$},
  xmin=-2.6, xmax=2.6, ymin=-0.6, ymax=2.6,
  xtick={-2,-1,1,2}, ytick={1,2}
]
\addplot[popcurve,domain=-2.3:2.3,samples=100]{abs(x)};
\addlegendentry{$y = |x|$}
\addplot[poporange, dashed, thick, domain=0:2.3]{x};
\addlegendentry{demi-tangente droite, pente $1$}
\addplot[popgreen, dashed, thick, domain=-2.3:0]{-x};
\addlegendentry{demi-tangente gauche, pente $-1$}
\addplot[only marks, mark=*, popdark] coordinates {(0,0)};
\node[popdark, below] at (axis cs:0,0) {$O$};
\end{axis}
\end{tikzpicture}
\end{center}
\legende{La courbe de $x \mapsto |x|$ admet en $0$ deux demi-tangentes de pentes différentes : elle n'est pas dérivable en $0$.}
\end{popfigure}

\begin{methode}[Dérivabilité d'une fonction définie par morceaux en un point de raccordement]
Soit $f$ définie par deux expressions différentes de part et d'autre de $a$. Pour étudier sa dérivabilité en $a$ :
\begin{enumerate}
\item \textbf{Continuité d'abord} : vérifier que $f$ est continue en $a$ (les deux expressions doivent donner la même valeur en $a$). Si $f$ n'est pas continue en $a$, elle ne peut pas être dérivable en $a$.
\item Calculer le taux d'accroissement \textbf{à gauche} avec l'expression valable pour $x < a$, et en déduire $f'_g(a)$.
\item Calculer le taux d'accroissement \textbf{à droite} avec l'expression valable pour $x > a$, et en déduire $f'_d(a)$.
\item \textbf{Conclure} : si $f'_g(a) = f'_d(a)$ (finis), $f$ est dérivable en $a$ ; sinon, elle ne l'est pas (point anguleux).
\end{enumerate}
\end{methode}

\begin{exoresolu}[Un raccordement d'atelier]
\textbf{Énoncé.} Une tôle est pliée selon le profil $f(x) = \begin{cases} x^2 & \text{si } x \leq 1 \\ 2x - 1 & \text{si } x > 1 \end{cases}$. La fonction $f$ est-elle dérivable en $1$ ?
\tcblower
\textbf{Solution détaillée.}
\textbf{1. Continuité en 1.} $f(1) = 1^2 = 1$ et, pour $x > 1$, $2x-1$ tend vers $1$ quand $x$ tend vers $1$. Donc $f$ est continue en $1$.
\textbf{2. Dérivabilité à gauche.} Pour $h < 0$, on utilise $f(x) = x^2$ :
\[ \tau(h) = \frac{(1+h)^2 - 1}{h} = \frac{2h + h^2}{h} = 2 + h \xrightarrow[h \to 0^-]{} 2. \]
Donc $f'_g(1) = 2$.
\textbf{3. Dérivabilité à droite.} Pour $h > 0$, on utilise $f(x) = 2x - 1$ :
\[ \tau(h) = \frac{2(1+h) - 1 - 1}{h} = \frac{2h}{h} = 2 \xrightarrow[h \to 0^+]{} 2. \]
Donc $f'_d(1) = 2$.
\textbf{4. Conclusion.} $f'_g(1) = f'_d(1) = 2$, donc $f$ \textbf{est dérivable en 1} et $f'(1) = 2$ : le pliage de la tôle est « doux », sans point anguleux.
\end{exoresolu}

\courssub{Lien entre dérivabilité et continuité}

\begin{propriete}[Dérivable implique continue]
Si une fonction $f$ est dérivable en $a$, alors $f$ est \cle{continue} en $a$, c'est-à-dire :
\[ \lim_{h \to 0} f(a+h) = f(a). \]
\textbf{Attention : la réciproque est fausse.} Une fonction peut être continue en $a$ sans être dérivable en $a$.
\end{propriete}

\begin{experience}[Démonstration guidée : pourquoi dérivable implique continue]
Supposons $f$ dérivable en $a$. Nous allons montrer que $f(a+h) - f(a)$ tend vers $0$ quand $h$ tend vers $0$.
\begin{enumerate}
\item \textbf{Point de départ.} Pour $h \neq 0$, complétez l'égalité obtenue en multipliant le taux d'accroissement par $h$ :
\[ f(a+h) - f(a) = h \times \frac{f(a+h)-f(a)}{h} = h \cdot \tau(h). \]
\item \textbf{Passage à la limite.} Quand $h$ tend vers $0$ :
\begin{itemize}
\item $h$ tend vers $0$ ;
\item $\tau(h)$ tend vers $f'(a)$, qui est un nombre \emph{fini} (car $f$ est dérivable en $a$).
\end{itemize}
\item \textbf{Conclusion.} Le produit $h \cdot \tau(h)$ tend donc vers $0 \times f'(a) = 0$. Ainsi :
\[ \lim_{h \to 0} \bigl(f(a+h) - f(a)\bigr) = 0, \quad \text{c'est-à-dire} \quad \lim_{h \to 0} f(a+h) = f(a). \]
La fonction $f$ est donc continue en $a$. $\blacksquare$
\end{enumerate}
\end{experience}

\begin{attention}[Erreur fréquente : confondre « continue » et « dérivable »]
\begin{itemize}
\item \textbf{Continuité} : la courbe se trace « sans lever le crayon ».
\item \textbf{Dérivabilité} : la courbe se trace sans lever le crayon \emph{et} sans point anguleux ni tangente verticale.
\end{itemize}
La fonction $f(x) = |x|$ est \textbf{continue en 0} (sa courbe ne se brise pas) mais \textbf{pas dérivable en 0} (point anguleux, deux demi-tangentes de pentes $-1$ et $1$). Ne jamais écrire « $f$ est continue en $a$, donc $f$ est dérivable en $a$ » : c'est le piège classique du Probatoire ! L'implication correcte est : dérivable $\Rightarrow$ continue.
\end{attention}

\begin{aretenir}[L'essentiel de la section]
\begin{itemize}
\item Le taux d'accroissement de $f$ entre $a$ et $a+h$ est $\tau(h) = \dfrac{f(a+h)-f(a)}{h}$ : c'est la pente de la sécante $(AM)$.
\item $f$ est dérivable en $a$ si $\tau(h)$ admet une \textbf{limite finie} quand $h \to 0$ ; cette limite est $f'(a)$, pente de la tangente en $A(a\,;\,f(a))$.
\item Dérivabilité en $a$ $\iff$ $f'_g(a) = f'_d(a)$ (finis). Sinon : point anguleux (ex. $|x|$ en $0$) ou tangente verticale (ex. $\sqrt{x}$ en $0$).
\item Dérivable en $a$ $\Rightarrow$ continue en $a$ ; la réciproque est fausse.
\item À retenir : $(x^2)'$ en $a$ vaut $2a$ ; $\left(\dfrac{1}{x}\right)'$ en $a$ vaut $-\dfrac{1}{a^2}$ ; $(\sqrt{x})'$ en $a$ vaut $\dfrac{1}{2\sqrt{a}}$.
\end{itemize}
\end{aretenir}

\begin{savaistu}[De Newton aux ateliers camerounais]
La notion de dérivée a été introduite au XVII\textsuperscript{e} siècle, indépendamment, par Isaac Newton (pour étudier les vitesses des corps en mouvement) et Gottfried Wilhelm Leibniz (pour étudier les tangentes aux courbes). Trois siècles plus tard, elle est partout : l'ingénieur qui calcule la pente maximale d'une rampe d'accès à Bafoussam, l'électricien qui étudie la variation d'un courant, le gestionnaire qui cherche le prix maximisant le bénéfice — comme Mama Odile et ses beignets — utilisent, souvent sans le savoir, le nombre dérivé. Dans les prochaines sections, nous apprendrons à dériver rapidement grâce à des formules, puis à utiliser le signe de la dérivée pour trouver les variations d'une fonction et résoudre le problème de Mama Odile.
\end{savaistu}

\begin{exercice}[Exercices d'entraînement - Section 1]
\begin{enumerate}
\item Soit $f(x) = 3x^2 + x$. Calculer $f'(1)$ en utilisant la définition du nombre dérivé.
\item Soit $g(x) = \dfrac{1}{x}$. Calculer $g'(3)$ puis donner la pente de la tangente à la courbe de $g$ au point d'abscisse $3$.
\item Soit $h$ définie par $h(x) = \begin{cases} x + 1 & \text{si } x \leq 0 \\ 2x + 1 & \text{si } x > 0 \end{cases}$. Étudier la continuité puis la dérivabilité de $h$ en $0$.
\end{enumerate}
\end{exercice}

\begin{corrige}[Exercices section 1]
\textbf{1.} $\tau(h) = \dfrac{3(1+h)^2 + (1+h) - 4}{h} = \dfrac{3h^2 + 7h}{h} = 3h + 7 \to 7$. Donc $f'(1) = 7$.

\textbf{2.} D'après le cours, $g'(a) = -\dfrac{1}{a^2}$, donc $g'(3) = -\dfrac{1}{9}$. La tangente en $3$ a pour pente $-\dfrac{1}{9}$.

\textbf{3.} Continuité : $h(0) = 1$ et $2x+1 \to 1$ quand $x \to 0^+$ : $h$ est continue en $0$. À gauche : $\tau(h) = \dfrac{(h+1)-1}{h} = 1$, donc $h'_g(0) = 1$. À droite : $\tau(h) = \dfrac{(2h+1)-1}{h} = 2$, donc $h'_d(0) = 2$. Comme $h'_g(0) \neq h'_d(0)$, $h$ n'est \textbf{pas dérivable en 0} (point anguleux).
\end{corrige}

\coursec{Fonction dérivée sur un intervalle}

\courssub{Définition et notation}

\textbf{Intuition.} Au lieu de calculer le nombre dérivé point par point, on peut construire une \emph{nouvelle fonction} qui associe à chaque $x$ de l'intervalle le nombre dérivé $f'(x)$. C'est la \cle{fonction dérivée}.

\begin{definition}[Fonction dérivée sur un intervalle]
Soit $f$ une fonction définie sur un intervalle $I$. On dit que $f$ est \cle{dérivable sur $I$} si elle est dérivable en tout point de $I$ (pour les bornes de $I$, on considère la dérivabilité à droite ou à gauche selon le cas). La fonction qui associe à tout $x \in I$ le nombre dérivé $f'(x)$ est appelée \cle{fonction dérivée de $f$ sur $I$} et se note $f'$.
\end{definition}

\begin{aretenir}[Notations usuelles pour la dérivée]
Pour une fonction $y = f(x)$, on note indifféremment :
\[ f'(x) \quad ; \quad y' \quad ; \quad \frac{dy}{dx} \quad ; \quad \frac{d}{dx}f(x). \]
La notation $\frac{dy}{dx}$ (notation de Leibniz) rappelle que la dérivée est une limite de rapport de variations $\frac{\Delta y}{\Delta x}$.
\end{aretenir}

\courssub{Tableau des dérivées usuelles (à connaître par cœur)}

Les dérivées des fonctions de référence, obtenues par la définition (comme dans la section précédente), sont les « briques de base » pour dériver toute fonction élémentaire.

\begin{propriete}[Dérivées des fonctions de référence]
Soient $a$ un réel et $n$ un entier naturel non nul ($n \in \mathbb{N}^*$). Sur leurs domaines de définition respectifs :
\begin{center}
\begin{tabularx}{\linewidth}{|l|X|l|}\hline
\textbf{Fonction $f(x)$} & \textbf{Domaine $D_f$} & \textbf{Dérivée $f'(x)$} \\ \hline
$k$ (constante) & $\mathbb{R}$ & $0$ \\ \hline
$x^n$ & $\mathbb{R}$ & $n x^{n-1}$ \\ \hline
$\dfrac{1}{x} = x^{-1}$ & $\mathbb{R}^*$ & $-\dfrac{1}{x^2} = -x^{-2}$ \\ \hline
$\sqrt{x} = x^{1/2}$ & $[0, +\infty[$ & $\dfrac{1}{2\sqrt{x}} = \frac{1}{2}x^{-1/2}$ \quad ($x > 0$) \\ \hline
\end{tabularx}
\end{center}
\textbf{Remarque :} La formule $(x^n)' = n x^{n-1}$ reste valable pour $n \in \mathbb{Z}$ (entier relatif) et même $n \in \mathbb{Q}$ (rationnel) sur les domaines adaptés, mais au programme de Première Tech, on l'utilise principalement pour $n \in \mathbb{N}^*$ et $n=-1, -2, \frac{1}{2}$.
\end{propriete}

\begin{exemplebox}[Dérivée de $f(x) = x^3$ par la définition (généralisation)]
Soit $f(x) = x^3$. Pour $a \in \mathbb{R}$ :
\[ \tau(h) = \frac{(a+h)^3 - a^3}{h} = \frac{a^3 + 3a^2h + 3ah^2 + h^3 - a^3}{h} = 3a^2 + 3ah + h^2 \xrightarrow[h\to 0]{} 3a^2. \]
On retrouve bien la formule $n x^{n-1}$ avec $n=3$ : $f'(x) = 3x^2$.
\end{exemplebox}

\coursec{Opérations sur les fonctions dérivables}

\courssub{Règles de dérivation (somme, produit, quotient)}

\textbf{Intuition.} On ne veut pas refaire la définition à chaque fois. Heureusement, la dérivée « respecte » les opérations algébriques usuelles.

\begin{propriete}[Opérations sur les fonctions dérivables]
Soient $u$ et $v$ deux fonctions dérivables sur un intervalle $I$.
\begin{enumerate}
\item \textbf{Somme :} $u+v$ est dérivable sur $I$ et $(u+v)' = u' + v'$.
\item \textbf{Produit :} $uv$ est dérivable sur $I$ et $(uv)' = u'v + uv'$.
\item \textbf{Quotient :} Si $v(x) \neq 0$ sur $I$, $\dfrac{u}{v}$ est dérivable sur $I$ et $\left(\dfrac{u}{v}\right)' = \dfrac{u'v - uv'}{v^2}$.
\item \textbf{Multiple par une constante :} $(\lambda u)' = \lambda u'$ ($\lambda \in \mathbb{R}$).
\end{enumerate}
\end{propriete}

\begin{attention}[Pièges classiques sur les opérations]
\begin{itemize}
\item \textbf{Produit :} $(uv)' \neq u'v'$ ! Il faut \emph{les deux termes} $u'v + uv'$. Mnémotechnique : « dérivée du premier fois le second \textbf{plus} le premier fois dérivée du second ».
\item \textbf{Quotient :} L'ordre au numérateur est crucial : $u'v - uv'$ (dérivée du numérateur $\times$ dénominateur \textbf{moins} numérateur $\times$ dérivée du dénominateur). Le dénominateur est $v^2$ (jamais $v$ seul).
\item \textbf{Domaine :} Pour le quotient, il faut \emph{impérativement} préciser le domaine : $\{ x \in I \mid v(x) \neq 0 \}$.
\end{itemize}
\end{attention}

\begin{exoresolu}[Application des règles de dérivation]
\textbf{Énoncé.} Dériver les fonctions suivantes sur leur domaine de définition.
\begin{enumerate}
\item $f(x) = 3x^4 - 5x^2 + 2x - 7$
\item $g(x) = (x^2 + 1)(x - 3)$
\item $h(x) = \dfrac{x^2 + 1}{x - 1}$
\end{enumerate}
\tcblower
\textbf{Solution.}
\begin{enumerate}
\item $f$ est une somme de monômes. $f'(x) = 3(4x^3) - 5(2x) + 2 - 0 = \boxed{12x^3 - 10x + 2}$.
\item $g$ est un produit. $u(x)=x^2+1 \Rightarrow u'=2x$ ; $v(x)=x-3 \Rightarrow v'=1$.
$g'(x) = u'v + uv' = 2x(x-3) + (x^2+1)\cdot 1 = 2x^2 - 6x + x^2 + 1 = \boxed{3x^2 - 6x + 1}$.
\item $h$ est un quotient. $u(x)=x^2+1 \Rightarrow u'=2x$ ; $v(x)=x-1 \Rightarrow v'=1$. Domaine : $\mathbb{R} \setminus \{1\}$.
$h'(x) = \dfrac{2x(x-1) - (x^2+1)\cdot 1}{(x-1)^2} = \dfrac{2x^2 - 2x - x^2 - 1}{(x-1)^2} = \boxed{\dfrac{x^2 - 2x - 1}{(x-1)^2}}$.
\end{enumerate}
\end{exoresolu}

\courssub{Dérivée de la composée d'une fonction affine par une fonction donnée}

\textbf{Intuition.} On rencontre souvent des fonctions de la forme $f(ax+b)$ (ex : $\sqrt{2x+1}$, $(3x-2)^4$, $\frac{1}{5x+1}$). Il existe une règle simple pour les dériver sans tout développer.

\begin{propriete}[Règle de la chaîne pour une fonction affine intérieure]
Soit $u$ une fonction dérivable sur un intervalle $J$ et $a,b \in \mathbb{R}$ avec $a \neq 0$. La fonction $x \mapsto u(ax+b)$ est dérivable sur l'intervalle $I = \{ x \in \mathbb{R} \mid ax+b \in J \}$ et sa dérivée est :
\[ \bigl( u(ax+b) \bigr)' = a \cdot u'(ax+b). \]
\textbf{En mots :} on dérive la fonction extérieure $u$ en gardant l'intérieur $(ax+b)$, et on multiplie par la dérivée de l'intérieur (qui est $a$).
\end{propriete}

\begin{exemplebox}[Applications de la règle de la chaîne (affine)]
\begin{enumerate}
\item $f(x) = \sqrt{2x+1} = (2x+1)^{1/2}$. $f'(x) = 2 \cdot \dfrac{1}{2}(2x+1)^{-1/2} = \boxed{\dfrac{1}{\sqrt{2x+1}}}$ (pour $2x+1 > 0$).
\item $g(x) = (3x-2)^4$. $g'(x) = 3 \cdot 4(3x-2)^3 = \boxed{12(3x-2)^3}$.
\item $h(x) = \dfrac{1}{5x+1} = (5x+1)^{-1}$. $h'(x) = 5 \cdot (-1)(5x+1)^{-2} = \boxed{-\dfrac{5}{(5x+1)^2}}$ (pour $x \neq -1/5$).
\end{enumerate}
\end{exemplebox}

\begin{methode}[Dériver une fonction élémentaire : démarche complète]
Pour dériver une fonction $f$ donnée par une expression algébrique :
\begin{enumerate}
\item \textbf{Identifier la structure} : somme, produit, quotient, composée $u(ax+b)$ ? Souvent c'est un mélange.
\item \textbf{Écrire le domaine de définition $D_f$} (zéros des dénominateurs, radicands positifs pour racine carrée).
\item \textbf{Appliquer les règles} étape par étape, en notant les fonctions intermédiaires $u, v$ si nécessaire.
\item \textbf{Simplifier l'expression} de $f'(x)$ (factoriser, mettre au même dénominateur).
\item \textbf{Rappeler le domaine de $f'$} (même que $D_f$ sauf points où $f'$ n'est pas définie, ex: racine au dénominateur).
\end{enumerate}
\end{methode}

\begin{exoresolu}[Dérivation complète : mélange d'opérations]
\textbf{Énoncé.} Soit $f(x) = \dfrac{\sqrt{x}}{x+1}$. Déterminer $D_f$, calculer $f'(x)$ et donner $f'(1)$.
\tcblower
\textbf{Solution.}
\textbf{1. Domaine.} $\sqrt{x}$ impose $x \ge 0$. Dénominateur $x+1 \neq 0 \Rightarrow x \neq -1$ (déjà exclu). Donc $D_f = [0, +\infty[$.
\textbf{2. Structure.} Quotient $u/v$ avec $u(x)=\sqrt{x}$, $v(x)=x+1$.
$u'(x) = \dfrac{1}{2\sqrt{x}}$ (pour $x>0$), $v'(x) = 1$.
\textbf{3. Dérivée (pour $x > 0$) :}
\[ f'(x) = \frac{u'v - uv'}{v^2} = \frac{\frac{1}{2\sqrt{x}}(x+1) - \sqrt{x}\cdot 1}{(x+1)^2}. \]
\textbf{4. Simplification :} Multiplier numérateur et dénominateur par $2\sqrt{x}$ :
\[ f'(x) = \frac{(x+1) - 2x}{2\sqrt{x}(x+1)^2} = \frac{1 - x}{2\sqrt{x}(x+1)^2}. \]
\textbf{5. Domaine de $f'$ :} $]0, +\infty[$ (car $\sqrt{x}$ au dénominateur).
\textbf{6. Valeur en 1 :} $f'(1) = \dfrac{1-1}{2\cdot 1 \cdot 4} = \boxed{0}$.
\end{exoresolu}

\coursec{Tangente et approximation affine}

\courssub{Équation de la tangente}

\textbf{Intuition.} Puisque $f'(a)$ est la pente de la tangente en $A(a\,;\,f(a))$, et que cette tangente passe par $A$, on peut écrire son équation cartésienne immédiatement.

\begin{propriete}[Équation de la tangente]
Soit $f$ dérivable en $a$. La tangente $\mathcal{T}$ à la courbe $\mathcal{C}_f$ au point $A(a\,;\,f(a))$ a pour équation :
\[ \mathcal{T} : \quad y = f'(a)(x - a) + f(a). \]
\end{propriete}

\begin{exemplebox}[Tangente à une courbe]
Soit $f(x) = x^3 - 2x + 1$. Déterminer l'équation de la tangente en $a=1$.
$f'(x) = 3x^2 - 2 \Rightarrow f'(1) = 1$. $f(1) = 1 - 2 + 1 = 0$.
Tangente : $y = 1 \cdot (x - 1) + 0 \iff \boxed{y = x - 1}$.
\end{exemplebox}

\courssub{Approximation affine (ou linéaire)}

\textbf{Intuition.} Près de $a$, la courbe $\mathcal{C}_f$ est presque confondue avec sa tangente. On peut donc approximer $f(a+h)$ par la valeur de la tangente en $a+h$. C'est l'\cle{approximation affine}.

\begin{propriete}[Approximation affine]
Si $f$ est dérivable en $a$, pour $h$ « petit » (proche de $0$) :
\[ f(a+h) \approx f(a) + h f'(a). \]
L'erreur commise est d'autant plus petite que $h$ est proche de $0$. C'est l'écriture « fonctionnelle » de $f(a+h) = f(a) + h f'(a) + o(h)$.
\end{propriete}

\begin{exemplebox}[Calcul approché de $\sqrt{4,1}$]
On veut approximer $\sqrt{4,1}$. On prend $f(x) = \sqrt{x}$, $a = 4$, $h = 0,1$.
$f(4) = 2$, $f'(x) = \frac{1}{2\sqrt{x}} \Rightarrow f'(4) = \frac{1}{4}$.
Approximation : $f(4,1) \approx 2 + 0,1 \times \frac{1}{4} = 2 + 0,025 = 2,025$.
Valeur exacte (calculatrice) : $\sqrt{4,1} \approx 2,02484...$ L'erreur est de l'ordre de $10^{-4}$.
\end{exemplebox}

\begin{attention}[Conditions de validité]
L'approximation affine n'est valable que \textbf{au voisinage immédiat} de $a$. Ne l'utilisez pas pour $h$ grand (ex: approximer $\sqrt{100}$ à partir de $a=4$ donnerait un résultat absurde). Au Probatoire, on demande souvent : « Donner une valeur approchée de $f(a+h)$ par approximation affine » $\Rightarrow$ appliquer la formule $f(a) + h f'(a)$.
\end{attention}

\coursec{Dérivée et sens de variation}

\courssub{Théorème fondamental : signe de la dérivée et variations}

\textbf{Intuition.} Si la pente de la tangente est positive, la courbe « monte » : la fonction croît. Si elle est négative, la courbe « descend » : la fonction décroît. Si elle est nulle, la courbe est plate : extremum possible.

\begin{propriete}[Dérivée et variations (Théorème fondamental)]
Soit $f$ une fonction dérivable sur un intervalle $I$.
\begin{itemize}
\item Si $f'(x) \ge 0$ sur $I$, alors $f$ est \cle{croissante} sur $I$.
\item Si $f'(x) \le 0$ sur $I$, alors $f$ est \cle{décroissante} sur $I$.
\item Si $f'(x) = 0$ sur $I$, alors $f$ est \cle{constante} sur $I$.
\end{itemize}
\textbf{Variantes strictes :} Si $f'(x) > 0$ sur $I$ (sauf éventuellement en un nombre fini de points), alors $f$ est \textbf{strictement croissante}. De même pour $< 0$ et strictement décroissante.
\end{propriete}

\begin{experience}[Démonstration guidée : pourquoi $f' \ge 0 \Rightarrow f$ croissante]
Soit $x_1 < x_2$ dans $I$. On veut montrer $f(x_1) \le f(x_2)$.
\begin{enumerate}
\item Appliquer le \textbf{Théorème des accroissements finis} (Théorème de Rolle généralisé / Lagrange) à $f$ sur $[x_1, x_2]$ : il existe $c \in ]x_1, x_2[$ tel que $f'(c) = \dfrac{f(x_2)-f(x_1)}{x_2-x_1}$.
\item Comme $x_2 - x_1 > 0$ et $f'(c) \ge 0$ (hypothèse), le quotient est $\ge 0$.
\item Donc $f(x_2) - f(x_1) \ge 0$, c'est-à-dire $f(x_1) \le f(x_2)$. $\blacksquare$
\end{enumerate}
\textit{Note : Ce théorème est admis en Première Technique, mais sa démonstration via les accroissements finis est exigible en Terminale.}
\end{experience}

\courssub{Tableau de variations : construction et lecture}

\begin{methode}[Construire le tableau de variations d'une fonction dérivable]
\begin{enumerate}
\item Calculer $f'(x)$ et la simplifier au maximum (forme factorisée de préférence).
\item Déterminer le \textbf{signe de $f'(x)$} sur l'intervalle d'étude (tableau de signes).
\item Tracer le tableau de variations :
\begin{itemize}
\item Première ligne : valeurs de $x$ (bornes de l'intervalle + valeurs annulant $f'$).
\item Deuxième ligne : signe de $f'$ ($+$, $-$, $0$) et flèches ($\nearrow$ pour $+$, $\searrow$ pour $-$).
\item Troisième ligne : variations de $f$ (valeurs aux bornes et aux extrema) et flèches.
\end{itemize}
\end{enumerate}
\end{methode}

\begin{exoresolu}[Tableau de variations complet]
\textbf{Énoncé.} Étudier les variations de $f(x) = x^3 - 3x + 1$ sur $\mathbb{R}$.
\tcblower
\textbf{Solution.}
\textbf{1. Dérivée :} $f'(x) = 3x^2 - 3 = 3(x^2 - 1) = 3(x-1)(x+1)$.
\textbf{2. Signe de $f'$ :} Polynôme du second degré, coefficient directeur $3>0$, racines $-1$ et $1$.
$f'(x) > 0$ sur $]-\infty, -1[ \cup ]1, +\infty[$ ; $f'(x) < 0$ sur $]-1, 1[$ ; $f'(-1)=f'(1)=0$.
\textbf{3. Tableau de variations :}
\begin{center}
\begin{tabular}{|c|ccccccc|}\hline
$x$ & $-\infty$ & & $-1$ & & $1$ & & $+\infty$ \\ \hline
$f'(x)$ & & $+$ & $0$ & $-$ & $0$ & $+$ & \\ \hline
$f(x)$ & $-\infty$ & $\nearrow$ & $3$ & $\searrow$ & $-1$ & $\nearrow$ & $+\infty$ \\ \hline
\end{tabular}
\end{center}
\textbf{Conclusion :} $f$ croît sur $]-\infty, -1]$, décroît sur $[-1, 1]$, croît sur $[1, +\infty[$. Maximum local $3$ en $-1$, minimum local $-1$ en $1$.
\end{exoresolu}

\begin{attention}[Rédaction du tableau de variations au Probatoire]
\begin{itemize}
\item \textbf{3 lignes exactes} : $x$, $f'(x)$ (signe + zéros), $f(x)$ (valeurs + flèches).
\item \textbf{Colonnes :} Une pour chaque valeur remarquable (bornes, zéros de $f'$) ET une colonne \textbf{entre} deux valeurs pour le signe/flèche.
\item \textbf{Flèches UNIQUEMENT sur la ligne $f(x)$}. Sur la ligne $f'(x)$, on met seulement $+$, $-$, $0$ (ou double barre $||$ pour valeur interdite).
\item \textbf{Valeurs interdites :} Si $f$ n'est pas définie en $a$ (ex: $1/x$ en $0$), on double la colonne : $||$ dans le trait vertical, et on écrit $-\infty$ / $+\infty$ ou la limite.
\end{itemize}
\end{attention}

\coursec{Applications à des situations concrètes}

\courssub{Recherche d'optimum (maximisation / minimisation)}

\textbf{Intuition.} C'est le problème de Mama Odile : trouver le prix qui maximise la recette. Mathématiquement, un extremum (max ou min) d'une fonction dérivable se produit là où la tangente est horizontale, c'est-à-dire où la dérivée s'annule.

\begin{propriete}[Condition nécessaire d'extremum]
Si $f$ est dérivable sur $I$ et admet un extremum (maximum ou minimum) en $a \in I$ (intérieur à $I$), alors $f'(a) = 0$.
\textbf{Attention :} $f'(a)=0$ est \textbf{nécessaire} mais pas \textbf{suffisante} (ex: $f(x)=x^3$ en $0$, $f'(0)=0$ mais pas d'extremum). Il faut vérifier le \textbf{changement de signe} de $f'$ (tableau de variations).
\end{propriete}

\begin{exoresolu}[Le prix optimal des beignets de Mama Odile]
\textbf{Énoncé.} Mama Odile modélise sa recette journalière (en milliers de FCFA) en fonction du prix unitaire $p$ (en centaines de FCFA) par la fonction $R(p) = -2p^2 + 20p - 30$, pour $p \in [2; 8]$. Quel prix maximise la recette ? Quelle est cette recette maximale ?
\tcblower
\textbf{Solution.}
\textbf{1. Dérivée :} $R'(p) = -4p + 20 = -4(p - 5)$.
\textbf{2. Signe de $R'$ :} $R'(p) > 0 \iff p < 5$ ; $R'(p) < 0 \iff p > 5$ ; $R'(5)=0$.
\textbf{3. Tableau de variations sur $[2; 8]$ :}
\begin{center}
\begin{tabular}{|c|ccccc|}\hline
$p$ & $2$ & & $5$ & & $8$ \\ \hline
$R'(p)$ & & $+$ & $0$ & $-$ & \\ \hline
$R(p)$ & $R(2)=-6$ & $\nearrow$ & $R(5)=20$ & $\searrow$ & $R(8)=-14$ \\ \hline
\end{tabular}
\end{center}
\textbf{4. Conclusion :} La recette est maximale pour $p = 5$ (soit \textbf{500 FCFA} le beignet). La recette maximale est de \textbf{20 000 FCFA}.
\end{exoresolu}

\courssub{Applications physiques : vitesse et accélération}

\begin{propriete}[Cinématique]
Soit $x(t)$ la position (abscisse) d'un mobile sur une droite à l'instant $t$.
\begin{itemize}
\item La \cle{vitesse instantanée} est $v(t) = x'(t)$.
\item L'\cle{accélération} est $a(t) = v'(t) = x''(t)$ (dérivée seconde).
\item Si $v(t_0) > 0$, le mobile avance dans le sens positif ; si $v(t_0) < 0$, il recule.
\item Si $v(t_0) = 0$ et $v$ change de signe, le mobile s'arrête et repart en sens inverse (inversion de mouvement).
\end{itemize}
\end{propriete}

\begin{exemplebox}[Chute d'un objet depuis un immeuble]
Un objet est lâché sans vitesse initiale du sommet d'un immeuble de 80 m de haut. Modèle : $x(t) = 80 - 5t^2$ (position en m, $t$ en s, $g \approx 10$ m/s$^2$).
\begin{itemize}
\item Vitesse : $v(t) = x'(t) = -10t$. (Négative : l'objet descend).
\item Accélération : $a(t) = v'(t) = -10$ m/s$^2$ (constante, pesanteur).
\item Heure d'arrivée au sol : $x(t)=0 \Rightarrow 80-5t^2=0 \Rightarrow t=4$ s.
\item Vitesse à l'impact : $v(4) = -40$ m/s (soit 144 km/h).
\end{itemize}
\end{exemplebox}

\courssub{Applications économiques : coût marginal, revenu marginal}

\begin{definition}[Grandeurs marginales]
Soit $C(q)$ le coût total de production de $q$ unités, et $R(q)$ le revenu total.
\begin{itemize}
\item Le \cle{coût marginal} est $C'(q)$ : coût approximatif de production de l'unité supplémentaire (la $q+1$-ième).
\item Le \cle{revenu marginal} est $R'(q)$ : revenu apporté par la vente de l'unité supplémentaire.
\item Le \cle{bénéfice} est $B(q) = R(q) - C(q)$. Il est maximal quand $B'(q)=0 \iff R'(q) = C'(q)$ (Revenu marginal = Coût marginal).
\end{itemize}
\end{definition}

\begin{exoresolu}[Optimisation du bénéfice en atelier]
\textbf{Énoncé.} Un atelier produit $q$ pièces. Coût total : $C(q) = 0,1q^2 + 10q + 500$ (kFCFA). Prix de vente unitaire fixe : $p = 50$ kFCFA/pièce. Déterminer la production $q$ maximisant le bénéfice.
\tcblower
\textbf{Solution.}
Revenu : $R(q) = 50q$.
Bénéfice : $B(q) = R(q) - C(q) = 50q - (0,1q^2 + 10q + 500) = -0,1q^2 + 40q - 500$.
Dérivée : $B'(q) = -0,2q + 40 = -0,2(q - 200)$.
$B'(q) > 0$ pour $q < 200$, $B'(q) < 0$ pour $q > 200$.
Maximum pour $q = 200$ pièces.
Bénéfice max : $B(200) = -0,1(40000) + 8000 - 500 = 3500$ kFCFA = \textbf{3 500 000 FCFA}.
\end{exoresolu}

\begin{aretenir}[Synthèse de la leçon : Dérivation]
\begin{itemize}
\item \textbf{Nombre dérivé $f'(a)$} : limite du taux d'accroissement. Pente de la tangente. Vitesse instantanée.
\item \textbf{Dérivabilité} $\iff$ $f'_g(a) = f'_d(a)$ (finis). Dérivable $\Rightarrow$ Continue.
\item \textbf{Fonction dérivée $f'$} : calculée via tableau des dérivées usuelles + règles de dérivation (somme, produit, quotient, chaîne affine $a \cdot u'(ax+b)$).
\item \textbf{Tangente en $a$} : $y = f'(a)(x-a) + f(a)$.
\item \textbf{Approximation affine} : $f(a+h) \approx f(a) + h f'(a)$.
\item \textbf{Variations} : Signe de $f'$ $\to$ Tableau de variations $\to$ Extrema.
\item \textbf{Applications} : Optimisation (annulation de $f'$ + changement de signe), Cinématique ($v=x', a=v'$), Économie ($C', R', B'=R'-C'$).
\end{itemize}
\end{aretenir}

\begin{savaistu}[La dérivée seconde et la concavité (Aperçu Terminale)]
La dérivée de la fonction dérivée $f'$ s'appelle la \emph{dérivée seconde} et se note $f''$. Elle mesure la variation de la pente : si $f'' > 0$, la courbe est \emph{convexe} (elle « tient l'eau », comme $x^2$) ; si $f'' < 0$, elle est \emph{concave} (elle « fuit », comme $-x^2$). Un point où $f''$ change de signe est un \emph{point d'inflexion} (la tangente coupe la courbe). Cela permet d'affiner le tracé des courbes en Terminale.
\end{savaistu}

\begin{exercice}[Exercices de synthèse - Leçon complète]
\begin{enumerate}
\item \textbf{Dérivabilité par la définition.} Soit $f(x) = \dfrac{1}{x+1}$. Montrer par la définition que $f$ est dérivable sur $]-1, +\infty[$ et donner $f'(x)$.
\item \textbf{Règles de dérivation.} Dériver sur leur domaine :
\begin{enumerate}
\item $f(x) = (x^2+3x)(2x-1)$
\item $g(x) = \dfrac{x^2-1}{x+2}$
\item $h(x) = \sqrt{3x-2}$
\item $k(x) = \dfrac{1}{(2x+1)^2}$
\end{enumerate}
\item \textbf{Tangente et approximation.} Soit $f(x) = x^3 - 2x$.
\begin{enumerate}
\item Trouver l'équation de la tangente $\mathcal{T}$ en $a=1$.
\item Par approximation affine, donner une valeur approchée de $f(1,02)$.
\end{enumerate}
\item \textbf{Variations et optimisation.} Soit $f(x) = \dfrac{x}{x^2+1}$ sur $\mathbb{R}$.
\begin{enumerate}
\item Calculer $f'(x)$ et étudier son signe.
\item Dresser le tableau de variations de $f$.
\item Quels sont le maximum et le minimum de $f$ sur $\mathbb{R}$ ?
\end{enumerate}
\item \textbf{Problème concret : Cylindre de volume maximal.} On veut fabriquer une boîte cylindrique (avec couvercle) de volume $V = 1000$ cm$^3$. Soit $r$ le rayon de la base et $h$ la hauteur. La surface totale de tôle utilisée est $S(r) = 2\pi r^2 + 2\pi r h$.
\begin{enumerate}
\item Exprimer $h$ en fonction de $r$ et en déduire $S(r)$.
\item Étudier les variations de $S$ sur $]0, +\infty[$.
\item Quel rayon $r$ minimise la quantité de tôle ? Donner la hauteur correspondante.
\end{enumerate}
\end{enumerate}
\end{exercice}

\begin{corrige}[Exercices de synthèse]
\textbf{1.} $f(x) = \frac{1}{x+1}$, $a > -1$. $\tau(h) = \frac{\frac{1}{a+h+1} - \frac{1}{a+1}}{h} = \frac{a+1 - (a+h+1)}{h(a+h+1)(a+1)} = \frac{-h}{h(a+h+1)(a+1)} = \frac{-1}{(a+h+1)(a+1)} \xrightarrow[h\to 0]{} \frac{-1}{(a+1)^2}$. Donc $f'(x) = -\frac{1}{(x+1)^2}$ sur $]-1, +\infty[$.

\textbf{2.}
\begin{enumerate}
\item $f'(x) = (2x+3)(2x-1) + (x^2+3x)\cdot 2 = 4x^2+4x-3 + 2x^2+6x = 6x^2+10x-3$.
\item $g'(x) = \frac{2x(x+2) - (x^2-1)\cdot 1}{(x+2)^2} = \frac{2x^2+4x - x^2 + 1}{(x+2)^2} = \frac{x^2+4x+1}{(x+2)^2}$.
\item $h'(x) = 3 \cdot \frac{1}{2\sqrt{3x-2}} = \frac{3}{2\sqrt{3x-2}}$ (domaine $x > 2/3$).
\item $k(x) = (2x+1)^{-2}$. $k'(x) = 2 \cdot (-2)(2x+1)^{-3} = -\frac{4}{(2x+1)^3}$ (domaine $x \neq -1/2$).
\end{enumerate}

\textbf{3.}
\begin{enumerate}
\item $f'(x) = 3x^2 - 2$. $f'(1) = 1$. $f(1) = -1$. Tangente : $y = 1(x-1) - 1 = x - 2$.
\item $f(1,02) \approx f(1) + 0,02 f'(1) = -1 + 0,02 \times 1 = -0,98$. (Exact : $-0,97979...$)
\end{enumerate}

\textbf{4.}
\begin{enumerate}
\item $f'(x) = \frac{1\cdot(x^2+1) - x\cdot 2x}{(x^2+1)^2} = \frac{1 - x^2}{(x^2+1)^2} = \frac{(1-x)(1+x)}{(x^2+1)^2}$.
\item Signe de $f'$ : $+$ sur $]-1, 1[$, $-$ sur $]-\infty, -1[ \cup ]1, +\infty[$. Zéros en $-1$ et $1$.
\item Tableau :
\begin{center}
\begin{tabular}{|c|ccccccc|}\hline
$x$ & $-\infty$ & & $-1$ & & $1$ & & $+\infty$ \\ \hline
$f'(x)$ & & $-$ & $0$ & $+$ & $0$ & $-$ & \\ \hline
$f(x)$ & $0^-$ & $\searrow$ & $-\frac{1}{2}$ & $\nearrow$ & $\frac{1}{2}$ & $\searrow$ & $0^+$ \\ \hline
\end{tabular}
\end{center}
\item Minimum global $-\frac{1}{2}$ en $-1$. Maximum global $\frac{1}{2}$ en $1$.
\end{enumerate}

\textbf{5.}
\begin{enumerate}
\item $V = \pi r^2 h = 1000 \Rightarrow h = \frac{1000}{\pi r^2}$.
$S(r) = 2\pi r^2 + 2\pi r \frac{1000}{\pi r^2} = 2\pi r^2 + \frac{2000}{r}$.
\item $S'(r) = 4\pi r - \frac{2000}{r^2} = \frac{4\pi r^3 - 2000}{r^2}$.
$S'(r) = 0 \iff 4\pi r^3 = 2000 \iff r^3 = \frac{500}{\pi} \iff r = \sqrt[3]{\frac{500}{\pi}} \approx 5,42$ cm.
Signe : $S'(r) < 0$ pour $r < r_0$, $>0$ pour $r > r_0$. Minimum en $r_0$.
\item $r \approx 5,42$ cm. $h = \frac{1000}{\pi r^2} = \frac{1000}{\pi (500/\pi)^{2/3}} = 2 \sqrt[3]{\frac{500}{\pi}} = 2r \approx 10,84$ cm. (La hauteur vaut le diamètre pour le minimum de surface).
\end{enumerate}
\end{corrige}

\coursec{Fonction dérivée sur un intervalle}

Dans la section précédente, nous avons défini le \cle{nombre dérivé} d'une fonction $f$ en un réel $a$ : c'est la limite du taux d'accroissement $\dfrac{f(a+h)-f(a)}{h}$ lorsque $h$ tend vers $0$, lorsque cette limite existe et est finie. Ce nombre $f'(a)$ donne la pente de la tangente à la courbe au point d'abscisse $a$.

Mais voici une idée qui va tout changer : au lieu de calculer le nombre dérivé point par point, un par un, pourquoi ne pas le calculer \emph{en une seule fois pour tous les points} ? Si une fonction est dérivable partout sur un intervalle, chaque réel $x$ de cet intervalle possède son propre nombre dérivé $f'(x)$. On obtient ainsi une \emph{nouvelle fonction}, qui à chaque $x$ associe la pente de la tangente en $x$. C'est cette nouvelle fonction, appelée \cle{fonction dérivée}, que nous étudions dans cette section. Elle est l'outil central de toute l'analyse : c'est elle qui nous permettra, dans la section 5, de connaître le sens de variation d'une fonction sans tracer sa courbe.

\courssub{Dérivabilité sur un intervalle}

\textbf{Intuition.} Dire qu'une fonction est dérivable en un point, c'est dire que sa courbe admet en ce point une tangente non verticale : localement, la courbe « ressemble » à une droite. Dire qu'elle est dérivable sur un intervalle, c'est dire que ce phénomène a lieu \emph{en tout point} de l'intervalle : la courbe est « lisse », sans point anguleux ni tangente verticale, sur tout cet intervalle.

\begin{definition}[Fonction dérivable sur un intervalle]
Soit $f$ une fonction définie sur un intervalle $I$.
\begin{itemize}
\item On dit que $f$ est \cle{dérivable sur l'intervalle ouvert} $I$ si $f$ est dérivable en \textbf{tout} réel $x$ de $I$.
\item Si $I = [a ; b]$ est un intervalle fermé, on dit que $f$ est dérivable sur $[a ; b]$ si $f$ est dérivable en tout point de $]a ; b[$, dérivable \textbf{à droite} en $a$ et dérivable \textbf{à gauche} en $b$.
\end{itemize}
\end{definition}

\begin{attention}[Aux bornes, on ne regarde que d'un seul côté]
En une borne de l'intervalle, la fonction n'est définie que d'un côté : on ne peut donc exiger que la dérivabilité \emph{du côté où la fonction existe}. Par exemple, la fonction racine carrée $x \mapsto \sqrt{x}$ est définie sur $[0 ; +\infty[$ ; on ne pourra parler en $0$ que de dérivabilité \textbf{à droite} (et nous verrons qu'elle échoue !).
\end{attention}

\begin{exemplebox}[Vérifier une dérivabilité sur un intervalle]
\begin{itemize}
\item La fonction $f : x \mapsto x^2$ est dérivable en tout réel, avec $f'(x) = 2x$ : elle est dérivable sur $\mathbb{R}$.
\item La fonction $g : x \mapsto \dfrac{1}{x}$ est dérivable en tout réel non nul : elle est dérivable sur $]-\infty ; 0[$ et sur $]0 ; +\infty[$, mais \textbf{pas} en $0$ (elle n'y est même pas définie).
\item La fonction $h : x \mapsto |x|$ est dérivable sur $]-\infty ; 0[$ et sur $]0 ; +\infty[$, mais pas en $0$ (point anguleux : dérivée à gauche $-1$, dérivée à droite $+1$).
\end{itemize}
\end{exemplebox}

\courssub{La fonction dérivée $f'$}

Lorsque $f$ est dérivable sur un intervalle $I$, chaque $x$ de $I$ possède un nombre dérivé $f'(x)$. On fabrique donc une nouvelle fonction.

\begin{definition}[Fonction dérivée]
Soit $f$ une fonction dérivable sur un intervalle $I$. La \cle{fonction dérivée} de $f$ (ou simplement la \cle{dérivée} de $f$) est la fonction, notée $f'$, définie sur $I$ par :
\[ f' : x \longmapsto f'(x) = \lim_{h \to 0} \frac{f(x+h)-f(x)}{h}. \]
L'ensemble des réels où $f$ est dérivable s'appelle l'\cle{ensemble de dérivabilité} de $f$.
\end{definition}

\begin{aretenir}[Deux objets à ne pas confondre]
\begin{itemize}
\item $f'(a)$ est un \textbf{nombre} : la pente de la tangente à la courbe de $f$ au point d'abscisse $a$.
\item $f'$ est une \textbf{fonction} : elle donne la pente de la tangente en \emph{n'importe quel} point $x$.
\end{itemize}
Connaître la fonction $f'$, c'est connaître d'un coup tous les nombres dérivés : il suffit ensuite de remplacer $x$ par la valeur voulue.
\end{aretenir}

\begin{exemplebox}[Du nombre dérivé à la fonction dérivée]
Soit $f(x) = x^2$. Pour un réel $x$ quelconque et $h \neq 0$ :
\[ \frac{f(x+h)-f(x)}{h} = \frac{(x+h)^2 - x^2}{h} = \frac{2xh + h^2}{h} = 2x + h \xrightarrow[h \to 0]{} 2x. \]
Donc $f$ est dérivable sur $\mathbb{R}$ et sa fonction dérivée est $f' : x \mapsto 2x$. On retrouve alors instantanément : $f'(1) = 2$, $f'(3) = 6$, $f'(-2) = -4$, sans refaire aucun calcul de limite.
\end{exemplebox}

\courssub{Dérivées des fonctions usuelles}

Refaire un calcul de limite à chaque fois serait épuisant. Les mathématiciens l'ont fait une fois pour toutes pour les fonctions de base : il faut \textbf{connaître par cœur} le tableau suivant.

\begin{propriete}[Tableau des dérivées usuelles — à connaître par cœur]
\begin{center}
\renewcommand{\arraystretch}{1.6}
\begin{tabularx}{\linewidth}{|c|X|c|c|}
\hline
\rowcolor{popblueL} \textbf{Fonction $f(x)$} & \textbf{Dérivée $f'(x)$} & \textbf{$f$ définie sur} & \textbf{$f$ dérivable sur} \\
\hline
$k$ (constante) & $0$ & $\mathbb{R}$ & $\mathbb{R}$ \\
\hline
$x$ & $1$ & $\mathbb{R}$ & $\mathbb{R}$ \\
\hline
$x^n$ \ ($n \in \mathbb{N}^*$) & $n\,x^{\,n-1}$ & $\mathbb{R}$ & $\mathbb{R}$ \\
\hline
$\dfrac{1}{x}$ & $-\dfrac{1}{x^2}$ & $\mathbb{R}^*$ & $\mathbb{R}^*$ \\
\hline
$\sqrt{x}$ & $\dfrac{1}{2\sqrt{x}}$ & $[0 ; +\infty[$ & $]0 ; +\infty[$ \\
\hline
$\sin x$ & $\cos x$ & $\mathbb{R}$ & $\mathbb{R}$ \\
\hline
$\cos x$ & $-\sin x$ & $\mathbb{R}$ & $\mathbb{R}$ \\
\hline
\end{tabularx}
\end{center}
\end{propriete}

\begin{attention}[La racine carrée n'est pas dérivable en 0]
Erreur très fréquente au Probatoire : écrire que $\sqrt{x}$ est dérivable sur $[0 ; +\infty[$. C'est \textbf{faux} ! Regardons le taux d'accroissement en $0$ :
\[ \frac{\sqrt{0+h}-\sqrt{0}}{h} = \frac{\sqrt{h}}{h} = \frac{1}{\sqrt{h}} \xrightarrow[h \to 0^+]{} +\infty. \]
La limite est infinie, donc le nombre dérivé en $0$ n'existe pas. Graphiquement, la courbe de $x \mapsto \sqrt{x}$ admet en $0$ une \textbf{demi-tangente verticale}. Conséquence pratique : la formule $\left(\sqrt{x}\right)' = \dfrac{1}{2\sqrt{x}}$ n'est valable que pour $x > 0$ (d'ailleurs, pour $x = 0$, le dénominateur serait nul : la formule elle-même « refuse » $0$ !).
\end{attention}

\courssub{Démonstrations des formules usuelles}

Ces démonstrations sont formatrices : elles montrent comment on passe du taux d'accroissement à la formule. La première est classiquement exigible.

\begin{experience}[Démonstration guidée : dérivée de $x \mapsto x^n$ pour $n \in \mathbb{N}^*$]
\textbf{Théorème.} Pour tout entier naturel $n \geq 1$, la fonction $f : x \mapsto x^n$ est dérivable sur $\mathbb{R}$ et $f'(x) = n\,x^{\,n-1}$.

\textbf{Étape 1 — Un outil algébrique.} Pour tous réels $u$ et $v$ et tout entier $n \geq 1$, on a la factorisation :
\[ u^n - v^n = (u - v)\left(u^{n-1} + u^{n-2}v + u^{n-3}v^2 + \dots + u\,v^{n-2} + v^{n-1}\right). \]
Vérifions-la pour $n = 3$ : $(u-v)(u^2 + uv + v^2) = u^3 + u^2v + uv^2 - u^2v - uv^2 - v^3 = u^3 - v^3$. L'égalité est vérifiée. La grande parenthèse contient exactement $n$ termes.

\textbf{Étape 2 — Taux d'accroissement.} Fixons un réel $a$ et prenons $h \neq 0$. Appliquons la factorisation avec $u = a+h$ et $v = a$ :
\[ \frac{f(a+h)-f(a)}{h} = \frac{(a+h)^n - a^n}{h} = \frac{h \times \big[(a+h)^{n-1} + (a+h)^{n-2}a + \dots + a^{n-1}\big]}{h}. \]

\textbf{Étape 3 — Simplification et limite.} On simplifie par $h$ :
\[ \frac{f(a+h)-f(a)}{h} = (a+h)^{n-1} + (a+h)^{n-2}a + \dots + a^{n-1}. \]
Quand $h \to 0$, chacun des $n$ termes tend vers $a^{n-1}$. La somme tend donc vers $n \times a^{n-1}$.

\textbf{Conclusion.} $f$ est dérivable en tout réel $a$ et $f'(a) = n\,a^{\,n-1}$, c'est-à-dire $\left(x^n\right)' = n\,x^{\,n-1}$. $\blacksquare$
\end{experience}

\begin{exemplebox}[Application immédiate]
\begin{itemize}
\item $\left(x^3\right)' = 3x^2$ ; \quad $\left(x^5\right)' = 5x^4$ ; \quad $\left(x^{10}\right)' = 10x^9$.
\item Pour $n = 1$ : $(x)' = 1 \times x^0 = 1$. On retrouve la ligne « $x$ » du tableau.
\end{itemize}
\end{exemplebox}

\begin{experience}[Démonstration : dérivée de $x \mapsto \dfrac{1}{x}$]
Soit $f(x) = \dfrac{1}{x}$, définie sur $\mathbb{R}^*$. Fixons $a \neq 0$ et $h$ non nul assez petit pour que $a + h \neq 0$ :
\begin{align*}
\frac{f(a+h)-f(a)}{h} &= \frac{1}{h}\left(\frac{1}{a+h} - \frac{1}{a}\right) = \frac{1}{h} \times \frac{a - (a+h)}{a(a+h)} \\
&= \frac{1}{h} \times \frac{-h}{a(a+h)} = \frac{-1}{a(a+h)}.
\end{align*}
Quand $h \to 0$, le dénominateur $a(a+h)$ tend vers $a^2$, donc le taux tend vers $-\dfrac{1}{a^2}$. Ainsi $f$ est dérivable en tout $a \neq 0$ et :
\[ \left(\frac{1}{x}\right)' = -\frac{1}{x^2}. \quad \blacksquare \]
\end{experience}

\begin{experience}[Démonstration : dérivée de $x \mapsto \sqrt{x}$]
Soit $f(x) = \sqrt{x}$ et $a > 0$. Pour $h \neq 0$ tel que $a + h \geq 0$, l'astuce consiste à multiplier par la \cle{quantité conjuguée} :
\begin{align*}
\frac{\sqrt{a+h}-\sqrt{a}}{h} &= \frac{\left(\sqrt{a+h}-\sqrt{a}\right)\left(\sqrt{a+h}+\sqrt{a}\right)}{h\left(\sqrt{a+h}+\sqrt{a}\right)} \\
&= \frac{(a+h) - a}{h\left(\sqrt{a+h}+\sqrt{a}\right)} = \frac{h}{h\left(\sqrt{a+h}+\sqrt{a}\right)} = \frac{1}{\sqrt{a+h}+\sqrt{a}}.
\end{align*}
Quand $h \to 0$, le dénominateur tend vers $2\sqrt{a}$ (non nul car $a > 0$). Donc $f$ est dérivable en tout $a > 0$ et :
\[ \left(\sqrt{x}\right)' = \frac{1}{2\sqrt{x}}. \quad \blacksquare \]
Remarquez que cette démonstration exige $a > 0$ : en $a = 0$, le dénominateur tendrait vers $0$ et le taux divergerait, ce qui confirme l'attention signalée plus haut.
\end{experience}

\courssub{Visualiser le lien entre $f$ et $f'$}

La figure suivante montre la fonction $f(x) = x^3$ (en haut) et sa dérivée $f'(x) = 3x^2$ (en bas). Observez le lien point par point : là où la courbe de $f$ est « plate » (en $x = 0$), la dérivée vaut $0$ ; là où $f$ monte, la dérivée est positive ; plus la montée est raide, plus $f'(x)$ est grand.

\begin{popfigure}
\begin{center}
\begin{tikzpicture}
\begin{axis}[
  width=0.9\linewidth, height=5.2cm,
  axis lines=middle, xlabel={$x$}, ylabel={$f(x)=x^3$},
  domain=-1.6:1.6, samples=100,
  xtick={-1,0,1}, ytick=\empty,
  every axis x label/.style={at={(ticklabel* cs:1)},anchor=west},
]
\addplot[popcurve] {x^3};
\end{axis}
\end{tikzpicture}

\vspace{0.4cm}

\begin{tikzpicture}
\begin{axis}[
  width=0.9\linewidth, height=5.2cm,
  axis lines=middle, xlabel={$x$}, ylabel={$f'(x)=3x^2$},
  domain=-1.6:1.6, samples=100,
  xtick={-1,0,1}, ytick={3},
  every axis x label/.style={at={(ticklabel* cs:1)},anchor=west},
]
\addplot[popcurve,poporange] {3*x^2};
\end{axis}
\end{tikzpicture}
\end{center}
\legende{En haut : la courbe de $f(x) = x^3$. En bas : la courbe de sa dérivée $f'(x) = 3x^2$. La dérivée s'annule là où la tangente est horizontale ($x = 0$) et est positive partout ailleurs, car $f$ est croissante.}
\end{popfigure}

\begin{exoresolu}[Calculer une fonction dérivée et l'exploiter]
\textbf{Énoncé.} Un atelier de menuiserie à Douala découpe des panneaux carrés de côté $x$ (en mètres). L'aire d'un panneau est $A(x) = x^2$ et le coût de fabrication d'un lot est modélisé par $C(x) = x^3$ (en milliers de francs CFA) pour $x \in [0 ; 3]$.
\begin{enumerate}
\item Justifier que $A$ et $C$ sont dérivables sur $[0 ; 3]$ et donner leurs fonctions dérivées.
\item Calculer $A'(2)$ et $C'(2)$. Interpréter $C'(2)$ pour le menuisier.
\end{enumerate}
\tcblower
\textbf{Solution détaillée.}
\begin{enumerate}
\item $A(x) = x^2$ et $C(x) = x^3$ sont des fonctions puissances $x^n$ avec $n \in \mathbb{N}^*$ : elles sont dérivables sur $\mathbb{R}$, donc en particulier sur $[0 ; 3]$ (dérivabilité à droite en $0$, à gauche en $3$). D'après la formule $\left(x^n\right)' = n x^{n-1}$ :
\[ A'(x) = 2x \qquad \text{et} \qquad C'(x) = 3x^2. \]
\item $A'(2) = 2 \times 2 = 4$ et $C'(2) = 3 \times 2^2 = 12$.

\textbf{Interprétation.} $C'(2) = 12$ signifie qu'autour de $x = 2$ m, augmenter le côté du panneau d'une petite quantité fait augmenter le coût d'environ $12$ fois cette quantité (en milliers de francs) : par exemple, passer de $2$ m à $2{,}1$ m augmente le coût d'environ $12 \times 0{,}1 = 1{,}2$ millier de francs, soit environ 1200 F CFA. C'est le \emph{coût marginal} du menuisier.
\end{enumerate}
\end{exoresolu}

\begin{attention}[Pièges fréquents sur les dérivées usuelles]
\begin{itemize}
\item La dérivée d'une constante est $0$, pas la constante elle-même : $(5)' = 0$, $(\pi)' = 0$.
\item $\left(\dfrac{1}{x}\right)' = -\dfrac{1}{x^2}$ : n'oubliez pas le signe moins, et n'écrivez pas $\dfrac{1}{x^2}$ tout court.
\item $\left(\sqrt{x}\right)' = \dfrac{1}{2\sqrt{x}}$, et non $\dfrac{1}{\sqrt{x}}$ ni $\dfrac{1}{2x}$.
\item $(\cos x)' = -\sin x$ : c'est $\cos$ qui « prend le moins », pas $\sin$.
\item Toujours préciser l'ensemble de dérivabilité : écrire « $f'(x) = \dfrac{1}{2\sqrt{x}}$ » sans dire que $x > 0$ est une rédaction incomplète.
\end{itemize}
\end{attention}

\begin{savaistu}[Lagrange, Leibniz et les notations de la dérivée]
La notation $f'(x)$ que nous utilisons est due au mathématicien franco-italien \textbf{Joseph-Louis Lagrange} (1736--1813), qui parlait de « fonction dérivée » — c'est de lui que vient le mot même ! De son côté, le philosophe et mathématicien allemand \textbf{Gottfried Wilhelm Leibniz} (1646--1716), co-inventeur du calcul différentiel avec Newton, écrivait la dérivée $\dfrac{\mathrm{d}f}{\mathrm{d}x}$, comme un « quotient de toutes petites variations ». Cette notation de Leibniz est encore massivement utilisée aujourd'hui en physique et en électrotechnique : quand un technicien écrit que l'intensité est $i = \dfrac{\mathrm{d}q}{\mathrm{d}t}$, il dit que le courant est la dérivée de la charge électrique par rapport au temps. Deux notations, une même idée géniale : mesurer comment une grandeur change instantanément.
\end{savaistu}

\begin{savaistu}[Pour aller plus loin — hors exigibilité au Probatoire]
La formule $\left(x^n\right)' = n\,x^{\,n-1}$ reste valable pour un entier $n$ \textbf{négatif}, à condition de se placer sur $\mathbb{R}^*$. Par exemple, avec $n = -1$ : $x^{-1} = \dfrac{1}{x}$ et la formule donne $\left(x^{-1}\right)' = -1 \times x^{-2} = -\dfrac{1}{x^2}$ : on retrouve exactement la dérivée de $\dfrac{1}{x}$ démontrée plus haut ! De même, $\left(\dfrac{1}{x^2}\right)' = \left(x^{-2}\right)' = -2x^{-3} = -\dfrac{2}{x^3}$. Cette généralisation sera très utile en classe de Terminale technique.
\end{savaistu}

\begin{exercice}[À faire]
\begin{enumerate}
\item Donner la fonction dérivée de chacune des fonctions suivantes, en précisant l'ensemble de dérivabilité : $f(x) = x^7$ ; \ $g(x) = \dfrac{1}{x}$ ; \ $h(x) = \sqrt{x}$ ; \ $k(x) = \sin x$.
\item Soit $f(x) = x^4$. Calculer $f'(x)$, puis $f'(1)$, $f'(-1)$ et $f'(0)$.
\item La fonction $x \mapsto \sqrt{x}$ est-elle dérivable sur $[0 ; +\infty[$ ? Justifier.
\end{enumerate}
\end{exercice}

\begin{corrige}[Exercice n°1]
\begin{enumerate}
\item $f'(x) = 7x^6$ sur $\mathbb{R}$ ; \ $g'(x) = -\dfrac{1}{x^2}$ sur $\mathbb{R}^*$ ; \ $h'(x) = \dfrac{1}{2\sqrt{x}}$ sur $]0 ; +\infty[$ ; \ $k'(x) = \cos x$ sur $\mathbb{R}$.
\item $f'(x) = 4x^3$ ; donc $f'(1) = 4$, $f'(-1) = -4$, $f'(0) = 0$.
\item Non : en $0$, le taux d'accroissement $\dfrac{\sqrt{h}}{h} = \dfrac{1}{\sqrt{h}}$ tend vers $+\infty$ quand $h \to 0^+$. La limite n'est pas finie, donc $\sqrt{x}$ n'est pas dérivable en $0$ (demi-tangente verticale). Elle est dérivable seulement sur $]0 ; +\infty[$.
\end{enumerate}
\end{corrige}

Nous savons désormais dériver toutes les fonctions usuelles. Mais en pratique, les fonctions rencontrées sont des \emph{mélanges} de ces fonctions de base : sommes, produits, quotients, compositions avec des fonctions affines. C'est l'objet de la section suivante : les opérations sur les fonctions dérivables.

\coursec{Opérations sur les fonctions dérivables}

\courssub{Transition et rappel}

Dans la section précédente, nous avons défini la \cle{fonction dérivée} $f'$ sur un intervalle $I$ lorsque $f$ est dérivable en tout point de $I$. Nous savons maintenant dériver les fonctions usuelles (puissances entières, racine carrée, fonctions trigonométriques de base) grâce au tableau de dérivées. Cependant, les fonctions rencontrées dans les problèmes techniques (électricité, mécanique, gestion) sont souvent des \emph{combinaisons} de ces fonctions de base : sommes, produits, quotients, compositions. Cette section établit les règles algébriques qui permettent de dériver ces combinaisons sans revenir à la définition par la limite du taux d'accroissement à chaque fois.

\courssub{Dérivée d'une somme et produit par un réel}

L'intuition est simple : la variation d'une somme est la somme des variations, et multiplier une fonction par une constante $k$ multiplie sa variation par $k$.

\begin{propriete}[Linéarité de la dérivation]
Soient $u$ et $v$ deux fonctions dérivables en un point $a$ (ou sur un intervalle $I$), et $k \in \mathbb{R}$.
\begin{enumerate}
    \item \textbf{Somme :} $u+v$ est dérivable en $a$ et $(u+v)'(a) = u'(a) + v'(a)$.
    \item \textbf{Produit par un réel :} $ku$ est dérivable en $a$ et $(ku)'(a) = k \cdot u'(a)$.
\end{enumerate}
\end{propriete}

\begin{experience}[Démonstration par les taux d'accroissement]
On utilise la définition du nombre dérivé comme limite du taux d'accroissement.
Pour $h \neq 0$ tel que $a+h$ appartienne au domaine de définition :
\[
\frac{(u+v)(a+h) - (u+v)(a)}{h} = \frac{u(a+h)-u(a)}{h} + \frac{v(a+h)-v(a)}{h}.
\]
En faisant tendre $h$ vers $0$, la limite de la somme est la somme des limites (théorème de la limite d'une somme), d'où $(u+v)'(a) = u'(a)+v'(a)$.
De même pour $ku$ :
\[
\frac{ku(a+h)-ku(a)}{h} = k \cdot \frac{u(a+h)-u(a)}{h} \xrightarrow[h\to 0]{} k \cdot u'(a).
\]
\end{experience}

\begin{exemplebox}[Dérivation d'un polynôme]
Soit $f(x) = 3x^4 - 5x^2 + 2x - 7$ sur $\mathbb{R}$.
Par linéarité et tableau de dérivées ($x^n \mapsto nx^{n-1}$) :
\[
f'(x) = 3(4x^3) - 5(2x) + 2(1) - 0 = 12x^3 - 10x + 2.
\]
\end{exemplebox}

\courssub{Dérivée d'un produit}

Contrairement à la somme, la dérivée d'un produit n'est \textbf{pas} le produit des dérivées. La variation d'un produit $uv$ dépend de la variation de $u$ \emph{et} de la valeur de $v$, et inversement.

\begin{propriete}[Formule du produit]
Soient $u$ et $v$ deux fonctions dérivables en $a$ (ou sur $I$). La fonction $uv$ est dérivable en $a$ et :
\[
(uv)'(a) = u'(a)v(a) + u(a)v'(a).
\]
\end{propriete}

\begin{experience}[Démonstration guidée : l'astuce de l'ajout-retrait]
On part du taux d'accroissement de $uv$ en $a$ :
\[
\frac{u(a+h)v(a+h) - u(a)v(a)}{h}.
\]
On ajoute et on retire le terme $u(a+h)v(a)$ au numérateur (c'est l'\cle{ajout-retrait}) :
\begin{align*}
\frac{u(a+h)v(a+h) - u(a+h)v(a) + u(a+h)v(a) - u(a)v(a)}{h}
&= u(a+h)\frac{v(a+h)-v(a)}{h} + v(a)\frac{u(a+h)-u(a)}{h}.
\end{align*}
Comme $u$ est dérivable en $a$, elle est continue en $a$, donc $\lim_{h\to 0} u(a+h) = u(a)$.
En passant à la limite $h \to 0$, on obtient bien $u(a)v'(a) + u'(a)v(a)$.
\end{experience}

\begin{attention}[Piège classique : $(uv)' \neq u'v'$]
Ne jamais écrire $(uv)' = u'v'$. C'est faux !
\textbf{Exemple :} $u(x)=x$, $v(x)=x$. $(uv)' = (x^2)' = 2x$. Mais $u'v' = 1 \times 1 = 1 \neq 2x$.
\emph{Moyen mnémotechnique :} « Dérive le premier, laisse le second + laisse le premier, dérive le second ».
\end{attention}

\begin{exemplebox}[Produit d'un polynôme et d'une racine]
Soit $f(x) = (x^2+1)\sqrt{x}$ sur $[0; +\infty[$.
On pose $u(x)=x^2+1$, $v(x)=\sqrt{x}$.
$u'(x)=2x$, $v'(x)=\dfrac{1}{2\sqrt{x}}$ (pour $x>0$).
\[
f'(x) = u'(x)v(x) + u(x)v'(x) = 2x\sqrt{x} + (x^2+1)\frac{1}{2\sqrt{x}}.
\]
On peut simplifier en mettant au même dénominateur $2\sqrt{x}$ :
\[
f'(x) = \frac{4x^2 + x^2 + 1}{2\sqrt{x}} = \frac{5x^2+1}{2\sqrt{x}} \quad (x>0).
\]
\end{exemplebox}

\courssub{Dérivée de l'inverse et d'un quotient}

\begin{propriete}[Dérivée de l'inverse]
Soit $v$ dérivable en $a$ avec $v(a) \neq 0$. La fonction $\frac{1}{v}$ est dérivable en $a$ et :
\[
\left(\frac{1}{v}\right)'(a) = -\frac{v'(a)}{v(a)^2}.
\]
\end{propriete}

\begin{experience}[Démonstration guidée]
\[
\frac{\frac{1}{v(a+h)} - \frac{1}{v(a)}}{h} = \frac{v(a) - v(a+h)}{h \cdot v(a+h)v(a)} = -\frac{v(a+h)-v(a)}{h} \cdot \frac{1}{v(a+h)v(a)}.
\]
Par continuité de $v$ en $a$, $v(a+h) \to v(a)$. La limite donne $-\dfrac{v'(a)}{v(a)^2}$.
\end{experience}

\begin{attention}[Erreur fréquente : $(1/v)' \neq 1/v'$]
La dérivée de l'inverse n'est pas l'inverse de la dérivée. Le signe moins et le carré au dénominateur sont essentiels.
\end{attention}

\begin{propriete}[Dérivée d'un quotient]
Soient $u$ et $v$ dérivables en $a$ avec $v(a) \neq 0$. La fonction $\frac{u}{v}$ est dérivable en $a$ et :
\[
\left(\frac{u}{v}\right)'(a) = \frac{u'(a)v(a) - u(a)v'(a)}{v(a)^2}.
\]
\end{propriete}

\begin{experience}[Démonstration comme conséquence]
On écrit $\frac{u}{v} = u \cdot \frac{1}{v}$. On applique la formule du produit et la formule de l'inverse :
\[
\left(u \cdot \frac{1}{v}\right)' = u' \cdot \frac{1}{v} + u \cdot \left(-\frac{v'}{v^2}\right) = \frac{u'v - uv'}{v^2}.
\]
\end{experience}

\begin{exoresolu}[Dérivation d'une fraction rationnelle]
Soit $f(x) = \frac{2x-1}{x+3}$ sur $]-3; +\infty[$.
\emph{Identifier $u$ et $v$ :} $u(x)=2x-1$, $v(x)=x+3$.
\emph{Calculer $u'$ et $v'$ :} $u'(x)=2$, $v'(x)=1$.
\emph{Appliquer la formule :}
\[
f'(x) = \frac{2(x+3) - (2x-1)\cdot 1}{(x+3)^2} = \frac{2x+6 - 2x + 1}{(x+3)^2} = \frac{7}{(x+3)^2}.
\]
\emph{Conclusion :} $f'(x) = \frac{7}{(x+3)^2} > 0$ sur $]-3; +\infty[$, donc $f$ est strictement croissante sur cet intervalle.
\end{exoresolu}

\courssub{Dérivée de la composée par une fonction affine}

C'est un cas particulier très fréquent de la dérivation des fonctions composées. Ici, la fonction intérieure est une fonction affine $g(x) = ax+b$.

\begin{propriete}[Dérivée de $f(ax+b)$]
Soit $f$ dérivable sur un intervalle $J$ et $g(x)=ax+b$ une fonction affine ($a,b \in \mathbb{R}$, $a \neq 0$) telle que $g(x) \in J$ pour $x \in I$. La fonction composée $f \circ g : x \mapsto f(ax+b)$ est dérivable sur $I$ et :
\[
(f(ax+b))' = a \cdot f'(ax+b).
\]
\end{propriete}

\begin{popfigure}
\begin{tikzpicture}[
    node distance=1.5cm and 2.5cm,
    every node/.style={font=\small},
    box/.style={draw=popblue, thick, rounded corners, fill=popblue!10, minimum width=2.5cm, minimum height=1cm},
    arrow/.style={-Stealth, thick, popdark}
]
\node[box] (x) {$x$};
\node[box, right=of x] (affine) {$ax+b$};
\node[box, right=of affine] (f) {$f$};
\node[box, right=of f] (y) {$f(ax+b)$};

\draw[arrow] (x) -- (affine) node[midway, above] {$g(x)=ax+b$};
\draw[arrow] (affine) -- (f) node[midway, above] {$f$};
\draw[arrow] (f) -- (y);

\node[below=0.8cm of affine, text=poporange, font=\small\bfseries] {$g'(x)=a$};
\node[below=0.8cm of f, text=poporange, font=\small\bfseries] {$f'(ax+b)$};
\node[below=1.6cm of affine, xshift=1.5cm, text=poppurple, font=\small\bfseries] {$(f \circ g)' = a \cdot f'(ax+b)$};
\end{tikzpicture}
\legende{Schéma de décomposition : la dérivée de la composée est le produit de la dérivée de l'affine ($a$) par la dérivée de $f$ évaluée en $ax+b$.}
\end{popfigure}

\begin{methode}[Dériver une composée $f(ax+b)$]
\begin{enumerate}
    \item Identifier la fonction extérieure $f$ et l'affine intérieure $ax+b$.
    \item Calculer $f'(X)$ (en traitant $X$ comme variable).
    \item Remplacer $X$ par $ax+b$ dans $f'$.
    \item Multiplier par le coefficient $a$ de $x$ dans l'affine.
\end{enumerate}
\end{methode}

\begin{exemplebox}[Exemples de composées par une affine]
\begin{enumerate}
    \item $f(x) = \sqrt{3x-2}$ sur $\left[\frac{2}{3}; +\infty\right[$.
    $f(x) = g(3x-2)$ avec $g(X)=\sqrt{X}$, $g'(X)=\dfrac{1}{2\sqrt{X}}$.
    $f'(x) = 3 \cdot \dfrac{1}{2\sqrt{3x-2}} = \dfrac{3}{2\sqrt{3x-2}}$.
    
    \item $f(x) = (2x+1)^5$ sur $\mathbb{R}$.
    $f(x) = g(2x+1)$ avec $g(X)=X^5$, $g'(X)=5X^4$.
    $f'(x) = 2 \cdot 5(2x+1)^4 = 10(2x+1)^4$.
    
    \item $f(x) = \dfrac{1}{4x-1}$ sur $\left]-\infty; \frac{1}{4}\right[ \cup \left]\frac{1}{4}; +\infty\right[$.
    $f(x) = g(4x-1)$ avec $g(X)=\dfrac{1}{X}$, $g'(X)=-\dfrac{1}{X^2}$.
    $f'(x) = 4 \cdot \dfrac{-1}{(4x-1)^2} = -\dfrac{4}{(4x-1)^2}$.
\end{enumerate}
\end{exemplebox}

\courssub{Synthèse et entraînement : dériver toute expression algébrique}

Grâce aux règles précédentes (somme, produit par constante, produit, quotient, composée affine) et au tableau des dérivées usuelles, on peut dériver toute fonction obtenue par combinaisons algébriques de $x^n$, $\sqrt{x}$, $\frac{1}{x}$, etc.

\begin{methode}[Stratégie de dérivation complète]
\begin{enumerate}
    \item \textbf{Simplifier} l'expression si possible (développer, factoriser, simplifier une fraction).
    \item \textbf{Identifier la structure principale} (somme ? produit ? quotient ? composée affine ?).
    \item \textbf{Appliquer la règle correspondante} en nommant les sous-fonctions ($u$, $v$, $f$, $g$).
    \item \textbf{Calculer les dérivées élémentaires} (tableau de référence).
    \item \textbf{Assembler et simplifier} le résultat final (mise au même dénominateur, factorisation).
\end{enumerate}
\end{methode}

\begin{exoresolu}[Entraînement mixte]
Dériver les fonctions suivantes sur leur domaine de définition.
\begin{enumerate}
    \item $f(x) = x^2\sqrt{x} - \frac{3}{x}$.
    \item $g(x) = \frac{x^2+1}{x-1}$.
    \item $h(x) = (x^2+3x)(2x-1)$.
\end{enumerate}

\tcblower
\textbf{Solution 1 :} $f(x) = x^{5/2} - 3x^{-1}$ sur $]0;+\infty[$.
$f'(x) = \frac{5}{2}x^{3/2} - 3(-1)x^{-2} = \frac{5}{2}x\sqrt{x} + \frac{3}{x^2}$.

\textbf{Solution 2 :} $g(x) = \frac{u}{v}$ avec $u=x^2+1$, $v=x-1$ sur $]-\infty;1[\cup]1;+\infty[$.
$u'=2x$, $v'=1$.
$g'(x) = \frac{2x(x-1) - (x^2+1)\cdot 1}{(x-1)^2} = \frac{2x^2-2x-x^2-1}{(x-1)^2} = \frac{x^2-2x-1}{(x-1)^2}$.

\textbf{Solution 3 :} $h(x) = u \cdot v$ avec $u=x^2+3x$, $v=2x-1$ sur $\mathbb{R}$.
$u'=2x+3$, $v'=2$.
$h'(x) = (2x+3)(2x-1) + (x^2+3x)\cdot 2 = (4x^2+4x-3) + (2x^2+6x) = 6x^2+10x-3$.
\end{exoresolu}

\courssub{Complément : ouverture vers la Terminale (formule de la chaîne)}

\begin{savaistu}[La formule générale de la chaîne]
La règle $(f(ax+b))' = a f'(ax+b)$ est un cas particulier d'un théorème fondamental : la \cle{dérivée d'une fonction composée}.
Si $g$ est dérivable en $a$ et $u$ dérivable en $g(a)$, alors $u \circ g$ est dérivable en $a$ et :
\[
(u \circ g)'(a) = g'(a) \cdot u'(g(a)).
\]
Cette formule permettra de dériver des fonctions plus complexes en Terminale technique. Retenez déjà l'idée : \emph{on dérive de l'extérieur vers l'intérieur en multipliant par la dérivée de l'intérieur}.
\end{savaistu}

\begin{exercice}[Entraînement à la rédaction]
Dériver les fonctions suivantes sur leur ensemble de définition. Présenter les étapes : identification de $u,v$ ou $f,g$, calculs intermédiaires, résultat final simplifié.
\begin{enumerate}
    \item $f(x) = (x+2)(x^2-3x+1)$
    \item $g(x) = \frac{3x-2}{2x+1}$
    \item $h(x) = \sqrt{5-2x}$
    \item $k(x) = \frac{x}{\sqrt{x+1}}$
    \item $p(x) = (x^2+1)^3$ \emph{(indication : écrire comme composée de $X^3$)}
\end{enumerate}
\end{exercice}

\begin{corrige}[Exercice n°1]
\begin{enumerate}
    \item $f(x) = u \cdot v$, $u=x+2$, $v=x^2-3x+1$. $u'=1$, $v'=2x-3$.
    $f'(x) = 1\cdot(x^2-3x+1) + (x+2)(2x-3) = x^2-3x+1 + 2x^2+x-6 = 3x^2-2x-5$.
    
    \item $g(x) = \frac{u}{v}$, $u=3x-2$, $v=2x+1$. $u'=3$, $v'=2$.
    $g'(x) = \frac{3(2x+1) - (3x-2)\cdot 2}{(2x+1)^2} = \frac{6x+3-6x+4}{(2x+1)^2} = \frac{7}{(2x+1)^2}$.
    
    \item $h(x) = f(5-2x)$ avec $f(X)=\sqrt{X}$, $f'(X)=\dfrac{1}{2\sqrt{X}}$.
    $h'(x) = -2 \cdot \dfrac{1}{2\sqrt{5-2x}} = -\dfrac{1}{\sqrt{5-2x}}$ (domaine $\left]-\infty; \frac{5}{2}\right[$).
    
    \item $k(x) = \frac{u}{v}$, $u=x$, $v=\sqrt{x+1}$. $u'=1$, $v'=\dfrac{1}{2\sqrt{x+1}}$.
    $k'(x) = \dfrac{\sqrt{x+1} - x\cdot\dfrac{1}{2\sqrt{x+1}}}{x+1} = \dfrac{\dfrac{2(x+1)-x}{2\sqrt{x+1}}}{x+1} = \dfrac{x+2}{2(x+1)\sqrt{x+1}}$ (domaine $]-1;+\infty[$).
    
    \item $p(x) = f(x^2+1)$ avec $f(X)=X^3$, $f'(X)=3X^2$.
    $p'(x) = 2x \cdot 3(x^2+1)^2 = 6x(x^2+1)^2$.
\end{enumerate}
\end{corrige}

\coursec{Tangente et approximation affine}

Après avoir maîtrisé les règles de dérivation (somme, produit, quotient, composée) qui nous permettent de calculer $f'(a)$ pour une grande variété de fonctions, nous abordons maintenant deux applications géométriques et numériques majeures de la dérivée : la \textbf{tangente} à une courbe et l'\textbf{approximation affine}. Ces deux notions sont les deux faces d'une même médaille : \emph{localement, une fonction dérivable se comporte comme une fonction affine}.

\courssub{La tangente à une courbe : intuition et définition rigoureuse}

\textbf{Intuition : la limite des sécantes.}\par
Soit $f$ une fonction définie sur un intervalle $I$ et $a \in I$. Le point $A(a\,;\,f(a))$ appartient à la courbe représentative $\mathcal{C}_f$. Prenons un autre point $M(x\,;\,f(x))$ sur $\mathcal{C}_f$, distinct de $A$. La droite $(AM)$ est une \textbf{sécante} à la courbe. Son coefficient directeur est le taux d'accroissement :
\[
\frac{f(x)-f(a)}{x-a}.
\]
Lorsque $M$ se rapproche de $A$ (c'est-à-dire $x \to a$), la sécante $(AM)$ pivote autour de $A$ et tend vers une position limite : c'est la \textbf{tangente} à $\mathcal{C}_f$ en $A$, \emph{si cette limite existe}. Or, cette limite est exactement la définition du nombre dérivé $f'(a)$.

\begin{experience}[Démonstration guidée : la tangente comme limite des sécantes]
On considère $f(x)=x^2$ et le point $A(1\,;\,1)$.
\begin{enumerate}
  \item Écrire le coefficient directeur de la sécante $(AM)$ pour $M(x\,;\,x^2)$, $x \neq 1$.
  \item Calculer la limite de ce coefficient lorsque $x \to 1$.
  \item En déduire l'équation de la tangente en $A$.
  \item Vérifier graphiquement (sur papier ou logiciel) que la droite obtenue « colle » à la courbe au voisinage de $A$.
\end{enumerate}
\tcblower
\textbf{Solution.}
\begin{enumerate}
  \item $m_{AM} = \dfrac{x^2-1}{x-1} = \dfrac{(x-1)(x+1)}{x-1} = x+1$ (pour $x \neq 1$).
  \item $\lim\limits_{x \to 1} m_{AM} = \lim\limits_{x \to 1} (x+1) = 2 = f'(1)$.
  \item La tangente a pour coefficient directeur $2$ et passe par $A(1\,;\,1)$ : $y-1 = 2(x-1) \iff y = 2x-1$.
  \item Au voisinage de $x=1$, la courbe $y=x^2$ et la droite $y=2x-1$ sont quasi confondues.
\end{enumerate}
\end{experience}

\begin{propriete}[Équation de la tangente]
Soit $f$ une fonction dérivable en $a$. La courbe $\mathcal{C}_f$ admet en $A(a\,;\,f(a))$ une tangente non verticale. Son équation cartésienne est :
\[
\boxed{\,y = f'(a)(x-a) + f(a)\,}
\]
\textbf{Preuve (exigible au Probatoire).} La tangente est la droite passant par $A$ de coefficient directeur $f'(a)$. Une droite d'équation $y = mx + p$ passant par $A$ vérifie $f(a) = m a + p$, d'où $p = f(a) - m a$. Avec $m = f'(a)$, on obtient $y = f'(a)x + f(a) - f'(a)a = f'(a)(x-a) + f(a)$. \qed
\end{propriete}

\begin{methode}[Déterminer l'équation de la tangente en un point]
Pour une fonction $f$ dérivable en $a$ :
\begin{enumerate}
  \item Calculer $f(a)$ (ordonnée du point de contact).
  \item Calculer $f'(a)$ (coefficient directeur de la tangente).
  \item Substituer dans la formule $y = f'(a)(x-a) + f(a)$ et simplifier sous la forme $y = mx + p$.
\end{enumerate}
\end{methode}

\begin{exemplebox}[Tangente à la parabole $f(x)=x^2$ en $x=3$]
$f(x)=x^2$, $a=3$.
\begin{enumerate}
  \item $f(3) = 9$. Le point de contact est $A(3\,;\,9)$.
  \item $f'(x)=2x$, donc $f'(3)=6$.
  \item Équation : $y = 6(x-3) + 9 = 6x - 18 + 9 = 6x - 9$.
\end{enumerate}
La tangente en $A$ a pour équation $\boxed{y = 6x - 9}$.
\end{exemplebox}

\begin{attention}[Piège classique : oublier le terme $-f'(a)\cdot a$]
\emph{Erreur fréquente :} écrire directement $y = f'(a)x + f(a)$.
\emph{Exemple :} pour $f(x)=x^2$ en $3$, écrire $y = 6x + 9$ (faux).
\emph{Correct :} $y = f'(a)(x-a) + f(a) = f'(a)x + \bigl(f(a) - a f'(a)\bigr)$.
Ici $p = 9 - 3\times 6 = -9$, d'où $y = 6x - 9$.
\textbf{Toujours développer $(x-a)$ ou utiliser la forme point-pente $y - f(a) = f'(a)(x-a)$.}
\end{attention}

\begin{popfigure}
\begin{tikzpicture}[scale=0.9, >=Stealth]
\begin{axis}[
    width=\linewidth, height=7cm,
    axis lines=middle,
    xlabel=$x$, ylabel=$y$,
    xmin=-0.5, xmax=5.5,
    ymin=-2, ymax=18,
    xtick={1,2,3,4,5}, ytick={0,5,10,15},
    clip=false,
    domain=-0.5:5.5, samples=200,
]
\addplot[popcurve, thick, domain=0:5] {x^2} node[right, pos=0.9] {$y=x^2$};
\addplot[poporange, thick, domain=-0.5:5.5] {6*x-9} node[right, pos=0.95] {$y=6x-9$};
\addplot[only marks, mark=*, mark size=2.5, popdark] coordinates {(3,9)};
\node[above left] at (axis cs:3,9) {$A(3,\,9)$};
% Zoom box
\draw[popblue, thick] (axis cs:2.5,4) rectangle (axis cs:3.5,14);
\end{axis}
% Zoom inset
\begin{axis}[
    width=5.5cm, height=4.5cm,
    at={(current axis.right of origin)},
    anchor=left of origin,
    xshift=1.2cm,
    axis lines=middle,
    xmin=2.7, xmax=3.3,
    ymin=7, ymax=11,
    xtick={2.8,3,3.2}, ytick={8,9,10},
    tick label style={font=\tiny},
    clip=true,
    domain=2.7:3.3, samples=200,
]
\addplot[popcurve, thick] {x^2};
\addplot[poporange, thick] {6*x-9};
\addplot[only marks, mark=*, mark size=2, popdark] coordinates {(3,9)};
\end{axis}
\node[popblue, font=\small] at (current axis.right of origin) [xshift=3.8cm, yshift=-1.5cm] {Zoom : courbe et tangente quasi confondues};
\end{tikzpicture}
\legende{Courbe $y=x^2$, tangente en $A(3,9)$ et zoom local montrant la confusion.}
\end{popfigure}

\courssub{L'approximation affine : le principe du « localement affine »}

\begin{definition}[Approximation affine d'une fonction en un point]
Soit $f$ dérivable en $a$. Pour $h$ proche de $0$, on a l'\textbf{approximation affine} (ou approximation linéaire) :
\[
\boxed{\,f(a+h) \approx f(a) + h\,f'(a)\,}
\]
La fonction affine $h \mapsto f(a) + h f'(a)$ est l'\textbf{application tangentielle} de $f$ en $a$. L'erreur commise est $E(h) = f(a+h) - \bigl(f(a) + h f'(a)\bigr)$ ; on a $\lim\limits_{h\to 0} \frac{E(h)}{h} = 0$ (l'erreur est négligeable devant $h$).
\end{definition}

\textbf{Interprétation géométrique.}\par
$f(a+h)$ est l'ordonnée du point de la courbe d'abscisse $a+h$. $f(a) + h f'(a)$ est l'ordonnée du point de la \emph{tangente} de même abscisse. L'approximation affine consiste à \textbf{remplacer la courbe par sa tangente} au voisinage de $A$.

\begin{popfigure}
\begin{tikzpicture}[scale=1.0, >=Stealth]
\begin{axis}[
    width=\linewidth, height=7cm,
    axis lines=middle,
    xlabel=$x$, ylabel=$y$,
    xmin=0, xmax=5.5,
    ymin=0, ymax=4,
    xtick={1,2,3,4,5}, ytick={1,2,3},
    clip=false,
    domain=0:5.5, samples=200,
]
% Courbe f(x) = sqrt(x)
\addplot[popcurve, thick, domain=0:5.5] {sqrt(max(0,x))} node[right, pos=0.9] {$\mathcal{C}_f : y=\sqrt{x}$};
% Tangente en a=4 : y = 2 + 0.25(x-4) = 0.25x + 1
\addplot[poporange, thick, domain=0:5.5] {0.25*x+1} node[above right, pos=0.95] {Tangente en $A$};
% Point A
\addplot[only marks, mark=*, mark size=2.5, popdark] coordinates {(4,2)};
\node[above left] at (axis cs:4,2) {$A(4,\,2)$};
% Point M sur la courbe (a+h) avec h=0.8
\addplot[only marks, mark=*, mark size=2.5, popdark] coordinates {(4.8,2.19089)};
\node[right] at (axis cs:4.8,2.19089) {$M(4.8,\,f(4.8))$};
% Point T sur la tangente
\addplot[only marks, mark=*, mark size=2.5, poporange] coordinates {(4.8,2.2)};
\node[above right] at (axis cs:4.8,2.2) {$T(4.8,\,f(4)+hf'(4))$};
% Segments verticaux pour erreur
\draw[dashed, popmuted] (axis cs:4.8,2.19089) -- (axis cs:4.8,2.2) node[midway, right] {$E(h)$};
\draw[dashed, popmuted] (axis cs:4,0) -- (axis cs:4,2) node[below] {$a=4$};
\draw[dashed, popmuted] (axis cs:4.8,0) -- (axis cs:4.8,2.19089) node[below] {$a+h=4.8$};
% h brace
\draw[<->, popblue] (axis cs:4,-0.2) -- (axis cs:4.8,-0.2) node[midway, below] {$h$};
% f(a) brace
\draw[<->, popblue] (axis cs:-0.2,0) -- (axis cs:-0.2,2) node[midway, left] {$f(a)$};
% Slope triangle on tangent
\draw[poporange, dashed] (axis cs:4,2) -- (axis cs:5,2) -- (axis cs:5,2.25) -- cycle;
\node[poporange, font=\small] at (axis cs:4.5, 2.1) {pente $f'(a)=0,25$};
\end{axis}
\end{tikzpicture}
\legende{Schéma de l'approximation affine pour $f(x)=\sqrt{x}$ en $a=4$ : $f(a+h) \approx f(a) + h f'(a)$. L'erreur $E(h)$ est l'écart vertical entre la courbe et sa tangente.}
\end{popfigure}

\courssub{Applications numériques : estimation sans calculatrice}

L'approximation affine permet d'estimer des valeurs numériques difficiles à calculer mentalement, en choisissant $a$ « commode » (proche de la valeur cherchée, avec $f(a)$ et $f'(a)$ simples).

\begin{exoresolu}[Estimer $\sqrt{4,02}$, $\dfrac{1}{0,98}$ et $(1,001)^3$]
\textbf{1. $\sqrt{4,02}$}\\
$f(x)=\sqrt{x}$, $a=4$, $h=0,02$. $f(4)=2$, $f'(x)=\dfrac{1}{2\sqrt{x}}$, $f'(4)=\dfrac{1}{4}=0,25$.
\[
\sqrt{4,02} \approx 2 + 0,02 \times 0,25 = 2 + 0,005 = \boxed{2,005}.
\]
(Valeur exacte : $2,00499\ldots$ ; erreur $\approx 5\times 10^{-7}$.)

\textbf{2. $\dfrac{1}{0,98}$}\\
$f(x)=\dfrac{1}{x}$, $a=1$, $h=-0,02$. $f(1)=1$, $f'(x)=-\dfrac{1}{x^2}$, $f'(1)=-1$.
\[
\frac{1}{0,98} \approx 1 + (-0,02)\times(-1) = 1 + 0,02 = \boxed{1,02}.
\]
(Valeur exacte : $1,020408\ldots$ ; erreur $\approx 4\times 10^{-4}$.)

\textbf{3. $(1,001)^3$}\\
$f(x)=x^3$, $a=1$, $h=0,001$. $f(1)=1$, $f'(x)=3x^2$, $f'(1)=3$.
\[
(1,001)^3 \approx 1 + 0,001 \times 3 = \boxed{1,003}.
\]
(Valeur exacte : $1,003003001$ ; erreur $\approx 3\times 10^{-9}$.)
\end{exoresolu}

\begin{attention}[Choix de $a$ et ordre de grandeur de l'erreur]
\begin{itemize}
  \item $a$ doit être \textbf{proche} de la valeur cherchée et \textbf{simple} ($f(a)$ et $f'(a)$ calculables mentalement).
  \item L'erreur est d'ordre $h^2$ (proportionnelle à $h^2$ pour les fonctions $C^2$). Plus $|h|$ est petit, plus l'approximation est précise.
  \item \textbf{Qualitativement} : si la courbe est \emph{au-dessus} de sa tangente (convexe, $f''>0$), l'approximation affine \textbf{sous-estime} la valeur vraie. Si elle est \emph{en-dessous} (concave, $f''<0$), elle \textbf{sur-estime}.
\end{itemize}
\end{attention}

\courssub{Estimation qualitative de l'erreur (lecture graphique)}

Sans calculer $f''$, on peut deviner le signe de l'erreur en regardant la courbe :
\begin{itemize}
  \item Courbe \textbf{convexe} (creusée vers le haut, ex. $x^2$, $1/x$ pour $x>0$) : la courbe est \emph{au-dessus} de la tangente $\Rightarrow$ approximation affine \textbf{sous-estime} ($f(a+h) > f(a)+hf'(a)$).
  \item Courbe \textbf{concave} (bombée vers le haut, ex. $\sqrt{x}$, $\ln x$) : la courbe est \emph{en-dessous} de la tangente $\Rightarrow$ approximation affine \textbf{sur-estime}.
\end{itemize}

\begin{exemplebox}[Signe de l'erreur pour les exemples précédents]
\begin{itemize}
  \item $\sqrt{x}$ est \textbf{concave} sur $\mathbb{R}_+^*$ $\Rightarrow$ tangente \emph{au-dessus} de la courbe $\Rightarrow$ $\sqrt{4,02} < 2,005$ : l'approximation $2,005$ est une \textbf{sur-estimation}. Vérification : $2,00499 < 2,005$. ✓
  \item $1/x$ est \textbf{convexe} sur $\mathbb{R}_+^*$ $\Rightarrow$ tangente \emph{en-dessous} $\Rightarrow$ approximation $1,02$ est une \textbf{sous-estimation}. Vérification : $1,0204 > 1,02$. ✓
  \item $x^3$ est \textbf{convexe} sur $\mathbb{R}_+^*$ ($f''(x)=6x>0$) $\Rightarrow$ tangente \emph{en-dessous} $\Rightarrow$ approximation $1,003$ est une \textbf{sous-estimation}. Vérification : $1,003003 > 1,003$. ✓
\end{itemize}
\end{exemplebox}

\courssub{Applications aux situations concrètes (métiers techniques)}

\begin{exoresolu}[Dilatation thermique d'une poutre métallique]
Une poutre en acier a une longueur $L_0 = 10,000$ m à $20^\circ\text{C}$. Son coefficient de dilatation linéique est $\alpha = 12 \times 10^{-6}\,\text{K}^{-1}$. La longueur à la température $T$ est $L(T) = L_0\bigl[1 + \alpha(T-20)\bigr]$.
On veut estimer la longueur à $T = 25^\circ\text{C}$ sans calculatrice.
\begin{enumerate}
  \item Identifier $f$, $a$, $h$.
  \item Calculer $f(a)$ et $f'(a)$.
  \item Donner l'estimation par approximation affine.
  \item Comparer à la valeur exacte.
\end{enumerate}
\tcblower
\textbf{Solution.}
\begin{enumerate}
  \item $f(T)=L(T)$, $a=20$, $h=5$.
  \item $f(20)=10,000$ m. $f'(T)=L_0\alpha = 10 \times 12\times 10^{-6} = 1,2\times 10^{-4}$ m/$^\circ$C. Donc $f'(20)=0,00012$ m/$^\circ$C.
  \item $L(25) \approx 10,000 + 5 \times 0,00012 = 10,000 + 0,0006 = \boxed{10,0006\,\text{m}}$ (soit $10,0006$ m ou $10$ m $0,6$ mm).
  \item Exact : $10[1+12\times10^{-6}\times5] = 10[1+60\times10^{-6}] = 10,0006$ m. \textbf{Égalité exacte ici car $f$ est affine !} L'approximation affine est \emph{exacte} pour une fonction affine.
\end{enumerate}
\end{exoresolu}

\begin{exoresolu}[Chute de tension dans un câble (électricité)]
La résistance $R$ d'un fil de cuivre varie avec la température selon $R(T) = R_0[1 + \alpha(T-T_0)]$ avec $\alpha = 3,9\times 10^{-3}\,\text{K}^{-1}$. Pour un câble de $R_0=2,50\,\Omega$ à $T_0=20^\circ\text{C}$, estimer $R$ à $22^\circ\text{C}$.
\tcblower
$f(T)=R(T)$, $a=20$, $h=2$. $f(20)=2,50\,\Omega$. $f'(T)=R_0\alpha = 2,50\times 3,9\times 10^{-3} = 9,75\times 10^{-3}\,\Omega/^\circ\text{C}$.
$R(22) \approx 2,50 + 2\times 0,00975 = 2,50 + 0,0195 = \boxed{2,5195\,\Omega}$.
(Valeur exacte : $2,50[1+3,9\times10^{-3}\times2] = 2,5195\,\Omega$ — encore exact car modèle affine.)
\end{exoresolu}

\begin{exoresolu}[Coût marginal et approximation de coût (gestion)]
Le coût total de production (en milliers de FCFA) est modélisé par $C(x) = 0,02x^3 - 1,5x^2 + 60x + 500$ pour $x$ unités produites (en centaines). Estimer le coût supplémentaire pour passer de $1000$ à $1005$ unités.
\tcblower
$x$ en centaines : $a=10$ (1000 unités), $h=0,05$ (1005 unités = 10,05 centaines $\Rightarrow h=0,05$).
$C(10) = 0,02(1000) - 1,5(100) + 600 + 500 = 20 - 150 + 1100 = 970$ kFCFA.
$C'(x) = 0,06x^2 - 3x + 60$. $C'(10) = 0,06(100) - 30 + 60 = 6 + 30 = 36$ kFCFA par centaine.
Coût supplémentaire $\approx h \cdot C'(10) = 0,05 \times 36 = 1,8$ kFCFA = \boxed{1\,800\,\text{FCFA}}.
(Valeur exacte : $C(10,05)-C(10) \approx 1\,801,35$ FCFA. Erreur $\approx 1,35$ FCFA, négligeable.)
\end{exoresolu}

\begin{savaistu}[Le tracé des courbes par les machines]
Les calculatrices graphiques, les logiciels de CAO (Conception Assistée par Ordinateur) et les moteurs de jeux vidéo ne « calculent » pas la courbe en chaque pixel. Ils utilisent l'\textbf{approximation affine} (ou sa version améliorée : interpolation polynomiale) : ils calculent quelques points exacts, puis relient ces points par de \textbf{petits segments de droite} (tangentes ou cordes). Si le pas est assez petit, l'œil ne voit qu'une courbe lisse. C'est le principe de la \emph{discrétisation} : remplacer le continu par du linéaire par morceaux. En impression 3D, c'est pareil : la courbe est découpée en facettes planes (triangles) — l'approximation affine en 3D !
\end{savaistu}

\begin{exercice}[Entraînement — Tangente et approximation affine]
\begin{enumerate}
  \item Déterminer l'équation de la tangente à la courbe de $f(x)=\dfrac{1}{x}$ au point d'abscisse $2$.
  \item En déduire une valeur approchée de $\dfrac{1}{2,03}$.
  \item La courbe de $f$ est-elle au-dessus ou en-dessous de sa tangente en $2$ ? Justifier sans calculer $f''$. Quel est le signe de l'erreur d'approximation ?
  \item Une entreprise a un chiffre d'affaires modélisé par $CA(t) = 500\sqrt{t+4}$ (en millions de FCFA), $t$ en années depuis 2020. Estimer le CA au 1er juillet 2024 ($t=4,5$) par approximation affine au voisinage de $t=4$.
\end{enumerate}
\end{exercice}

\begin{corrige}[Exercice 1]
\begin{enumerate}
  \item $f(x)=1/x$, $a=2$. $f(2)=0,5$. $f'(x)=-1/x^2$, $f'(2)=-0,25$.
      Tangente : $y = -0,25(x-2) + 0,5 = -0,25x + 0,5 + 0,5 = \boxed{y = -0,25x + 1}$.
  \item $h = 0,03$. $\dfrac{1}{2,03} \approx 0,5 + 0,03\times(-0,25) = 0,5 - 0,0075 = \boxed{0,4925}$.
  \item $f(x)=1/x$ est convexe sur $\mathbb{R}_+^*$ (sa courbe est une branche d'hyperbole « creusée vers le haut »). La courbe est \textbf{au-dessus} de sa tangente. Donc l'approximation affine (valeur sur la tangente) \textbf{sous-estime} la valeur réelle : $\dfrac{1}{2,03} > 0,4925$. (Vérification : $1/2,03 \approx 0,49261$.)
  \item $f(t)=500\sqrt{t+4}$, $a=4$, $h=0,5$. $f(4)=500\sqrt{8}=500\times 2\sqrt{2}\approx 1414,2$. $f'(t)=\dfrac{500}{2\sqrt{t+4}}=\dfrac{250}{\sqrt{t+4}}$, $f'(4)=\dfrac{250}{\sqrt{8}}=\dfrac{250}{2\sqrt{2}}\approx 88,39$.
      $CA(4,5) \approx 1414,2 + 0,5\times 88,39 = 1414,2 + 44,195 = \boxed{1458,4\,\text{millions FCFA}}$.
\end{enumerate}
\end{corrige}

\coursec{Dérivée et sens de variation}

Dans la section précédente, nous avons vu que la dérivée $f'(a)$ donne la \cle{pente de la tangente} à la courbe au point d'abscisse $a$, et que la tangente fournit une excellente approximation de la courbe au voisinage de ce point (approximation affine). Regardons maintenant les choses plus globalement : si, sur tout un intervalle, les tangentes « montent » (pente positive), la courbe elle-même doit monter ; si elles « descendent » (pente négative), la courbe descend. Autrement dit, le \cle{signe de la fonction dérivée} $f'$ sur un intervalle renseigne sur le \cle{sens de variation} de la fonction $f$ sur cet intervalle. C'est l'un des résultats les plus puissants et les plus utilisés de toute l'année : il permet de dresser le tableau de variations d'une fonction, de trouver ses maximums et minimums, et donc de résoudre des problèmes concrets d'optimisation (coût minimal d'un atelier, bénéfice maximal d'une production, rendement maximal d'une machine).

\courssub{Le théorème fondamental : signe de $f'$ et sens de variation de $f$}

\textbf{Intuition.} Imagine un technicien qui relève la température d'un four toutes les minutes. La fonction $f$ donne la température, et $f'(x)$ donne la \emph{vitesse de variation} de cette température à l'instant $x$. Si cette vitesse est positive, la température augmente ; si elle est négative, elle diminue ; si elle est nulle en permanence, la température reste constante. La dérivée joue exactement ce rôle de « vitesse de variation instantanée ».

\begin{propriete}[Théorème fondamental (admis)]
Soit $f$ une fonction dérivable sur un intervalle $I$.
\begin{itemize}
\item Si $f'(x) \geq 0$ pour tout $x$ de $I$, alors $f$ est \cle{croissante} sur $I$.
\item Si $f'(x) \leq 0$ pour tout $x$ de $I$, alors $f$ est \cle{décroissante} sur $I$.
\item Si $f'(x) = 0$ pour tout $x$ de $I$, alors $f$ est \cle{constante} sur $I$.
\end{itemize}
\textbf{Version stricte :} si $f'(x) > 0$ pour tout $x$ de $I$, sauf éventuellement en des points isolés où $f'$ s'annule, alors $f$ est \cle{strictement croissante} sur $I$ (même énoncé avec $<0$ pour strictement décroissante).
\par\smallskip
\emph{Ce théorème est admis en classe de Première technique : sa démonstration n'est pas exigible au Probatoire.}
\end{propriete}

\begin{attention}[Bien lire les hypothèses]
Deux conditions sont indispensables avant d'appliquer le théorème :
\begin{enumerate}
\item $f$ doit être \textbf{dérivable sur tout l'intervalle} $I$ considéré ;
\item le signe de $f'$ doit être étudié \textbf{sur cet intervalle}, pas seulement en un point.
\end{enumerate}
On ne peut pas conclure « $f'(2) > 0$ donc $f$ est croissante » : la positivité en un seul point ne dit rien sur les variations globales.
\end{attention}

\begin{exemplebox}[Lecture directe]
Soit $f$ définie sur $\mathbb{R}$ par $f(x) = x^2$. Alors $f'(x) = 2x$.
\begin{itemize}
\item Sur $]-\infty ; 0]$, $f'(x) \leq 0$ : $f$ est décroissante.
\item Sur $[0 ; +\infty[$, $f'(x) \geq 0$ : $f$ est croissante.
\end{itemize}
On retrouve bien les variations connues de la fonction carré, avec un minimum en $0$.
\end{exemplebox}

\courssub{Méthode : dresser un tableau de variations complet}

\begin{methode}[Étudier les variations d'une fonction en 4 étapes]
Pour étudier les variations d'une fonction $f$ :
\begin{enumerate}
\item \textbf{Ensemble} : déterminer l'ensemble de définition de $f$ et vérifier que $f$ y est dérivable (sur chaque intervalle).
\item \textbf{Dérivée} : calculer $f'(x)$ à l'aide des opérations sur les dérivées (somme, produit, quotient, composée par une fonction affine).
\item \textbf{Signe} : résoudre $f'(x) = 0$, puis étudier le \textbf{signe de $f'(x)$} (tableau de signes, factorisation éventuelle).
\item \textbf{Tableau} : dresser le tableau de variations de $f$ : flèche montante $\nearrow$ là où $f' > 0$, flèche descendante $\searrow$ là où $f' < 0$, et calculer les valeurs de $f$ aux points remarquables (bornes et zéros de $f'$).
\end{enumerate}
\end{methode}

\begin{exoresolu}[Variations de $f(x) = x^3 - 3x + 1$ sur $\mathbb{R}$]
Énoncé : étudier les variations de la fonction $f$ définie sur $\mathbb{R}$ par $f(x) = x^3 - 3x + 1$, et dresser son tableau de variations.
\tcblower
\textbf{Solution détaillée.}
\par\textbf{Étape 1.} $f$ est une fonction polynôme, donc définie et dérivable sur $\mathbb{R}$.
\par\textbf{Étape 2.} Pour tout réel $x$ :
\[ f'(x) = 3x^2 - 3 = 3(x^2 - 1) = 3(x-1)(x+1). \]
\par\textbf{Étape 3.} $f'(x) = 0 \iff x = -1$ ou $x = 1$. Le trinôme $3(x-1)(x+1)$ est du signe de $(x-1)(x+1)$ :
\begin{itemize}
\item $f'(x) > 0$ sur $]-\infty ; -1[$ et sur $]1 ; +\infty[$ ;
\item $f'(x) < 0$ sur $]-1 ; 1[$.
\end{itemize}
\par\textbf{Étape 4.} Valeurs remarquables :
\[ f(-1) = -1 + 3 + 1 = 3 \qquad ; \qquad f(1) = 1 - 3 + 1 = -1. \]
Tableau de variations :
\begin{center}
\begin{tabular}{|c|ccccccc|}\hline
$x$ & $-\infty$ & & $-1$ & & $1$ & & $+\infty$ \\ \hline
$f'(x)$ & & $+$ & $0$ & $-$ & $0$ & $+$ & \\ \hline
$f$ & & $\nearrow$ & $3$ & $\searrow$ & $-1$ & $\nearrow$ & \\ \hline
\end{tabular}
\end{center}
\par\textbf{Conclusion.} $f$ est croissante sur $]-\infty ; -1]$, décroissante sur $[-1 ; 1]$, croissante sur $[1 ; +\infty[$. Elle admet un \textbf{maximum local} égal à $3$ en $x = -1$ et un \textbf{minimum local} égal à $-1$ en $x = 1$.
\end{exoresolu}

\begin{popfigure}
\begin{tikzpicture}
\begin{axis}[
  width=0.85\linewidth, height=6.5cm,
  axis lines=middle, xlabel=$x$, ylabel=$y$,
  xmin=-2.6, xmax=2.6, ymin=-3.5, ymax=4.5,
  xtick={-2,-1,1,2}, ytick={-1,3},
  grid=both, grid style={popline!40},
]
\addplot[popcurve,domain=-2.25:2.25,samples=200]{x^3 - 3*x + 1};
\addplot[poporange,thick,domain=-1.7:-0.3]{3};
\addplot[poppurple,thick,domain=0.3:1.7]{-1};
\addplot[only marks,mark=*,popdark] coordinates {(-1,3) (1,-1)};
\node[popdark,anchor=south west] at (axis cs:-1,3) {\small $M(-1\,;\,3)$};
\node[popdark,anchor=north west] at (axis cs:1,-1) {\small $N(1\,;\,-1)$};
\end{axis}
\end{tikzpicture}
\legende{Courbe de $f(x) = x^3 - 3x + 1$. Les tangentes en $x = -1$ et $x = 1$ sont horizontales car $f'(-1) = f'(1) = 0$.}
\end{popfigure}

\courssub{Extremums locaux et dérivée}

\begin{definition}[Extremum local]
Soit $f$ une fonction définie sur un intervalle $I$ et $c$ un point de $I$.
\begin{itemize}
\item $f$ admet un \cle{maximum local} en $c$ s'il existe un intervalle ouvert contenant $c$ sur lequel $f(x) \leq f(c)$.
\item $f$ admet un \cle{minimum local} en $c$ s'il existe un intervalle ouvert contenant $c$ sur lequel $f(x) \geq f(c)$.
\end{itemize}
Un extremum local est un maximum local ou un minimum local.
\end{definition}

\textbf{Observation.} Sur la figure précédente, aux points $M$ et $N$ où $f$ atteint ses extremums locaux, la tangente est \emph{horizontale} : cela signifie $f'(-1) = 0$ et $f'(1) = 0$. Ce n'est pas un hasard.

\begin{propriete}[Condition nécessaire d'extremum local]
Soit $f$ une fonction dérivable sur un intervalle ouvert $I$ et $c \in I$. Si $f$ admet un extremum local en $c$, alors $f'(c) = 0$.
\end{propriete}

\begin{experience}[Démonstration guidée : pourquoi $f'(c) = 0$ en un extremum ?]
Supposons que $f$ admet un \textbf{maximum local} en $c$ (le cas du minimum est analogue). Alors, pour $h$ assez proche de $0$, on a $f(c+h) \leq f(c)$, c'est-à-dire $f(c+h) - f(c) \leq 0$.
\begin{enumerate}
\item \textbf{Taux à droite.} Pour $h > 0$ petit :
\[ \frac{f(c+h)-f(c)}{h} \leq 0 \quad \text{(numérateur} \leq 0,\ \text{dénominateur} > 0\text{)}. \]
En faisant tendre $h$ vers $0$ par valeurs positives, on obtient $f'(c) \leq 0$.
\item \textbf{Taux à gauche.} Pour $h < 0$ petit :
\[ \frac{f(c+h)-f(c)}{h} \geq 0 \quad \text{(numérateur} \leq 0,\ \text{dénominateur} < 0\text{)}. \]
En faisant tendre $h$ vers $0$ par valeurs négatives, on obtient $f'(c) \geq 0$.
\item \textbf{Conclusion.} $f'(c) \leq 0$ et $f'(c) \geq 0$, donc $f'(c) = 0$. \hfill $\blacksquare$
\end{enumerate}
Retiens l'idée : au sommet d'une « colline », la pente ne peut être ni positive (on monterait encore) ni négative (on serait déjà descendu) : elle est nulle.
\end{experience}

\begin{attention}[Erreur fréquente : $f'(c) = 0$ ne suffit pas !]
La réciproque du théorème précédent est \textbf{fausse} : on peut avoir $f'(c) = 0$ sans extremum en $c$.
\par\textbf{Contre-exemple.} Soit $f(x) = x^3$. Alors $f'(x) = 3x^2$, donc $f'(0) = 0$. Pourtant $f'(x) \geq 0$ partout : $f$ est strictement croissante sur $\mathbb{R}$ et n'a \textbf{aucun extremum} en $0$ (la tangente en $0$ est horizontale, mais la courbe la traverse : c'est un \emph{point d'inflexion à tangente horizontale}).
\par\textbf{Règle pratique.} Pour affirmer un extremum local en $c$, il faut vérifier que $f'$ \textbf{s'annule en $c$ EN CHANGEANT DE SIGNE} :
\begin{itemize}
\item $f'$ passe de $+$ à $-$ en $c$ $\Rightarrow$ maximum local en $c$ ;
\item $f'$ passe de $-$ à $+$ en $c$ $\Rightarrow$ minimum local en $c$ ;
\item $f'$ s'annule sans changer de signe $\Rightarrow$ pas d'extremum.
\end{itemize}
\end{attention}

\courssub{Correspondance graphique entre $f$ et $f'$}

Une excellente façon de vérifier un tableau de variations est de comparer la courbe de $f$ avec celle de $f'$ tracée juste en dessous, avec les mêmes abscisses : là où la courbe de $f'$ est \emph{au-dessus} de l'axe des abscisses, la courbe de $f$ \emph{monte} ; là où elle est en dessous, $f$ \emph{descend} ; là où $f'$ coupe l'axe, $f$ change de sens de variation.

\begin{popfigure}
\begin{tikzpicture}
\begin{axis}[
  width=0.85\linewidth, height=5.2cm,
  axis lines=middle, ylabel=$f(x)$,
  xmin=-2.6, xmax=2.6, ymin=-3.5, ymax=4.5,
  xtick={-1,1}, xticklabels={$-1$,$1$}, ytick=\empty,
  grid=both, grid style={popline!40},
]
\addplot[popcurve,domain=-2.25:2.25,samples=200]{x^3 - 3*x + 1};
\addplot[popmuted,dashed] coordinates {(-1,-3.5) (-1,4.5)};
\addplot[popmuted,dashed] coordinates {(1,-3.5) (1,4.5)};
\end{axis}
\end{tikzpicture}

\vspace{0.2cm}
\begin{tikzpicture}
\begin{axis}[
  width=0.85\linewidth, height=4.2cm,
  axis lines=middle, xlabel=$x$, ylabel=$f'(x)$,
  xmin=-2.6, xmax=2.6, ymin=-4.5, ymax=10.5,
  xtick={-1,1}, xticklabels={$-1$,$1$}, ytick=\empty,
  grid=both, grid style={popline!40},
]
\addplot[popcurve,poporange,domain=-2.05:2.05,samples=200]{3*x^2 - 3};
\addplot[popmuted,dashed] coordinates {(-1,-4.5) (-1,10.5)};
\addplot[popmuted,dashed] coordinates {(1,-4.5) (1,10.5)};
\addplot[only marks,mark=*,popdark] coordinates {(-1,0) (1,0)};
\end{axis}
\end{tikzpicture}
\legende{En haut : $f(x) = x^3 - 3x + 1$. En bas : $f'(x) = 3x^2 - 3$. Quand $f' > 0$ (courbe orange au-dessus de l'axe), $f$ croît ; quand $f' < 0$, $f$ décroît ; aux zéros de $f'$ ($x = -1$ et $x = 1$), $f$ atteint ses extremums locaux.}
\end{popfigure}

\courssub{Recherche de maximum ou de minimum sur un intervalle}

\begin{methode}[{Trouver le maximum et le minimum de $f$ sur $[a ; b]$}]
\begin{enumerate}
\item Dresser le tableau de variations de $f$ sur $[a ; b]$ (signe de $f'$).
\item Repérer les extremums locaux à l'intérieur de l'intervalle.
\item \textbf{Comparer} les valeurs de $f$ aux extremums locaux \textbf{avec les valeurs aux bornes} $f(a)$ et $f(b)$.
\item Le plus grand de ces nombres est le maximum de $f$ sur $[a ; b]$, le plus petit est le minimum.
\end{enumerate}
\end{methode}

\begin{attention}[Ne pas oublier les bornes]
Un extremum \emph{local} n'est pas forcément un extremum \emph{global} sur l'intervalle. Par exemple, pour $f(x) = x^3 - 3x + 1$ sur $[-2 ; 2]$ : le maximum local vaut $f(-1) = 3$, mais $f(2) = 8 - 6 + 1 = 3$ également, et $f(-2) = -8 + 6 + 1 = -1$. Ici le maximum sur $[-2;2]$ vaut $3$ (atteint en $-1$ et en $2$) et le minimum vaut $-1$ (atteint en $-2$ et en $1$). Toujours confronter extremums locaux et valeurs aux bornes.
\end{attention}

\begin{exoresolu}[Bénéfice maximal d'un atelier]
Énoncé : un atelier de Douala produit $x$ centaines de pièces par semaine ($x \in [0 ; 5]$). Le bénéfice, en centaines de milliers de francs CFA, est modélisé par
\[ B(x) = -x^3 + 6x^2 - 9x + 4. \]
Déterminer la production qui maximise le bénéfice hebdomadaire, et la valeur de ce bénéfice maximal.
\tcblower
\textbf{Solution détaillée.}
\par\textbf{Étape 1.} $B$ est un polynôme, dérivable sur $[0 ; 5]$.
\par\textbf{Étape 2.} $B'(x) = -3x^2 + 12x - 9 = -3(x^2 - 4x + 3) = -3(x-1)(x-3)$.
\par\textbf{Étape 3.} $B'(x) = 0 \iff x = 1$ ou $x = 3$. Signe de $B'$ sur $[0 ; 5]$ : $B'(x) < 0$ sur $[0 ; 1[$, $B'(x) > 0$ sur $]1 ; 3[$, $B'(x) < 0$ sur $]3 ; 5]$.
\par\textbf{Étape 4.} Valeurs :
\begin{align*}
B(0) &= 4, & B(1) &= -1 + 6 - 9 + 4 = 0, \\
B(3) &= -27 + 54 - 27 + 4 = 4, & B(5) &= -125 + 150 - 45 + 4 = -16.
\end{align*}
Tableau de variations :
\begin{center}
\begin{tabular}{|c|ccccccc|}\hline
$x$ & $0$ & & $1$ & & $3$ & & $5$ \\ \hline
$B'(x)$ & & $-$ & $0$ & $+$ & $0$ & $-$ & \\ \hline
$B$ & $4$ & $\searrow$ & $0$ & $\nearrow$ & $4$ & $\searrow$ & $-16$ \\ \hline
\end{tabular}
\end{center}
\par\textbf{Comparaison.} Candidats au maximum : $B(0) = 4$, $B(3) = 4$ ; candidats au minimum : $B(1) = 0$, $B(5) = -16$.
\par\textbf{Conclusion.} Le bénéfice maximal est de $4$ centaines de milliers de francs, soit \SI{400000}{} FCFA, atteint pour $x = 0$ (pas de production, cas limite) ou pour $x = 3$, c'est-à-dire \textbf{300 pièces par semaine}. Au-delà de 300 pièces, les coûts dégradent le bénéfice ; à 500 pièces, l'atelier perd de l'argent ($B(5) = -16$).
\end{exoresolu}

\begin{exercice}[À toi de jouer]
Soit $g$ la fonction définie sur $\mathbb{R}$ par $g(x) = x^3 - 12x + 5$.
\begin{enumerate}
\item Calculer $g'(x)$ et étudier son signe.
\item Dresser le tableau de variations de $g$ sur $\mathbb{R}$.
\item $g$ admet-elle des extremums locaux ? Préciser leur nature et leur valeur.
\item Déterminer le maximum et le minimum de $g$ sur l'intervalle $[-3 ; 4]$.
\end{enumerate}
\end{exercice}

\begin{corrige}[Exercice 1]
\begin{enumerate}
\item $g'(x) = 3x^2 - 12 = 3(x^2 - 4) = 3(x-2)(x+2)$. Donc $g'(x) > 0$ sur $]-\infty ; -2[$ et $]2 ; +\infty[$, et $g'(x) < 0$ sur $]-2 ; 2[$.
\item $g(-2) = -8 + 24 + 5 = 21$ et $g(2) = 8 - 24 + 5 = -11$. $g$ est croissante sur $]-\infty ; -2]$, décroissante sur $[-2 ; 2]$, croissante sur $[2 ; +\infty[$.
\item $g'$ s'annule en changeant de signe en $-2$ (de $+$ à $-$) : maximum local $g(-2) = 21$ ; et en $2$ (de $-$ à $+$) : minimum local $g(2) = -11$.
\item Sur $[-3 ; 4]$ : $g(-3) = -27 + 36 + 5 = 14$, $g(4) = 64 - 48 + 5 = 21$. En comparant avec $g(-2) = 21$ et $g(2) = -11$ : le maximum de $g$ sur $[-3 ; 4]$ vaut $21$ (atteint en $-2$ et en $4$) et le minimum vaut $-11$ (atteint en $2$).
\end{enumerate}
\end{corrige}

\begin{aretenir}[L'essentiel de la section]
\begin{itemize}
\item $f' \geq 0$ sur $I$ $\Rightarrow$ $f$ croissante sur $I$ ; $f' \leq 0$ $\Rightarrow$ $f$ décroissante ; $f' = 0$ $\Rightarrow$ $f$ constante (théorème admis).
\item Extremum local en un point intérieur dérivable $\Rightarrow$ $f'(c) = 0$ (condition \textbf{nécessaire} mais \textbf{pas suffisante} : penser à $x^3$ en $0$).
\item Un extremum est confirmé quand $f'$ \textbf{s'annule en changeant de signe}.
\item Tableau de variations en 4 étapes : ensemble, dérivée, signe de $f'$, tableau avec les valeurs remarquables.
\item Sur un segment $[a ; b]$, comparer toujours les extremums locaux avec $f(a)$ et $f(b)$.
\end{itemize}
\end{aretenir}

\begin{savaistu}[Optimisation : un enjeu industriel]
La recherche d'extremums par la dérivée est au cœur de l'ingénierie moderne : dimensionner une tôle pour fabriquer une caisse de volume maximal, choisir la vitesse d'un camion qui minimise la consommation de carburant sur l'axe Douala--Yaoundé, régler la tension d'un circuit pour un rendement maximal. Au Cameroun, les bureaux d'études techniques utilisent quotidiennement ces outils, souvent via des logiciels qui... calculent des dérivées ! Le mathématicien français Joseph-Louis Lagrange (1736--1813) a donné à la dérivée son rôle central dans l'étude des variations ; c'est à lui que l'on doit la notation $f'$ que tu utilises.
\end{savaistu}

\coursec{Applications à des situations concrètes}

Dans la section précédente, nous avons établi le lien fondamental entre le signe de la fonction dérivée $f'$ et le sens de variation de la fonction $f$. Ce résultat est un théorème central du programme, exigible au Probatoire.

\begin{propriete}[Dérivée et sens de variation -- Théorème fondamental]
Soit $f$ une fonction dérivable sur un intervalle $I$.
\begin{itemize}
\item Si $f'(x) \geq 0$ sur $I$, alors $f$ est croissante sur $I$.
\item Si $f'(x) \leq 0$ sur $I$, alors $f$ est décroissante sur $I$.
\item Si $f'(x) > 0$ sur $I$ (sauf éventuellement en un nombre fini de points), alors $f$ est \textbf{strictement} croissante sur $I$.
\item Si $f'(x) < 0$ sur $I$ (sauf éventuellement en un nombre fini de points), alors $f$ est \textbf{strictement} décroissante sur $I$.
\end{itemize}
\emph{Démonstration exigible au Probatoire (utilisation du théorème des accroissements finis).}
\end{propriete}

Ce théorème est bien plus qu'un outil de contrôle : c'est la clé qui permet de résoudre des problèmes concrets de la vie professionnelle et quotidienne. Quel prix faut-il fixer pour que la recette soit la plus grande possible ? Quelle quantité de tôle faut-il pour fabriquer une boîte de volume donné avec le moins de matériau possible ? À quelle vitesse roule exactement un moto-taxi à l'instant $t = 5$ secondes ? Toutes ces questions se traduisent en langage mathématique par une fonction, puis se résolvent par la dérivation. C'est l'objet de cette dernière section.

\courssub{Modéliser une situation par une fonction}

\begin{definition}[Modélisation]
\cle{Modéliser} une situation concrète, c'est la traduire en langage mathématique : on choisit une \cle{variable} $x$ (prix, longueur, temps, quantité...), on précise son \cle{ensemble de définition} $I$ imposé par le contexte, et on exprime la grandeur étudiée (recette, coût, aire, distance...) comme une \cle{fonction} $f$ de $x$, définie sur $I$.
\end{definition}

L'intuition est simple : dans un problème concret, une grandeur dépend d'une autre. La recette de Mama Odile dépend du prix de vente ; l'aire d'un champ dépend de sa largeur ; la distance parcourue dépend du temps écoulé. Dès que l'on écrit cette dépendance sous la forme $y = f(x)$, toute la machinerie de la dérivation devient disponible.

\begin{attention}[L'ensemble de définition vient du contexte]
La variable n'est pas libre de prendre n'importe quelle valeur réelle ! Un prix est positif, une longueur est positive, une quantité fabriquée est un entier positif, un temps est positif. Par exemple, si la recette est $R(p) = p(200 - 4p)$, le prix $p$ doit vérifier $p \geq 0$ et $200 - 4p \geq 0$ (la quantité vendue ne peut pas être négative), donc $p \in [0\,;\,50]$. \textbf{Oublier cette étape est la première cause d'erreur} dans les problèmes d'optimisation.
\end{attention}

\courssub{Problèmes d'optimisation : maximiser ou minimiser}

\begin{definition}[Problème d'optimisation]
Un problème d'\cle{optimisation} consiste à chercher la valeur de la variable qui rend une fonction \cle{maximale} (recette, bénéfice, aire, volume) ou \cle{minimale} (coût, quantité de matériau, distance, temps de parcours) sur son ensemble de définition.
\end{definition}

L'idée géométrique est parlante : sur la courbe de la fonction, le maximum correspond au point le plus haut, le minimum au point le plus bas. Or, en un sommet ou en un creux, la tangente est horizontale, c'est-à-dire que le nombre dérivé y est nul. On cherche donc les valeurs qui annulent $f'$, puis on vérifie à l'aide du tableau de variations qu'il s'agit bien d'un maximum (ou d'un minimum) sur l'ensemble de définition.

\begin{methode}[Résoudre un problème d'optimisation en 5 étapes]
\begin{enumerate}
\item \textbf{Choisir la variable} : nommer la grandeur inconnue (par exemple $x$ ou $p$) et déterminer son ensemble de définition $I$ à partir des contraintes concrètes.
\item \textbf{Construire la fonction} : exprimer la grandeur à optimiser en fonction de la variable, $f(x)$, en utilisant les relations de l'énoncé (périmètre fixé, prix unitaire, dimensions...).
\item \textbf{Dériver} : calculer $f'(x)$ à l'aide des formules de dérivation.
\item \textbf{Étudier les variations} : résoudre $f'(x) = 0$, étudier le signe de $f'$ sur $I$, dresser le tableau de variations et repérer le maximum ou le minimum.
\item \textbf{Conclure dans le contexte} : vérifier que la valeur trouvée appartient bien à $I$, puis rédiger une phrase de conclusion qui répond à la question posée, avec les unités.
\end{enumerate}
\end{methode}

\begin{exoresolu}[La recette de Mama Odile]
Mama Odile vend des beignets au marché de Mokolo. Elle a remarqué que si elle fixe le prix unitaire à $p$ francs CFA, elle en vend $200 - 4p$ par jour. On note $R(p)$ sa recette journalière, en francs CFA. Déterminer le prix $p$ qui rend la recette maximale, et calculer cette recette maximale.
\tcblower
\textbf{Étape 1 — Variable et ensemble de définition.} La variable est le prix $p$. On doit avoir $p \geq 0$ et $200 - 4p \geq 0$, c'est-à-dire $p \leq 50$. Donc $p \in [0\,;\,50]$.

\textbf{Étape 2 — Fonction.} La recette est le prix multiplié par la quantité vendue :
\[ R(p) = p(200 - 4p) = -4p^{2} + 200p. \]

\textbf{Étape 3 — Dérivée.} $R$ est une fonction polynôme, dérivable sur $\mathbb{R}$, donc sur $[0\,;\,50]$ :
\[ R'(p) = -8p + 200. \]

\textbf{Étape 4 — Variations.} $R'(p) = 0 \iff -8p + 200 = 0 \iff p = 25$. De plus, $R'(p) > 0$ pour $p < 25$ et $R'(p) < 0$ pour $p > 25$. D'où le tableau de variations :
\begin{center}
\begin{tabular}{|c|ccccc|}\hline
$p$ & $0$ & & $25$ & & $50$ \\ \hline
$R'(p)$ & & $+$ & $0$ & $-$ & \\ \hline
$R$ & $0$ & $\nearrow$ & $2500$ & $\searrow$ & $0$ \\ \hline
\end{tabular}
\end{center}
avec $R(25) = 25 \times (200 - 4 \times 25) = 25 \times 100 = 2500$.

\textbf{Étape 5 — Conclusion.} La valeur $p = 25$ appartient bien à $[0\,;\,50]$. D'après le tableau de variations, $R$ admet sur cet intervalle un maximum atteint en $p = 25$. \textbf{Mama Odile doit fixer le prix du beignet à $25$~FCFA ; sa recette journalière maximale est alors de $2500$~FCFA.}
\end{exoresolu}

\begin{popfigure}
\begin{center}
\begin{tikzpicture}
\begin{axis}[
    width=0.85\linewidth, height=6.5cm,
    axis lines=middle,
    xlabel={$p$ (prix en FCFA)}, ylabel={$R(p)$ (recette en FCFA)},
    xmin=-3, xmax=55, ymin=-300, ymax=3200,
    xtick={0,10,20,25,30,40,50},
    ytick={0,1000,2000,2500,3000},
    grid=both, grid style={popline!40},
    tick label style={font=\small}
]
\addplot[popcurve,domain=0:50,samples=200]{-4*x^2+200*x};
\addplot[mark=*,mark size=2.5pt,color=poporange] coordinates {(25,2500)};
\draw[dashed,poporange] (axis cs:25,0) -- (axis cs:25,2500) -- (axis cs:0,2500);
\node[poporange,anchor=south west,font=\small] at (axis cs:26,2550) {Sommet $(25\,;\,2500)$};
\end{axis}
\end{tikzpicture}
\end{center}
\legende{Parabole de la recette $R(p) = p(200 - 4p)$ : le sommet correspond au maximum, atteint pour $p = 25$ FCFA.}
\end{popfigure}

\begin{exoresolu}[Clôturer un champ contre la rivière]
Un planteur de la région de Bafoussam veut clôturer un champ rectangulaire adossé à une rivière (le côté le long de la rivière n'a pas besoin de clôture). Il dispose de $120$ m de grillage. Quelles dimensions doit-il choisir pour que l'aire du champ soit maximale ?
\tcblower
\textbf{Étape 1.} Notons $x$ la largeur du champ (côté perpendiculaire à la rivière) et $y$ sa longueur (côté parallèle à la rivière), en mètres. On a $x > 0$ et $y > 0$.

\textbf{Étape 2.} Le grillage sert pour deux largeurs et une longueur : $2x + y = 120$, donc $y = 120 - 2x$. Comme $y > 0$, on a $x < 60$, donc $x \in ]0\,;\,60[$. L'aire est :
\[ A(x) = x \times y = x(120 - 2x) = -2x^{2} + 120x. \]

\textbf{Étape 3.} $A$ est dérivable sur $]0\,;\,60[$ et $A'(x) = -4x + 120$.

\textbf{Étape 4.} $A'(x) = 0 \iff x = 30$. $A'(x) > 0$ pour $x < 30$ et $A'(x) < 0$ pour $x > 30$ :
\begin{center}
\begin{tabular}{|c|ccccc|}\hline
$x$ & $0$ & & $30$ & & $60$ \\ \hline
$A'(x)$ & & $+$ & $0$ & $-$ & \\ \hline
$A$ & $0$ & $\nearrow$ & $1800$ & $\searrow$ & $0$ \\ \hline
\end{tabular}
\end{center}
avec $A(30) = 30 \times (120 - 60) = 30 \times 60 = 1800$.

\textbf{Étape 5.} $x = 30 \in ]0\,;\,60[$ : la valeur est acceptable. L'aire est maximale pour une largeur de $30$ m et une longueur de $120 - 2 \times 30 = 60$ m. \textbf{L'aire maximale du champ est $1800$~m$^{2}$.}
\end{exoresolu}

\begin{popfigure}
\begin{center}
\begin{tikzpicture}[scale=0.9]
% Rivière
\fill[popblueL] (0,3.2) rectangle (8,4.6);
\draw[decorate,decoration={snake,amplitude=1.5pt,segment length=12pt},popblue,thick] (0,3.5) -- (8,3.5);
\draw[decorate,decoration={snake,amplitude=1.5pt,segment length=12pt},popblue,thick] (0,4.2) -- (8,4.2);
\node[popblue,font=\small\itshape] at (4,3.9) {Rivière (pas de clôture)};
% Champ
\fill[popgreenL] (1.5,0) rectangle (6.5,3.2);
\draw[popgreen,very thick] (1.5,0) rectangle (6.5,3.2);
\draw[popgreen,very thick,dashed] (1.5,3.2) -- (6.5,3.2);
% Cotes
\draw[<->,popdark] (1.5,-0.4) -- (6.5,-0.4) node[midway,below,font=\small] {$y$ (longueur)};
\draw[<->,popdark] (1.1,0) -- (1.1,3.2) node[midway,left,font=\small] {$x$};
\draw[<->,popdark] (6.9,0) -- (6.9,3.2) node[midway,right,font=\small] {$x$};
\node[popdark,font=\small] at (4,1.6) {Champ : aire $A = x \times y$};
\node[popdark,font=\small] at (4,-1) {Grillage utilisé : $2x + y = 120$ m};
\end{tikzpicture}
\end{center}
\legende{Champ rectangulaire adossé à une rivière : seuls trois côtés sont clôturés.}
\end{popfigure}

\begin{attention}[Vérifier que la solution appartient à l'ensemble de définition]
Erreur fréquente à l'examen : on résout $f'(x) = 0$, on trouve une valeur, et on conclut immédiatement. Mais si cette valeur \textbf{n'appartient pas} à l'ensemble de définition imposé par le contexte, elle ne peut pas être la réponse ! Par exemple, si l'on cherche une longueur et que l'équation $f'(x) = 0$ donne $x = -5$, cette valeur est à rejeter. Il faut alors étudier les variations de $f$ sur l'intervalle autorisé : l'extremum est souvent atteint à une borne de l'intervalle. \textbf{Toujours conclure en vérifiant : la valeur trouvée est-elle acceptable pour le problème concret ?}
\end{attention}

\courssub{Vitesse moyenne et vitesse instantanée}

Imaginons un moto-taxi (benskin) qui part du carrefour et dont la position sur une route droite est repérée, à l'instant $t$ (en secondes), par le nombre $x(t)$ (en mètres). Entre les instants $t$ et $t + h$, il parcourt la distance $x(t+h) - x(t)$ pendant la durée $h$.

\begin{definition}[Vitesse moyenne et vitesse instantanée]
Soit $x(t)$ la position d'un mobile à l'instant $t$.
\begin{itemize}
\item La \cle{vitesse moyenne} entre les instants $t$ et $t + h$ (avec $h \neq 0$) est :
\[ v_{\text{moy}} = \frac{x(t+h) - x(t)}{h}. \]
\item La \cle{vitesse instantanée} à l'instant $t$ est la limite de la vitesse moyenne quand $h$ tend vers $0$, c'est-à-dire le nombre dérivé de $x$ en $t$ :
\[ v(t) = x'(t). \]
\end{itemize}
\end{definition}

L'intuition est celle du compteur de vitesse : la vitesse moyenne sur un trajet Yaoundé--Douala ne dit rien de la vitesse à un instant précis ; la vitesse instantanée, elle, est celle qu'affiche le compteur à l'instant $t$. Mathématiquement, on reconnaît dans la vitesse moyenne le \textbf{taux d'accroissement} de la fonction position : passer à la limite, c'est exactement dériver.

\begin{exemplebox}[Vitesse instantanée d'un moto-taxi]
La position d'un moto-taxi est donnée par $x(t) = 2t^{2} + 3t$ (en mètres, $t$ en secondes).
\begin{enumerate}
\item Vitesse moyenne entre $t = 5$ et $t = 6$ :
\[ v_{\text{moy}} = \frac{x(6) - x(5)}{1} = \frac{(72 + 18) - (50 + 15)}{1} = 90 - 65 = 25 \text{ m/s}. \]
\item Vitesse instantanée : $x$ est dérivable sur $\mathbb{R}$ et $v(t) = x'(t) = 4t + 3$. À l'instant $t = 5$ :
\[ v(5) = 4 \times 5 + 3 = 23 \text{ m/s}. \]
\end{enumerate}
On remarque que la vitesse moyenne ($25$ m/s) est proche de la vitesse instantanée ($23$ m/s) : plus l'intervalle de temps est petit, meilleure est l'approximation.
\end{exemplebox}

\courssub{Coût marginal en économie}

\begin{definition}[Coût marginal]
Une entreprise produit $n$ unités d'un article pour un coût total $C(n)$. Le \cle{coût marginal} de la $(n+1)$-ième unité est le coût supplémentaire pour produire une unité de plus :
\[ C_{m}(n) = C(n+1) - C(n). \]
Lorsque $n$ est grand, on approche ce coût marginal par le nombre dérivé :
\[ C_{m}(n) \approx C'(n). \]
\end{definition}

D'où vient cette approximation ? De l'\textbf{approximation affine} vue dans la section précédente : au voisinage de $n$,
\[ C(n + h) \approx C(n) + h \times C'(n). \]
En prenant $h = 1$, on obtient $C(n+1) - C(n) \approx C'(n)$. Le nombre dérivé $C'(n)$ donne donc une bonne estimation du coût de production d'une unité supplémentaire.

\begin{exoresolu}[Coût marginal d'un menuisier]
Un menuisier de Douala fabrique des chaises. Le coût total de production de $n$ chaises est, en milliers de FCFA :
\[ C(n) = 0{,}1n^{2} + 2n + 10. \]
\begin{enumerate}
\item Calculer le coût marginal exact de la $21^{\text{e}}$ chaise, $C(21) - C(20)$.
\item Calculer $C'(20)$ et comparer.
\end{enumerate}
\tcblower
\textbf{1.} $C(21) = 0{,}1 \times 441 + 42 + 10 = 96{,}1$ et $C(20) = 0{,}1 \times 400 + 40 + 10 = 90$. Donc :
\[ C_{m}(20) = C(21) - C(20) = 96{,}1 - 90 = 6{,}1 \text{ milliers de FCFA, soit } 6100 \text{ FCFA}. \]
\textbf{2.} $C$ est dérivable et $C'(n) = 0{,}2n + 2$, donc $C'(20) = 4 + 2 = 6$ milliers de FCFA, soit $6000$ FCFA.

\textbf{Comparaison.} L'approximation $C'(20) = 6000$ FCFA est très proche du coût exact $6100$ FCFA (écart de $100$ FCFA seulement). Le menuisier peut donc utiliser $C'(n)$ pour décider rapidement si produire une chaise de plus est rentable.
\end{exoresolu}

\begin{savaistu}[Le coût marginal, outil de décision des entreprises]
En économie, une règle de décision fondamentale affirme qu'une entreprise a intérêt à augmenter sa production tant que le coût marginal reste inférieur au prix de vente : chaque unité supplémentaire rapporte alors plus qu'elle ne coûte. Les entreprises camerounaises — brasseries, minoteries, ateliers de confection — utilisent ce raisonnement pour fixer leurs quantités produites et maximiser leur bénéfice. Les mathématiciens et économistes ont formalisé cette idée dès le XIX$^{\text{e}}$ siècle, et elle repose entièrement sur la notion de dérivée que tu viens d'étudier !
\end{savaistu}

\courssub{Rédaction type examen}

Au Probatoire et au Baccalauréat technique, un problème d'optimisation ou de modélisation est noté autant sur la \textbf{rédaction} que sur les calculs. Voici la structure attendue.

\begin{methode}[Rédiger une réponse complète et justifiée]
\begin{enumerate}
\item \textbf{Annoncer} clairement ce que l'on cherche et nommer la variable : « Soit $x$ la largeur du champ, en mètres. »
\item \textbf{Préciser l'ensemble de définition} en le justifiant par le contexte : « Comme $x > 0$ et $120 - 2x > 0$, on a $x \in ]0\,;\,60[$. »
\item \textbf{Justifier la dérivabilité} avant de dériver : « $A$ est une fonction polynôme, donc dérivable sur $]0\,;\,60[$, et $A'(x) = \ldots$ »
\item \textbf{Justifier l'extremum par le théorème} : « $A'(x) > 0$ sur $]0\,;\,30[$ et $A'(x) < 0$ sur $]30\,;\,60[$, donc $A$ est croissante puis décroissante : elle admet un maximum en $x = 30$. »
\item \textbf{Conclure dans le contexte}, avec une phrase complète et les unités : « Le planteur doit choisir une largeur de $30$ m et une longueur de $60$ m ; l'aire maximale est $1800$ m$^{2}$. »
\end{enumerate}
\end{methode}

\begin{exercice}[Pour s'entraîner]
Un artisan fabrique des gâteaux qu'il vend $q$ centaines de francs l'unité. Sa recette pour $q$ est $R(q) = -2q^{2} + 80q$ (en centaines de FCFA), avec $q \in [0\,;\,40]$.
\begin{enumerate}
\item Étudier les variations de $R$ sur $[0\,;\,40]$.
\item En déduire le prix de vente qui maximise la recette, et la recette maximale.
\end{enumerate}
\end{exercice}

\begin{corrige}[Exercice]
\textbf{1.} $R$ est une fonction polynôme, dérivable sur $[0\,;\,40]$, et $R'(q) = -4q + 80$. On a $R'(q) = 0 \iff q = 20$ ; $R'(q) > 0$ pour $q < 20$ et $R'(q) < 0$ pour $q > 20$. Donc $R$ est strictement croissante sur $[0\,;\,20]$ et strictement décroissante sur $[20\,;\,40]$.

\textbf{2.} $q = 20 \in [0\,;\,40]$ : d'après le sens de variation, $R$ admet un maximum en $q = 20$, qui vaut $R(20) = -2 \times 400 + 80 \times 20 = 800$. L'artisan doit vendre le gâteau $2000$~FCFA ; sa recette maximale est alors de $80\,000$~FCFA.
\end{corrige}

\begin{aretenir}[L'essentiel de la section]
\begin{itemize}
\item Modéliser, c'est choisir une variable, déterminer son ensemble de définition à partir du contexte, puis construire la fonction $f$.
\item Pour optimiser : dériver, résoudre $f'(x) = 0$, dresser le tableau de variations, \textbf{vérifier que la solution appartient à l'ensemble de définition}, puis conclure dans le contexte avec les unités.
\item La vitesse instantanée est la dérivée de la position : $v(t) = x'(t)$.
\item Le coût marginal $C(n+1) - C(n)$ est approché par $C'(n)$ grâce à l'approximation affine.
\item À l'examen : annoncer, justifier la dérivabilité, invoquer le théorème (signe de $f'$ et variations), conclure par une phrase rédigée.
\end{itemize}
\end{aretenir}

\newpage

\coursec{Activité d'intégration}

\coursec{Leçon 14 : Dérivation — Application : Optimisation d'une clôture}

\courssub{Objectifs de l'activité}
Cette activité d'intégration vise à mobiliser l'ensemble des savoirs et savoir-faire de la leçon sur la dérivation dans une situation concrète d'optimisation. Elle permet à l'élève de :
\begin{itemize}
    \item Modéliser une situation réelle (clôture d'un champ) par une fonction numérique $A(x)$ sur un intervalle borné.
    \item Déterminer l'ensemble de définition physique de la fonction et calculer sa dérivée en appliquant les règles de dérivation (somme, produit, composée affine).
    \item Étudier le sens de variation de la fonction à l'aide d'un tableau de variations rigoureux.
    \item Résoudre un problème d'optimisation (recherche d'un maximum) et interpréter le résultat dans le contexte.
    \item Utiliser l'approximation affine (différentielle) pour estimer l'erreur commise en s'écartant de la valeur optimale, et comprendre les limites de cette approximation au voisinage d'un extremum.
    \item Rédiger une conclusion complète, justifiée et structurée, conforme aux attendus de l'examen du Probatoire.
\end{itemize}

\courssub{Situation}
À \textbf{Bafoussam}, Grand-Père \textsc{Noubi} souhaite aménager un potager rectangulaire pour cultiver du maïs et des haricots. Il dispose de \textbf{100 mètres de fil barbelé} pour réaliser la clôture. Son terrain longe un grand mur en parpaings qui servira de limite naturelle pour l'un des grands côtés du rectangle : il n'y a donc \emph{pas besoin de fil barbelé} du côté du mur.

Grand-Père \textsc{Noubi} veut que la surface cultivable soit la \textbf{plus grande possible} pour maximiser sa récolte. Il se demande quelles dimensions (largeur et longueur) il doit donner à son champ.

On note $x$ la \textbf{largeur} du champ (côté perpendiculaire au mur), exprimée en mètres.

\begin{popfigure}
\begin{tikzpicture}[scale=0.7, every node/.style={font=\small}]
    % Mur
    \draw[popdark, line width=2pt, pattern=north east lines, pattern color=popdark] (0,0) -- (8,0);
    \node[popdark, above] at (4,0.3) {\textbf{Mur (pas de clôture)}};
    % Clôture : 3 côtés
    \draw[poporange, line width=2pt, -{Stealth[length=3mm]}] (0,0) -- (0,-4) node[midway, left] {$x$};
    \draw[poporange, line width=2pt, -{Stealth[length=3mm]}] (0,-4) -- (8,-4) node[midway, below] {$L$};
    \draw[poporange, line width=2pt, -{Stealth[length=3mm]}] (8,-4) -- (8,0) node[midway, right] {$x$};
    % Fil barbelé (représentation symbolique)
    \draw[poporange, decorate, decoration={snake, amplitude=1pt, segment length=4pt}] (0,0) -- (0,-4) -- (8,-4) -- (8,0);
    \node[poporange, align=center] at (4,-5.5) {Longueur totale du fil barbelé : \SI{100}{m}};
\end{tikzpicture}
\legende{Configuration du champ de Grand-Père Noubi : le mur remplace un côté de la clôture.}
\end{popfigure}

\courssub{Consignes}
\begin{enumerate}
    \item \textbf{Modélisation :} Exprimer la longueur $L$ du champ en fonction de $x$. En déduire l'expression de l'aire $A(x)$ du champ en fonction de $x$.
    \item \textbf{Ensemble de définition :} Déterminer l'intervalle $I$ des valeurs possibles pour $x$ (contraintes physiques : longueurs positives). Préciser $A(x)$ sur $I$.
    \item \textbf{Dérivation :} Montrer que la fonction $A$ est dérivable sur $I$ et calculer sa fonction dérivée $A'(x)$.
    \item \textbf{Tableau de variations :} Dresser le tableau de variations complet de $A$ sur $I$.
    \item \textbf{Optimisation :} Déterminer la largeur $x_0$ qui maximise l'aire. Calculer l'aire maximale $A_{\max}$ et la longueur $L_0$ correspondante. Donner les dimensions optimales du champ.
    \item \textbf{Approximation affine (Prolongement) :} 
    \begin{enumerate}
        \item Écrire la formule de l'approximation affine de $A$ au voisinage de $x_0 = 25$.
        \item Grand-Père \textsc{Noubi} hésite et pense faire une largeur de $26$ m au lieu de $25$ m. Utiliser l'approximation affine pour estimer l'aire $A(26)$.
        \item Calculer la valeur exacte $A(26)$ et comparer avec l'estimation. Quelle est l'aire \emph{perdue} par rapport au maximum ? Commenter la précision de l'approximation affine au voisinage d'un extremum.
    \end{enumerate}
    \item \textbf{Rédaction finale :} Rédiger une réponse complète et justifiée (comme à l'examen) résolvant le problème initial de Grand-Père \textsc{Noubi}.
\end{enumerate}

\courssub{Ressources à mobiliser}
\begin{definition}[Nombre dérivé en un point]
Soit $f$ une fonction définie sur un intervalle $I$ et $a \in I$. Si la limite 
\[ \lim_{x \to a} \frac{f(x) - f(a)}{x - a} \]
existe et est finie, on dit que $f$ est \cle{dérivable} en $a$. Cette limite est appelée \cle{nombre dérivé} de $f$ en $a$ et se note $f'(a)$.
\end{definition}

\begin{propriete}[Dérivabilité et opérations sur un intervalle]
Soient $f$ et $g$ dérivables sur un intervalle $I$.
\begin{itemize}
    \item $f+g$ est dérivable sur $I$ et $(f+g)' = f' + g'$.
    \item $f \times g$ est dérivable sur $I$ et $(fg)' = f'g + fg'$.
    \item Si $u(x) = ax+b$ (affine) et $f$ dérivable sur $I$, alors $f \circ u$ est dérivable et $(f \circ u)'(x) = a f'(ax+b)$.
\end{itemize}
Fonctions usuelles dérivables sur $\mathbb{R}$ : $x \mapsto x^n$ ($n \in \mathbb{N}^*$) avec $(x^n)' = n x^{n-1}$.
\end{propriete}

\begin{propriete}[Dérivée et sens de variation]
Soit $f$ dérivable sur un intervalle $I$.
\begin{itemize}
    \item Si $f'(x) \ge 0$ sur $I$, alors $f$ est croissante sur $I$.
    \item Si $f'(x) \le 0$ sur $I$, alors $f$ est décroissante sur $I$.
    \item Si $f'(x) = 0$ sur $I$, alors $f$ est constante sur $I$.
\end{itemize}
Un extremum (maximum ou minimum) ne peut se produire qu'en un point où $f'(x)=0$ ou aux bornes de l'intervalle. Le changement de signe de $f'$ (de $+$ à $-$ pour un maximum, de $-$ à $+$ pour un minimum) permet de conclure.
\end{propriete}

\begin{propriete}[Approximation affine (Différentielle)]
Si $f$ est dérivable en $x_0$, au voisinage de $x_0$ :
\[ f(x_0 + h) \approx f(x_0) + h f'(x_0) \quad \text{(pour $h$ petit)}. \]
L'erreur commise est un infinitésimal d'ordre supérieur à $h$ : $f(x_0+h) - [f(x_0) + h f'(x_0)] = o(h)$.
\end{propriete}

\begin{methode}[Résolution d'un problème d'optimisation]
\begin{enumerate}
    \item Introduire la variable $x$ et exprimer la grandeur à optimiser $f(x)$.
    \item Déterminer l'ensemble de définition $I$ (contraintes du réel).
    \item Étudier la dérivabilité de $f$ sur $I$ et calculer $f'(x)$.
    \item Résoudre $f'(x)=0$ dans $I$ et dresser le tableau de variations.
    \item Conclure : valeur optimale, point où elle est atteinte, interprétation concrète avec unités.
\end{enumerate}
\end{methode}

\courssub{Démarche et corrigé commenté}
\begin{exoresolu}{La clôture optimale de Grand-Père Noubi}
\textbf{1. Modélisation}
Le périmètre clôturé utilise le fil barbelé pour les deux largeurs ($x$ chacune) et une longueur ($L$).
\[ 2x + L = 100 \quad \Rightarrow \quad L(x) = 100 - 2x. \]
L'aire $A$ du rectangle est le produit de la largeur par la longueur :
\[ A(x) = x \times L(x) = x(100 - 2x) = 100x - 2x^2. \]

\textbf{2. Ensemble de définition $I$}
Les dimensions physiques imposent $x > 0$ (largeur strictement positive) et $L(x) > 0$ (longueur strictement positive).
\[ L(x) > 0 \iff 100 - 2x > 0 \iff 2x < 100 \iff x < 50. \]
Donc $x \in ]0; 50[$. 
Pour l'étude de la fonction continue sur un intervalle fermé (recherche de maximum global), on prolonge par continuité aux bornes : $I = [0; 50]$.
Sur $I$, $A(x) = -2x^2 + 100x$.

\textbf{3. Dérivation}
$A$ est une fonction polynôme du second degré, elle est dérivable sur $\mathbb{R}$, donc sur $I = [0; 50]$.
On applique la règle de la somme et la règle de la puissance $(x^n)' = n x^{n-1}$ :
\[ A'(x) = (100x)' - (2x^2)' = 100 - 4x. \]
(On peut aussi voir $A(x) = x(100-2x)$ et utiliser $(uv)' = u'v + uv'$ : $1 \cdot (100-2x) + x \cdot (-2) = 100 - 2x - 2x = 100 - 4x$.)

\textbf{4. Tableau de variations de $A$ sur $[0; 50]$}
On étudie le signe de $A'(x) = 100 - 4x = 4(25 - x)$.
\begin{itemize}
    \item $A'(x) = 0 \iff x = 25$.
    \item $A'(x) > 0 \iff 25 - x > 0 \iff x < 25$ (sur $[0; 25[$).
    \item $A'(x) < 0 \iff x > 25$ (sur $]25; 50]$).
\end{itemize}
$A$ est continue sur $[0; 50]$, dérivable sur $]0; 50[$.
Valeurs aux points remarquables :
$A(0) = 0$ ; $A(25) = 100 \times 25 - 2 \times 25^2 = 2500 - 1250 = 1250$ ; $A(50) = 0$.

\begin{center}
\begin{tabular}{|c|ccccc|}\hline
$x$ & $0$ & & $25$ & & $50$ \\ \hline
$A'(x)$ & & $+$ & $0$ & $-$ & \\ \hline
$A(x)$ & $0$ & $\nearrow$ & $1250$ & $\searrow$ & $0$ \\ \hline
\end{tabular}
\end{center}

\textbf{5. Optimisation}
D'après le tableau, $A$ admet un \textbf{maximum} en $x_0 = 25$.
La valeur maximale est $A_{\max} = A(25) = \SI{1250}{m^2}$.
La longueur correspondante est $L_0 = L(25) = 100 - 2 \times 25 = \SI{50}{m}$.
\textbf{Dimensions optimales :} Largeur $= \SI{25}{m}$, Longueur $= \SI{50}{m}$ (le long du mur).

\textbf{6. Approximation affine au voisinage de $x_0 = 25$}
\begin{itemize}
    \item[a)] Formule : $A(25+h) \approx A(25) + h A'(25)$.
    On a $A(25) = 1250$ et $A'(25) = 100 - 4 \times 25 = 0$.
    Donc $A(25+h) \approx 1250 + h \times 0 = 1250$.
    \item[b)] Pour $x = 26$, on a $h = 1$. L'estimation affine donne : $A(26) \approx \SI{1250}{m^2}$.
    \item[c)] Valeur exacte : $A(26) = 100 \times 26 - 2 \times 26^2 = 2600 - 1352 = \SI{1248}{m^2}$.
    L'aire \textbf{perdue} (réelle) par rapport au maximum est $A(25) - A(26) = 1250 - 1248 = \SI{2}{m^2}$.
    L'approximation affine prédit une aire de $\SI{1250}{m^2}$ (perte nulle), alors que la perte réelle est de $\SI{2}{m^2}$.
    \textbf{Commentaire :} Au voisinage d'un extremum, la dérivée est nulle. L'approximation affine (qui est la tangente) est \emph{horizontale}. Elle ne « voit » pas la courbure de la parabole (liée à la dérivée seconde). L'erreur relative est faible ici ($\num{0,16}\%$), mais l'approximation affine \textbf{ne permet pas d'estimer la variation (la perte) au premier ordre} près d'un extremum ; il faudrait un développement d'ordre 2 (non au programme de Première).
\end{itemize}

\textbf{7. Rédaction finale (Type Probatoire)}
\begin{quote}
\textbf{Énoncé :} Grand-Père Noubi dispose de \SI{100}{m} de grillage pour clôturer un champ rectangulaire adossé à un mur. Quelles dimensions choisir pour maximiser la surface ?

\textbf{Résolution :}
Soit $x$ la largeur du champ (en m). La longueur est $L = 100 - 2x$.
L'aire est $A(x) = x(100-2x) = -2x^2 + 100x$.
Les contraintes $x>0$ et $L>0$ donnent $x \in [0; 50]$.
$A$ est polynôme, donc dérivable sur $[0; 50]$ avec $A'(x) = 100 - 4x = 4(25-x)$.
$A'(x) = 0 \iff x = 25$.
$A'(x) > 0$ sur $[0; 25[$ et $A'(x) < 0$ sur $]25; 50]$.
$A$ croît sur $[0; 25]$, décroît sur $[25; 50]$, avec $A(0)=0$, $A(25)=1250$, $A(50)=0$.
Le maximum de $A$ sur $[0; 50]$ est $\SI{1250}{m^2}$, atteint pour $x=25$.
La longueur correspondante est $L = 100 - 50 = \SI{50}{m}$.

\textbf{Conclusion :} Grand-Père Noubi doit faire un champ de \textbf{\SI{25}{m} de largeur} sur \textbf{\SI{50}{m} de longueur} (le long du mur) pour obtenir la surface maximale de \textbf{\SI{1250}{m²}}.
\end{quote}
\end{exoresolu}

\courssub{Bilan de l'activité}
Cette activité a permis de mettre en évidence les points clés suivants, essentiels pour le Probatoire :
\begin{itemize}
    \item \textbf{La modélisation rigoureuse :} Passer du texte à la fonction $A(x)$, en identifiant la variable, les contraintes (domaine $I$) et l'expression algébrique.
    \item \textbf{La mécanique de la dérivation :} Application fluide des règles de dérivation (linéarité, puissance, ou produit) sur une fonction polynôme simple.
    \item \textbf{Le lien Derivée / Variations :} La construction du tableau de variations est l'outil central pour trancher entre maximum et minimum. La lecture du tableau (croissance puis décroissance $\Rightarrow$ maximum) doit être explicite.
    \item \textbf{L'interprétation physique :} Retrouver les deux dimensions (largeur ET longueur) et l'aire max, avec les unités ($\text{m}$, $\text{m}^2$).
    \item \textbf{La rédaction d'examen :} La conclusion finale doit être une réponse courte, autonome, reprenant la question, les valeurs numériques et les unités, sans calculs intermédiaires.
    \item \textbf{La portée et les limites de l'approximation affine :} L'exercice 6 montre qu'au voisinage d'un extremum ($f'(x_0)=0$), l'approximation affine est constante (tangente horizontale). Elle est précise en valeur absolue (erreur $o(h)$) mais \textbf{inopérante pour prédire la variation (la perte) au premier ordre}. C'est une excellente introduction intuitive à la notion de dérivée seconde / concavité (programme de Terminale).
\end{itemize}

\begin{attention}[Pièges à éviter pour l'examen]
\begin{itemize}
    \item \textbf{Oublier l'ensemble de définition $I$ :} Ne pas écrire $x \in [0; 50]$ (ou $]0; 50[$) fait perdre des points. Les bornes $0$ et $50$ correspondent à un champ de surface nulle (dégénéré), elles sont cruciales pour l'étude sur un segment.
    \item \textbf{Confondre $x$ (largeur) et $L$ (longueur) :} La question demande « les dimensions ». Il faut donner les deux.
    \item \textbf{Tableau de variations incomplet :} Il faut 3 lignes ($x$, $A'(x)$, $A(x)$), les bornes de $I$, la valeur annulant la dérivée ($25$), les signes de $A'$ \textbf{entre} les valeurs (pas sur les cases des valeurs), les flèches \textbf{uniquement sur la ligne de $A$}, et les valeurs images ($0, 1250, 0$).
    \item \textbf{Dire « $A'(x) > 0$ donc maximum » :} C'est faux. C'est le \textbf{changement de signe} de $A'$ (de $+$ à $-$) qui prouve le maximum.
    \item \textbf{Approximation affine :} Ne pas écrire $A(26) = A(25) + 1 \times A'(25)$ avec un « $=$ ». C'est une approximation ($\approx$ ou $\simeq$).
\end{itemize}
\end{attention}

\begin{savaistu}[Le saviez-vous ? : L'isopérimétrique dans la pratique]
Ce problème est un cas particulier du \textbf{problème isopérimétrique} : parmi tous les rectangles de périmètre donné (ou semi-périmètre donné ici $x+L=50$), celui qui a la plus grande aire est le \textbf{carré}.
Ici, le « demi-périmètre » clôturé est $x + L = 50$. Le maximum d'aire est atteint pour $x = L = 25$. Le champ optimal est donc un \textbf{carré de 25 m de côté} (le mur faisant office du 4ème côté).
En agriculture au Cameroun, cette optimisation permet d'économiser le grillage (coût élevé) ou de maximiser le rendement pour une longueur de haie vive ou de fil barbelé donnée. C'est une application directe des mathématiques à la gestion économique d'une exploitation.
\end{savaistu}

\newpage

\coursec{Méthodes et exercices résolus}

\coursec{Dérivation}

\courssub{Nombre dérivé en un réel}

\begin{definition}[Nombre dérivé en un point]
Soit $f$ une fonction définie sur un intervalle $I$ et $a \in I$. Si la limite 
\[
\lim_{h \to 0} \frac{f(a+h)-f(a)}{h}
\]
existe et est finie, on dit que $f$ est \textbf{dérivable en $a$}. Cette limite est appelée \textbf{nombre dérivé de $f$ en $a$} et se note $f'(a)$.
\par
Géométriquement, $f'(a)$ est le coefficient directeur de la tangente à la courbe $\mathcal{C}_f$ au point $A(a; f(a))$.
\end{definition}

\begin{propriete}[Dérivabilité et continuité]
Si une fonction est dérivable en un point, elle est continue en ce point.
\par
\textbf{Attention :} La réciproque est fausse. Exemple : $f(x)=|x|$ est continue en $0$ mais non dérivable en $0$.
\end{propriete}

\begin{methode}[Calculer le nombre dérivé en un point par la limite du taux d'accroissement]
\textbf{Objectif :} Déterminer $f'(a)$ pour une fonction $f$ donnée par une expression algébrique, en $a$ appartenant à son ensemble de définition.
\begin{enumerate}
    \item \textbf{Vérifier} que $a$ appartient à l'ensemble de définition $D_f$ et que $f$ est continue en $a$ (condition nécessaire).
    \item \textbf{Écrire} le taux d'accroissement : $\tau(h) = \dfrac{f(a+h)-f(a)}{h}$ pour $h \neq 0$.
    \item \textbf{Simplifier} l'expression de $\tau(h)$ (factorisation, mise au même dénominateur, conjugaison si racine) pour faire disparaître le $h$ au dénominateur.
    \item \textbf{Calculer} la limite : $f'(a) = \lim_{h \to 0} \tau(h)$.
    \item \textbf{Conclure} : « La fonction $f$ est dérivable en $a$ et $f'(a) = \dots$ » ou « $f$ n'est pas dérivable en $a$ car la limite n'existe pas (ou est infinie) ».
\end{enumerate}
\textbf{Réflexe :} Pour $f(x)=\sqrt{x}$ en $a>0$, multiplier numérateur et dénominateur par $\sqrt{a+h}+\sqrt{a}$. Pour $f(x)=\frac{1}{x}$, mettre au même dénominateur.
\end{methode}

\begin{exoresolu}[Nombre dérivé d'une fonction racine — Type Probatoire]
Soit la fonction $f$ définie sur $[0; +\infty[$ par $f(x) = \sqrt{x+4}$.
\begin{enumerate}
    \item Étudier la dérivabilité de $f$ en $0$ et calculer $f'(0)$.
    \item En déduire une valeur approchée de $\sqrt{4,01}$ par approximation affine.
\end{enumerate}
\tcblower
\textbf{Solution détaillée}
\begin{enumerate}
    \item \textbf{Étude de la dérivabilité en $0$ :}
    La fonction $f$ est définie en $0$ car $0+4=4 \ge 0$. $f(0) = \sqrt{4} = 2$.
    Pour $h > -4$ (voisinage de $0$), le taux d'accroissement est :
    \[
    \tau(h) = \frac{f(0+h)-f(0)}{h} = \frac{\sqrt{h+4} - 2}{h}.
    \]
    On rationalise le numérateur en multipliant par le conjugué $\sqrt{h+4}+2$ :
    \begin{align*}
    \tau(h) &= \frac{(\sqrt{h+4} - 2)(\sqrt{h+4}+2)}{h(\sqrt{h+4}+2)} \\
    &= \frac{(h+4) - 4}{h(\sqrt{h+4}+2)} \\
    &= \frac{h}{h(\sqrt{h+4}+2)} \quad (h \neq 0) \\
    &= \frac{1}{\sqrt{h+4}+2}.
    \end{align*}
    On calcule la limite :
    \[
    \lim_{h \to 0} \tau(h) = \lim_{h \to 0} \frac{1}{\sqrt{h+4}+2} = \frac{1}{\sqrt{4}+2} = \frac{1}{4}.
    \]
    La limite est finie, donc $f$ est dérivable en $0$ et \textbf{$f'(0) = \num{0,25}$}.
    
    \item \textbf{Approximation affine :}
    On cherche une valeur approchée de $\sqrt{4,01} = f(0,01)$.
    Formule : $f(x_0+h) \approx f(x_0) + h f'(x_0)$ avec $x_0=0$ et $h=0,01$.
    \[
    f(0,01) \approx f(0) + 0,01 \times f'(0) = 2 + 0,01 \times 0,25 = 2 + 0,0025 = 2,0025.
    \]
    \textbf{Conclusion :} $\sqrt{4,01} \approx 2,0025$ (valeur exacte $\approx 2,002498\dots$).
\end{enumerate}
\end{exoresolu}

\begin{savaistu}[Histoire des mathématiques : La querelle Newton-Leibniz]
La notion de dérivée naît au XVII\textsuperscript{e} siècle pour résoudre le problème de la tangente et celui du mouvement instantané. \textbf{Isaac Newton} (1643--1727) utilise les « fluxions » (vitesses) pour la physique. \textbf{Gottfried Wilhelm Leibniz} (1646--1716) invente la notation $\frac{dy}{dx}$ et le $d$ différentiel, encore utilisée aujourd'hui. Une violente querelle de priorité opposa les deux savants ; l'histoire retient que les deux ont découvert le calcul infinitésimal indépendamment.
\end{savaistu}

\courssub{Dérivabilité à gauche et à droite}

\begin{definition}[Nombres dérivés à gauche et à droite]
Soit $f$ définie sur un intervalle $I$ et $a$ une borne de $I$ ou un point intérieur.
\begin{itemize}
    \item Le \textbf{nombre dérivé à gauche} en $a$ est $f'_g(a) = \lim_{h \to 0^-} \frac{f(a+h)-f(a)}{h}$ (si la limite existe finie).
    \item Le \textbf{nombre dérivé à droite} en $a$ est $f'_d(a) = \lim_{h \to 0^+} \frac{f(a+h)-f(a)}{h}$ (si la limite existe finie).
\end{itemize}
\end{definition}

\begin{propriete}[Critère de dérivabilité en un point]
Une fonction $f$ est dérivable en $a$ \textbf{si et seulement si} $f'_g(a)$ et $f'_d(a)$ existent et sont égaux. On a alors $f'(a) = f'_g(a) = f'_d(a)$.
\par
Cas de non-dérivabilité :
\begin{itemize}
    \item $f'_g(a) \neq f'_d(a)$ (finis) : \textbf{Point anguleux} (deux demi-tangentes distinctes).
    \item $f'_g(a) = +\infty$ et $f'_d(a) = -\infty$ (ou inversement) : \textbf{Point de rebroussement (cuspide) de première espèce} (tangente verticale commune).
    \item $f'_g(a) = f'_d(a) = +\infty$ (ou $-\infty$) : \textbf{Tangente verticale} (pas de nombre dérivé fini).
\end{itemize}
\end{propriete}

\begin{methode}[Étudier la dérivabilité à gauche et à droite — Points de non-dérivabilité]
\textbf{Objectif :} Déterminer si une fonction est dérivable en un point « suspect » (changement de définition, valeur absolue, racine en 0, partie entière).
\begin{enumerate}
    \item \textbf{Calculer le nombre dérivé à gauche :} $f'_g(a) = \lim_{h \to 0^-} \frac{f(a+h)-f(a)}{h}$.
    \item \textbf{Calculer le nombre dérivé à droite :} $f'_d(a) = \lim_{h \to 0^+} \frac{f(a+h)-f(a)}{h}$.
    \item \textbf{Comparer :}
        \begin{itemize}
            \item Si $f'_g(a) = f'_d(a) \in \mathbb{R}$ : $f$ est dérivable en $a$, $f'(a)$ = cette valeur.
            \item Si $f'_g(a) \neq f'_d(a)$ (finis) : \textbf{Point anguleux}. Tangentes distinctes à gauche et à droite.
            \item Si l'une des limites est $+\infty$ et l'autre $-\infty$ : \textbf{Cuspide}. Tangente verticale commune.
            \item Si les deux limites sont $+\infty$ (ou $-\infty$) : \textbf{Tangente verticale}. Pas de nombre dérivé.
        \end{itemize}
    \item \textbf{Rédiger :} « $f$ n'est pas dérivable en $a$ car $f'_g(a) = \dots \neq f'_d(a) = \dots$ ».
\end{enumerate}
\textbf{Exemples types :} $f(x)=|x|$ en $0$ : $f'_g(0)=-1$, $f'_d(0)=1$ (angle). $f(x)=x^{2/3}$ en $0$ : limites $\pm\infty$ (cuspide). $f(x)=x^{1/3}$ en $0$ : limite $+\infty$ des deux côtés (tangente verticale $x=0$).
\end{methode}

\begin{exoresolu}[Dérivabilité à gauche/droite : Valeur absolue et racine cubique — Type Probatoire]
Soient les fonctions $f(x) = |x^2 - 4|$ et $g(x) = \sqrt[3]{x} = x^{1/3}$.
\begin{enumerate}
    \item Étudier la dérivabilité de $f$ en $2$ et en $-2$.
    \item Étudier la dérivabilité de $g$ en $0$. Préciser la nature du point et l'équation de la tangente si elle existe.
\end{enumerate}
\tcblower
\textbf{Solution détaillée}
\begin{enumerate}
    \item $f(x) = |(x-2)(x+2)|$.
    \textbf{En $2$ :} $f(2)=0$.
    Pour $x \in ]0;2[$, $x^2-4 < 0 \implies f(x) = -(x^2-4) = 4-x^2$. $f'(x) = -2x$. $f'_g(2) = \lim_{x\to 2^-} -2x = -4$.
    Pour $x \in ]2; +\infty[$, $x^2-4 > 0 \implies f(x) = x^2-4$. $f'(x) = 2x$. $f'_d(2) = \lim_{x\to 2^+} 2x = 4$.
    $f'_g(2) = -4 \neq f'_d(2) = 4$. \textbf{$f$ n'est pas dérivable en $2$ (point anguleux).}
    Tangente gauche : $y = -4(x-2)$. Tangente droite : $y = 4(x-2)$.
    
    \textbf{En $-2$ :} Par parité ($f$ paire), $f'_g(-2) = 4$, $f'_d(-2) = -4$. \textbf{Non dérivable en $-2$ (point anguleux symétrique).}
    
    \item $g(x) = x^{1/3}$ définie sur $\mathbb{R}$. $g(0)=0$.
    Pour $h \neq 0$, $\tau(h) = \frac{h^{1/3}}{h} = h^{-2/3} = \frac{1}{h^{2/3}} = \frac{1}{\sqrt[3]{h^2}}$.
    $\lim_{h \to 0^-} \tau(h) = +\infty$ et $\lim_{h \to 0^+} \tau(h) = +\infty$.
    Les deux nombres dérivés sont infinis et de même signe.
    \textbf{$g$ n'est pas dérivable en $0$ (pas de nombre dérivé fini), mais admet une tangente verticale d'équation $\mathbf{x=0}$.}
    C'est un point à tangente verticale (pas un cuspide car les deux branches montent).
\end{enumerate}
\end{exoresolu}

\begin{popfigure}
\begin{tikzpicture}[scale=1.2, >=Stealth]
% Axes
\draw[->] (-3,0) -- (3,0) node[right] {$x$};
\draw[->] (0,-2) -- (0,2.5) node[above] {$y$};
% f(x) = |x^2-4| (parabola branches)
\draw[popblue, thick, domain=-2.8:-2, samples=50] plot (\x, {(\x*\x-4)});
\draw[popblue, thick, domain=-2:2, samples=50] plot (\x, {4-\x*\x});
\draw[popblue, thick, domain=2:2.8, samples=50] plot (\x, {(\x*\x-4)});
% Tangents at x=2
\draw[poporange, dashed, domain=0.5:3.5] plot (\x, {4*(\x-2)});
\draw[poporange, dashed, domain=0.5:3.5] plot (\x, {-4*(\x-2)});
% Points
\fill[popblue] (2,0) circle (2pt) node[below right] {$A(2,0)$};
\fill[popblue] (-2,0) circle (2pt) node[below left] {$B(-2,0)$};
% g(x) = x^(1/3)
\draw[popgreen, thick, domain=-2.5:2.5, samples=100] plot (\x, {sign(\x)*abs(\x)^(1/3)});
% Vertical tangent at 0
\draw[popgreen, dashed] (0,-2) -- (0,2.5);
\fill[popgreen] (0,0) circle (2pt) node[above right] {$O$};
\end{tikzpicture}
\legende{À gauche : $f(x)=|x^2-4|$ avec point anguleux en $x=2$ (demi-tangentes de pentes $-4$ et $4$). À droite : $g(x)=\sqrt[3]{x}$ avec tangente verticale $x=0$ en $x=0$.}
\end{popfigure}

\courssub{Fonction dérivée sur un intervalle}

\begin{definition}[Fonction dérivée sur un intervalle]
Soit $f$ une fonction définie sur un intervalle $I$. Si $f$ est dérivable en tout point de $I$, on appelle \textbf{fonction dérivée de $f$ sur $I$} la fonction $f'$ qui associe à tout $x \in I$ le nombre dérivé $f'(x)$.
\end{definition}

\begin{propriete}[Règles de dérivation usuelles — À connaître par cœur pour le Probatoire]
Soient $u$ et $v$ des fonctions dérivables sur un intervalle $I$, $n \in \mathbb{N}^*$, $a \in \mathbb{R}$.
\begin{center}
\begin{tabular}{|l|c|}\hline
\textbf{Fonction $f(x)$} & \textbf{Dérivée $f'(x)$} \\ \hline
$k$ (constante) & $0$ \\ \hline
$x^n$ & $n x^{n-1}$ \\ \hline
$\sqrt{u(x)}$ ($u(x)>0$) & $\dfrac{u'(x)}{2\sqrt{u(x)}}$ \\ \hline
$\dfrac{1}{u(x)}$ ($u(x)\neq 0$) & $-\dfrac{u'(x)}{u(x)^2}$ \\ \hline
$u(x) + v(x)$ & $u'(x) + v'(x)$ \\ \hline
$u(x) v(x)$ & $u'(x)v(x) + u(x)v'(x)$ \\ \hline
$\dfrac{u(x)}{v(x)}$ ($v(x)\neq 0$) & $\dfrac{u'(x)v(x) - u(x)v'(x)}{v(x)^2}$ \\ \hline
$\ln u(x)$ ($u(x)>0$) & $\dfrac{u'(x)}{u(x)}$ \\ \hline
$e^{u(x)}$ & $u'(x) e^{u(x)}$ \\ \hline
$\sin u(x)$ & $u'(x) \cos u(x)$ \\ \hline
$\cos u(x)$ & $-u'(x) \sin u(x)$ \\ \hline
\end{tabular}
\end{center}
\end{propriete}

\begin{methode}[Déterminer la fonction dérivée sur un intervalle]
\textbf{Objectif :} Trouver l'expression de $f'(x)$ sur un intervalle $I$ où $f$ est dérivable.
\begin{enumerate}
    \item \textbf{Identifier} la forme de $f$ : somme, produit, quotient, composée $u \circ v$, puissance $u^n$, racine, fonctions usuelles ($\exp, \ln, \sin, \cos$).
    \item \textbf{Appliquer} les règles de dérivation (linéarité, $(uv)'=u'v+uv'$, $\left(\frac{u}{v}\right)'=\frac{u'v-uv'}{v^2}$, $(u^n)'=n u' u^{n-1}$, $(\ln u)'=\frac{u'}{u}$, $(e^u)'=u'e^u$, etc.).
    \item \textbf{Simplifier} l'expression obtenue (factorisation, réduction de fraction).
    \item \textbf{Préciser} l'ensemble de définition de $f'$ (souvent $D_f$ privé des points où $f$ n'est pas dérivable : $v=0$ pour quotient, $u \le 0$ pour $\ln u$, etc.).
\end{enumerate}
\end{methode}

\begin{exoresolu}[Dérivée d'un quotient et d'une composée — Type Probatoire]
Soit $f(x) = \frac{x^2-1}{x^2+1}$ définie sur $\mathbb{R}$.
\begin{enumerate}
    \item Déterminer l'expression de $f'(x)$ sur $\mathbb{R}$.
    \item Résoudre l'équation $f'(x) = 0$ et dresser le tableau de signes de $f'$.
    \item En déduire les variations de $f$ et son tableau de variations.
\end{enumerate}
\tcblower
\textbf{Solution détaillée}
\begin{enumerate}
    \item $f$ est un quotient de fonctions polynomiales (dérivables sur $\mathbb{R}$) avec dénominateur $x^2+1 \neq 0$ sur $\mathbb{R}$. $f$ est dérivable sur $\mathbb{R}$.
    Posons $u(x)=x^2-1$ et $v(x)=x^2+1$. $u'(x)=2x$, $v'(x)=2x$.
    \[
    f'(x) = \frac{u'v - uv'}{v^2} = \frac{2x(x^2+1) - (x^2-1)2x}{(x^2+1)^2}.
    \]
    Simplifions le numérateur :
    \begin{align*}
    \text{Num} &= 2x^3 + 2x - (2x^3 - 2x) \\
    &= 2x^3 + 2x - 2x^3 + 2x = 4x.
    \end{align*}
    Donc \textbf{$f'(x) = \frac{4x}{(x^2+1)^2}$} sur $\mathbb{R}$.
    
    \item $f'(x) = 0 \iff 4x = 0 \iff x = 0$.
    Le dénominateur $(x^2+1)^2 > 0$ pour tout $x \in \mathbb{R}$.
    Le signe de $f'$ est le signe de $4x$ (donc de $x$).
    \begin{center}
    \begin{tabular}{|c|ccccc|}\hline
    $x$ & $-\infty$ & & $0$ & & $+\infty$ \\ \hline
    $f'(x)$ & & $-$ & $0$ & $+$ & \\ \hline
    $f$ & $1$ & $\searrow$ & $-1$ & $\nearrow$ & $1$ \\ \hline
    \end{tabular}
    \end{center}
    \textit{Note :} $\lim_{x\to\pm\infty} f(x) = 1$ (asymptote horizontale $y=1$). $f(0)=-1$.
    
    \item \textbf{Variations :} $f$ est décroissante sur $]-\infty; 0]$, minimale en $0$ avec $f(0)=-1$, croissante sur $[0; +\infty[$.
\end{enumerate}
\end{exoresolu}

\courssub{Opérations sur les fonctions dérivables — Composée d'une fonction affine}

\begin{propriete}[Dérivée de la composée d'une fonction affine]
Soit $f$ dérivable sur un intervalle $I$ et $u(x) = ax+b$ une fonction affine ($a,b \in \mathbb{R}, a \neq 0$) telle que $u(x) \in I$.
La fonction $g(x) = f(ax+b)$ est dérivable et \textbf{$g'(x) = a \cdot f'(ax+b)$}.
\par
\textit{Preuve :} $g'(x) = u'(x) f'(u(x)) = a \cdot f'(ax+b)$ (Règle de la chaîne cas particulier).
\end{propriete}

\begin{methode}[Dériver une fonction composée avec une fonction affine — Règle de la chaîne cas particulier]
\textbf{Objectif :} Calculer la dérivée de $g(x) = f(ax+b)$ où $a,b \in \mathbb{R}$, $a \neq 0$ et $f$ dérivable.
\begin{enumerate}
    \item \textbf{Identifier} la fonction intérieure $u(x) = ax+b$ (affine) et la fonction extérieure $f$.
    \item \textbf{Appliquer} la formule : $g'(x) = a \times f'(ax+b)$.
    \item \textbf{Attention :} Ne pas oublier le facteur $a$ (coefficient directeur de l'affine). C'est l'erreur n°1 sur ce chapitre.
    \item \textbf{Simplifier} l'expression finale.
\end{enumerate}
\textbf{Exemples :} Si $g(x) = \sin(3x-2)$, alors $g'(x) = 3 \cos(3x-2)$. Si $g(x) = (2x+1)^5$, $g'(x) = 5 \times 2 \times (2x+1)^4 = 10(2x+1)^4$.
\end{methode}

\begin{exoresolu}[Dérivée de composée affine et recherche d'antécédents — Type Probatoire]
Soit $f(x) = \ln(2x+3)$ définie sur $]-\frac{3}{2}; +\infty[$.
\begin{enumerate}
    \item Calculer $f'(x)$.
    \item Résoudre $f'(x) = 1$.
    \item Déterminer l'équation de la tangente $\mathcal{T}$ à la courbe $\mathcal{C}_f$ au point d'abscisse $0$.
\end{enumerate}
\tcblower
\textbf{Solution détaillée}
\begin{enumerate}
    \item $f$ est de la forme $\ln(u(x))$ avec $u(x)=2x+3$ (affine, $u'(x)=2$). $u(x)>0$ sur le domaine.
    Formule : $(\ln u)' = \frac{u'}{u}$.
    \textbf{$f'(x) = \frac{2}{2x+3}$} sur $]-\frac{3}{2}; +\infty[$.
    
    \item Équation $f'(x) = 1 \iff \frac{2}{2x+3} = 1$.
    $2x+3 > 0$ sur le domaine, on multiplie : $2 = 2x+3 \iff 2x = -1 \iff x = -0,5$.
    Vérification : $-0,5 > -1,5$, la solution est acceptée.
    \textbf{L'unique solution est $x = -0,5$.}
    
    \item Point d'abscisse $0$ : $f(0) = \ln(3)$. $f'(0) = \frac{2}{3}$.
    Équation de la tangente $\mathcal{T}$ : $y = f'(0)(x-0) + f(0)$.
    $y = \frac{2}{3}x + \ln 3$.
    \textbf{Équation de $\mathcal{T}$ : $y = \frac{2}{3}x + \ln 3$.}
\end{enumerate}
\end{exoresolu}

\courssub{Tangente et approximation affine}

\begin{definition}[Tangente à une courbe en un point]
Soit $f$ dérivable en $x_0$. La \textbf{tangente} à la courbe $\mathcal{C}_f$ au point $A(x_0; f(x_0))$ est la droite $\mathcal{T}$ d'équation :
\[
y = f'(x_0)(x - x_0) + f(x_0).
\]
C'est la droite qui « colle » le mieux à la courbe au voisinage de $A$.
\end{definition}

\begin{propriete}[Approximation affine (ou linéaire)]
Si $f$ est dérivable en $x_0$, pour $h$ « petit » (proche de $0$), on a l'approximation :
\[
f(x_0+h) \approx f(x_0) + h f'(x_0).
\]
Géométriquement, on remplace la courbe par sa tangente. L'erreur commise est l'écart entre la courbe et la tangente.
\par
Si $f$ est \textbf{convexe} sur un intervalle contenant $x_0$ (courbe au-dessus de la tangente), alors $f(x_0+h) \ge f(x_0) + h f'(x_0)$ (minorant).
Si $f$ est \textbf{concave} (courbe en-dessous de la tangente), alors $f(x_0+h) \le f(x_0) + h f'(x_0)$ (majorant).
\end{propriete}

\begin{methode}[Équation de la tangente et approximation affine au voisinage d'un point]
\textbf{Objectif :} Écrire l'équation de la tangente en un point $A(x_0; f(x_0))$ et utiliser l'approximation affine $f(x_0+h) \approx f(x_0) + h f'(x_0)$.
\begin{enumerate}
    \item \textbf{Vérifier} que $f$ est dérivable en $x_0$ (calculer $f'(x_0)$).
    \item \textbf{Équation de la tangente $\mathcal{T}$ :} $y = f'(x_0)(x - x_0) + f(x_0)$.
    \item \textbf{Approximation affine :} Pour $h$ « petit » (ex: $0,01$, $0,1$), $f(x_0+h) \approx f(x_0) + h f'(x_0)$.
    \item \textbf{Interprétation géométrique :} On remplace la courbe par sa tangente.
    \item \textbf{Encadrement (si demandé) :} Étudier la convexité (signe de $f''$ si connu, ou forme de $f'$) pour dire si l'approximation est un majorant ou un minorant.
\end{enumerate}
\textbf{Rédiger :} « Par approximation affine au voisinage de $x_0$, on a $f(x_0+h) \approx \dots$ »
\end{methode}

\begin{experience}[Zoom sur une courbe : découverte de la tangente]
\textbf{Activité :} Dans un logiciel de géométrie dynamique (GeoGebra, Python+Matplotlib), tracer $f(x)=x^2$ et zoomer successivement autour du point $A(1;1)$ avec des facteurs $10, 100, 1000$.
\par
\textbf{Observation :} La courbe devient indiscernable de la droite tangente $y=2x-1$.
\par
\textbf{Conclusion :} La dérivée $f'(1)=2$ est bien le coefficient directeur de la meilleure approximation affine locale.
\end{experience}

\begin{exoresolu}[Tangente, approximation et encadrement par convexité — Type Probatoire]
Soit $f(x) = x^3 - 3x + 1$ sur $\mathbb{R}$. On admet que $f$ est convexe sur $[0; +\infty[$.
\begin{enumerate}
    \item Déterminer l'équation de la tangente $\mathcal{T}$ en $x_0=1$.
    \item Donner une valeur approchée de $f(1,02)$ par approximation affine.
    \item Préciser si cette valeur est un majorant ou un minorant de $f(1,02)$.
\end{enumerate}
\tcblower
\textbf{Solution détaillée}
\begin{enumerate}
    \item $f'(x) = 3x^2 - 3$. $f'(1) = 0$. $f(1) = 1 - 3 + 1 = -1$.
    La tangente en $1$ a pour pente $0$ (tangente horizontale).
    Équation : $y = 0 \times (x-1) - 1 \iff \textbf{y = -1}$.
    
    \item Approximation affine en $x_0=1$ pour $h=0,02$ :
    $f(1,02) \approx f(1) + 0,02 \times f'(1) = -1 + 0,02 \times 0 = -1$.
    \textbf{Valeur approchée : $-1$.}
    
    \item $f$ est convexe sur $[0; +\infty[$. La courbe $\mathcal{C}_f$ est \textbf{au-dessus} de sa tangente sur cet intervalle.
    Donc pour $x=1,02 \in [0; +\infty[$, $f(1,02) \ge -1$.
    La valeur approchée $-1$ est un \textbf{minorant} de $f(1,02)$.
    \textit{(Vérification : $f(1,02) = 1,061208 - 3,06 + 1 = -0,998792 > -1$).}
\end{enumerate}
\end{exoresolu}

\courssub{Dérivée et sens de variation}

\begin{propriete}[Dérivée et sens de variation — Théorème fondamental]
Soit $f$ une fonction dérivable sur un intervalle $I$.
\begin{itemize}
    \item Si $f'(x) \ge 0$ sur $I$, alors $f$ est \textbf{croissante} sur $I$.
    \item Si $f'(x) \le 0$ sur $I$, alors $f$ est \textbf{décroissante} sur $I$.
    \item Si $f'(x) = 0$ sur $I$, alors $f$ est \textbf{constante} sur $I$.
\end{itemize}
\par
Réciproque (utile pour les extrema) : Si $f$ a un extremum local en $x_0 \in I$ (intérieur à $I$) et $f$ dérivable en $x_0$, alors $f'(x_0)=0$.
\par
\textbf{Attention :} $f'(x_0)=0$ ne suffit pas pour garantir un extremum (ex: $f(x)=x^3$ en $0$, point d'inflexion à tangente horizontale). Il faut vérifier le \textbf{changement de signe} de $f'$.
\end{propriete}

\begin{methode}[Étudier les variations d'une fonction à l'aide de sa dérivée — Tableau de variations]
\textbf{Objectif :} Construire le tableau de variations complet d'une fonction dérivable sur un intervalle.
\begin{enumerate}
    \item \textbf{Calculer} $f'(x)$ et simplifier (factoriser si possible).
    \item \textbf{Étudier le signe de $f'(x)$} : résoudre $f'(x)=0$, dresser un tableau de signes (valeurs interdites = bornes du domaine ou zéros du dénominateur).
    \item \textbf{Construire le tableau de variations} :
        \begin{itemize}
            \item Ligne $x$ : valeurs interdites (double barre ||) et zéros de $f'$.
            \item Ligne $f'$ : signes ($+$, $-$, $0$).
            \item Ligne $f$ : flèches $\nearrow$ (si $f'>0$), $\searrow$ (si $f'<0$), valeurs $f(x)$ aux points remarquables, limites aux bornes du domaine.
        \end{itemize}
    \item \textbf{Conclure} sur les extrema (min/max) et le sens de variation sur chaque intervalle.
\end{enumerate}
\textbf{Piège :} Ne pas oublier les limites aux bornes du domaine (asymptotes) et les valeurs interdites (double barre dans le tableau).
\end{methode}

\begin{exoresolu}[Étude complète : variations, extrema, asymptotes — Type Probatoire]
Soit $f(x) = x + \frac{4}{x}$ définie sur $D_f = \mathbb{R}^* = ]-\infty;0[ \cup ]0;+\infty[$.
\begin{enumerate}
    \item Calculer $f'(x)$ et étudier son signe.
    \item Dresser le tableau de variations de $f$ sur $D_f$.
    \item Montrer que la droite $\mathcal{D}$ d'équation $y=x$ est une asymptote.
    \item Tracer $\mathcal{C}_f$ et $\mathcal{D}$ dans un repère (qualitatif).
\end{enumerate}
\tcblower
\textbf{Solution détaillée}
\begin{enumerate}
    \item $f'(x) = 1 - \frac{4}{x^2} = \frac{x^2-4}{x^2} = \frac{(x-2)(x+2)}{x^2}$.
    $x^2 > 0$ sur $D_f$. Signe de $f'$ = signe de $(x-2)(x+2)$.
    Zéros : $x=-2$ et $x=2$. Valeur interdite : $x=0$.
    
    \item \textbf{Tableau de variations :}
    \begin{center}
    \begin{tabular}{|c|cccccccccc|}\hline
    $x$ & $-\infty$ & & $-2$ & & $0$ & & $2$ & & $+\infty$ \\ \hline
    $f'(x)$ & & $+$ & $0$ & $+$ & $||$ & $-$ & $0$ & $+$ & \\ \hline
    $f(x)$ & $-\infty$ & $\nearrow$ & $-4$ & $\nearrow$ & $||$ & $+\infty$ & $\searrow$ & $4$ & $\nearrow$ & $+\infty$ \\ \hline
    \end{tabular}
    \end{center}
    \textit{Calculs :} $f(-2)=-2-2=-4$ (max local). $f(2)=2+2=4$ (min local).
    Limites : $\lim_{x\to 0^-} f(x) = -\infty$, $\lim_{x\to 0^+} f(x) = +\infty$. $\lim_{x\to\pm\infty} f(x) = \pm\infty$.
    
    \item Asymptote : $f(x) - x = \frac{4}{x}$. $\lim_{x\to\pm\infty} (f(x)-x) = \lim_{x\to\pm\infty} \frac{4}{x} = 0$.
    Donc $y=x$ est asymptote (oblique) aux deux infinis. Position relative : $f(x)-x = \frac{4}{x}$. Pour $x>0$, $f(x)>x$ (courbe au-dessus). Pour $x<0$, $f(x)<x$ (courbe en-dessous).
    
    \item \textbf{Croquis :} Deux branches hyperboliques. Branche $x>0$ : part de $+\infty$ (axe $y$), descend jusqu'au min $4$ en $2$, remonte vers $+\infty$ en suivant $y=x$ par au-dessus. Branche $x<0$ : part de $-\infty$, monte jusqu'au max $-4$ en $-2$, redescend vers $-\infty$ en suivant $y=x$ par en-dessous.
\end{enumerate}
\end{exoresolu}

\begin{popfigure}
\begin{tikzpicture}[scale=0.7, >=Stealth]
% Axes
\draw[->] (-6,0) -- (6,0) node[right] {$x$};
\draw[->] (0,-8) -- (0,8) node[above] {$y$};
% Asymptote y=x
\draw[popmuted, dashed, thick] (-6,-6) -- (6,6) node[above right] {$y=x$};
% Branch x>0
\draw[popblue, thick, domain=0.3:5.5, samples=100] plot (\x, {\x + 4/\x});
% Branch x<0
\draw[popblue, thick, domain=-5.5:-0.3, samples=100] plot (\x, {\x + 4/\x});
% Min/Max points
\fill[popblue] (2,4) circle (3pt) node[above right] {$(2,4)$ min};
\fill[popblue] (-2,-4) circle (3pt) node[below left] {$(-2,-4)$ max};
% Vertical asymptote x=0
\draw[popmuted, dashed] (0,-8) -- (0,8);
\end{tikzpicture}
\legende{Courbe $\mathcal{C}_f$ de $f(x)=x+\frac{4}{x}$ et son asymptote oblique $y=x$.}
\end{popfigure}

\courssub{Applications à des situations concrètes}

\begin{methode}[Résoudre un problème concret d'optimisation ou de taux de variation instantané]
\textbf{Objectif :} Modéliser une situation technique (électricité, mécanique, économie, géométrie) par une fonction $f$, puis utiliser la dérivée pour optimiser ou calculer un taux.
\begin{enumerate}
    \item \textbf{Lire} l'énoncé, identifier la grandeur à optimiser (coût, surface, puissance, volume) ou le taux à calculer (vitesse, intensité, variation de température).
    \item \textbf{Introduire} une variable $x$ (longueur, temps, quantité) et exprimer la grandeur $y = f(x)$.
    \item \textbf{Préciser} l'intervalle $I$ de variation de $x$ (contraintes physiques : $x>0$, $x < L$, etc.).
    \item \textbf{Étudier} les variations de $f$ sur $I$ (tableau de variations).
    \item \textbf{Conclure} : « Le maximum/minimum est atteint pour $x = \dots$, la valeur optimale est $\dots$ » ou « Le taux de variation instantané à l'instant $t_0$ est $f'(t_0) = \dots$ unité/unité ».
    \item \textbf{Vérifier} la cohérence du résultat (ordre de grandeur, unité).
\end{enumerate}
\textbf{Exemples classiques :} Puissance maximale dans un circuit ($P=RI^2$), volume maximal d'une boîte, coût minimal de production, vitesse d'une réaction chimique.
\end{methode}

\begin{exoresolu}[Optimisation géométrique : Cylindre inscrit dans un cône — Type Probatoire]
Un cône droit a pour hauteur $H=12$ cm et pour rayon de base $R=6$ cm. On inscrit un cylindre droit de hauteur $h$ et de rayon $r$ dans ce cône (une base du cylindre sur la base du cône, l'autre base touchant la génératrice).
\begin{enumerate}
    \item Exprimer $r$ en fonction de $h$ (théorème de Thalès dans une section axiale).
    \item Exprimer le volume $V(h)$ du cylindre en fonction de $h$. Préciser l'intervalle de variation de $h$.
    \item Déterminer les dimensions $r$ et $h$ du cylindre de volume maximal.
\end{enumerate}
\tcblower
\textbf{Solution détaillée}
\begin{enumerate}
    \item Section axiale : triangle rectangle (hauteur $H$, base $R$) et rectangle inscrit (hauteur $h$, largeur $2r$). Par similarité des triangles (Thalès) :
    \[
    \frac{r}{R} = \frac{H-h}{H} \implies r = R\left(1 - \frac{h}{H}\right) = 6\left(1 - \frac{h}{12}\right) = 6 - \frac{h}{2}.
    \]
    
    \item Volume cylindre : $V = \pi r^2 h$.
    $V(h) = \pi \left(6 - \frac{h}{2}\right)^2 h = \pi \left(36 - 6h + \frac{h^2}{4}\right)h = \pi \left(36h - 6h^2 + \frac{h^3}{4}\right)$.
    Contraintes : $h \ge 0$ et $r \ge 0 \implies 6 - h/2 \ge 0 \implies h \le 12$.
    Domaine d'étude : \textbf{$h \in [0; 12]$}.
    
    \item Dérivée sur $[0; 12]$ :
    $V'(h) = \pi \left(36 - 12h + \frac{3h^2}{4}\right) = \frac{3\pi}{4} \left(48 - 16h + h^2\right) = \frac{3\pi}{4} (h^2 - 16h + 48)$.
    Résolvons $h^2 - 16h + 48 = 0$. $\Delta = 256 - 192 = 64$. $\sqrt{\Delta}=8$.
    $h_1 = \frac{16-8}{2} = 4$ ; $h_2 = \frac{16+8}{2} = 12$.
    $V'(h) = \frac{3\pi}{4}(h-4)(h-12)$.
    Tableau de signes sur $[0;12]$ : $V' \ge 0$ sur $[0;4]$, $V' \le 0$ sur $[4;12]$.
    $V$ croît sur $[0;4]$, décroît sur $[4;12]$. Maximum en $h=4$.
    $h_{\text{opt}} = 4$ cm. $r_{\text{opt}} = 6 - \frac{4}{2} = 4$ cm.
    $V_{\text{max}} = V(4) = \pi \times 4^2 \times 4 = 64\pi \approx 201,06 \text{ cm}^3$.
    \textbf{Le cylindre de volume maximal a pour dimensions : hauteur $4$ cm et rayon $4$ cm.}
\end{enumerate}
\end{exoresolu}

\begin{exoresolu}[Application électricité : Puissance maximale transférée — Type Probatoire]
Un générateur réel de tension $E = 12$ V et de résistance interne $r = 2\ \Omega$ alimente un récepteur de résistance variable $R$ (en $\Omega$). La puissance dissipée par le récepteur est $P(R) = R \left(\frac{E}{R+r}\right)^2$.
\begin{enumerate}
    \item Exprimer $P(R)$ et préciser son domaine de définition physique.
    \item Étudier les variations de $P$ et déterminer la valeur de $R$ pour laquelle la puissance est maximale.
    \item Calculer cette puissance maximale.
\end{enumerate}
\tcblower
\textbf{Solution détaillée}
\begin{enumerate}
    \item $P(R) = \frac{E^2 R}{(R+r)^2} = \frac{144 R}{(R+2)^2}$. Domaine physique : $R > 0$ (résistance strictement positive). On étudie sur $]0; +\infty[$.
    
    \item Dérivée (quotient) : $u(R)=144R$, $v(R)=(R+2)^2$. $u'=144$, $v'=2(R+2)$.
    \begin{align*}
    P'(R) &= \frac{144(R+2)^2 - 144R \cdot 2(R+2)}{(R+2)^4} \\
    &= \frac{144(R+2)\big[(R+2) - 2R\big]}{(R+2)^4} \\
    &= \frac{144(2 - R)}{(R+2)^3}.
    \end{align*}
    Le dénominateur $>0$ sur $]0;+\infty[$. Signe de $P'$ = signe de $2-R$.
    $P'(R) = 0 \iff R = 2\ \Omega$.
    \begin{center}
    \begin{tabular}{|c|ccccc|}\hline
    $R$ & $0$ & & $2$ & & $+\infty$ \\ \hline
    $P'(R)$ & & $+$ & $0$ & $-$ & \\ \hline
    $P(R)$ & $0$ & $\nearrow$ & $P_{\text{max}}$ & $\searrow$ & $0$ \\ \hline
    \end{tabular}
    \end{center}
    $P$ croît sur $]0;2]$, décroît sur $[2;+\infty[$. Maximum pour \textbf{$R = 2\ \Omega$} (égale à la résistance interne $r$).
    
    \item $P_{\text{max}} = P(2) = \frac{144 \times 2}{(4)^2} = \frac{288}{16} = 18$ W.
    \textbf{Puissance maximale : $18$ W pour $R = 2\ \Omega$.}
\end{enumerate}
\end{exoresolu}

\begin{attention}[Les 5 erreurs qui coûtent des points au Probatoire]
\begin{enumerate}
    \item \textbf{Oublier le coefficient directeur dans la dérivée d'une composée affine.} Ex: dériver $\ln(3x+1)$ et écrire $\frac{1}{3x+1}$ au lieu de $\frac{3}{3x+1}$. \textit{Règle : $(f(ax+b))' = a f'(ax+b)$.}
    \item \textbf{Confondre « dérivée nulle » et « extremum ».} $f'(x_0)=0$ est nécessaire mais pas suffisant. Il faut vérifier le \textbf{changement de signe} de $f'$ (ou signe de $f''$ si au programme) pour conclure à un max/min. Si $f'$ ne change pas de signe, c'est un point d'inflexion à tangente horizontale (ex: $x^3$ en $0$).
    \item \textbf{Négliger l'ensemble de définition (Domaine) avant de dériver.} Dériver $\ln(x-2)$ sans dire $x>2$, ou $\sqrt{x}$ sans $x\ge 0$, ou un quotient sans annuler le dénominateur. Le tableau de variations doit comporter les \textbf{doubles barres ||} pour les valeurs interdites.
    \item \textbf{Écrire l'équation de la tangente sous la forme $y = ax+b$ sans calculer $b$ correctement.} Formule à appliquer mécaniquement : $y = f'(x_0)(x-x_0) + f(x_0)$. Ne pas confondre $f'(x_0)$ (nombre) et $f'(x)$ (fonction).
    \item \textbf{Mauvaise rédaction de l'approximation affine.} Dire « $f(1,01) = f(1) + 0,01 f'(1)$ » (égalité fausse). Il faut écrire \textbf{$f(1,01) \approx f(1) + 0,01 f'(1)$} ou « Par approximation affine, $f(1,01) \simeq \dots$ ». Préciser toujours le point $x_0$ et l'incrément $h$.
\end{enumerate}
\end{attention}

\newpage

\coursec{Exercices}

\courssub{Je vérifie mes connaissances}

\begin{exercice}[QCM : nombre dérivé]
Pour chaque question, une seule réponse est exacte. Recopie la bonne lettre.
\begin{enumerate}
\item Soit $f(x)=x^2+3x$. Le nombre dérivé de $f$ en $1$ est :
\begin{enumerate}
\item[a)] $4$ \quad b)] $5$ \quad c)] $2$ \quad d)] $3$
\end{enumerate}
\item La fonction dérivée de $f(x)=\dfrac{1}{x}$ sur $]0;+\infty[$ est :
\begin{enumerate}
\item[a)] $f'(x)=\dfrac{1}{x^2}$ \quad b)] $f'(x)=-\dfrac{1}{x^2}$ \quad c)] $f'(x)=\ln x$ \quad d)] $f'(x)=-\dfrac{2}{x^3}$
\end{enumerate}
\item Si $f(x)=\sqrt{x}$, alors $f'(4)$ vaut :
\begin{enumerate}
\item[a)] $\dfrac{1}{2}$ \quad b)] $\dfrac{1}{4}$ \quad c)] $2$ \quad d)] $\dfrac{1}{8}$
\end{enumerate}
\item La tangente à la courbe de $f$ au point d'abscisse $a$ a pour équation :
\begin{enumerate}
\item[a)] $y=f(a)+f'(a)x$ \quad b)] $y=f'(a)(x-a)+f(a)$ \quad c)] $y=f(a)(x-a)+f'(a)$ \quad d)] $y=f'(a)x$
\end{enumerate}
\end{enumerate}
\end{exercice}

\begin{exercice}[Vrai ou faux, en justifiant]
Réponds par VRAI ou FAUX et justifie ta réponse par une démonstration ou un contre-exemple.
\begin{enumerate}
\item Toute fonction continue en un réel $a$ est dérivable en $a$.
\item La fonction $x\mapsto |x|$ est dérivable en $0$.
\item Si $f'(x)>0$ sur un intervalle $I$, alors $f$ est strictement croissante sur $I$.
\item Si $f$ et $g$ sont dérivables, alors $(f\times g)'=f'\times g'$.
\end{enumerate}
\end{exercice}

\begin{exercice}[Définitions et cours]
\begin{enumerate}
\item Donner la définition du nombre dérivé d'une fonction $f$ en un réel $a$ à l'aide d'une limite.
\item Qu'appelle-t-on \cle{approximation affine} de $f$ au voisinage de $a$ ? Écrire la formule correspondante.
\item Énoncer la condition sur $f'$ pour que $f$ soit décroissante sur un intervalle $I$.
\end{enumerate}
\end{exercice}

\begin{exercice}[Calculs directs de dérivées]
Calculer la fonction dérivée de chacune des fonctions suivantes, en précisant l'ensemble de dérivabilité :
\begin{enumerate}
\item $f(x)=5x^3-2x^2+7x-4$ ;
\item $g(x)=\dfrac{3}{x}+\sqrt{x}$ ;
\item $h(x)=(2x+1)^3$ ;
\item $k(x)=\sqrt{5x-10}$.
\end{enumerate}
\end{exercice}

\courssub{Je m'entraîne}

\begin{exercice}[Dérivabilité à gauche et à droite]
Soit $f$ la fonction définie sur $\mathbb{R}$ par :
\[ f(x)=\begin{cases} x^2 & \text{si } x\leq 1 \\ 2x-1 & \text{si } x>1 \end{cases} \]
\begin{enumerate}
\item Montrer que $f$ est continue en $1$.
\item Étudier la dérivabilité de $f$ à gauche en $1$, puis à droite en $1$.
\item $f$ est-elle dérivable en $1$ ? Interpréter graphiquement le résultat.
\end{enumerate}
\end{exercice}

\begin{exercice}[Opérations sur les dérivées]
Pour chacune des fonctions suivantes, préciser l'ensemble de définition, l'ensemble de dérivabilité, puis calculer la fonction dérivée :
\begin{enumerate}
\item $f(x)=(3x^2-5x)(x+2)$ ;
\item $g(x)=\dfrac{x-1}{2x+3}$ ;
\item $h(x)=\sqrt{4x-8}$.
\end{enumerate}
\end{exercice}

\begin{exercice}[Tangente et position relative]
Soit $f$ la fonction définie sur $\mathbb{R}$ par $f(x)=x^3-2x$, et $\mathcal{C}_f$ sa courbe représentative.
\begin{enumerate}
\item Calculer $f(-1)$ et $f'(-1)$.
\item Déterminer une équation de la tangente $T$ à $\mathcal{C}_f$ au point d'abscisse $-1$.
\item Étudier le signe de $f(x)-(x+2)$ pour $x$ au voisinage de $-1$, et en déduire la position relative de $\mathcal{C}_f$ et de $T$.
\end{enumerate}
\end{exercice}

\begin{exercice}[Approximation affine sans calculatrice]
\begin{enumerate}
\item À l'aide de l'approximation affine de $x\mapsto \sqrt{x}$ au voisinage de $16$, donner une valeur approchée de $\sqrt{16,1}$.
\item À l'aide de l'approximation affine de $x\mapsto \dfrac{1}{x^2}$ au voisinage de $1$, donner une valeur approchée de $\dfrac{1}{(1,02)^2}$.
\item Comparer ces valeurs approchées aux valeurs données par la calculatrice et commenter.
\end{enumerate}
\end{exercice}

\courssub{Je me prépare au Probatoire}

\begin{exercice}[Coût moyen dans une cimenterie]
Une entreprise de Douala fabrique $x$ sacs de ciment par jour ($x>0$). Le coût total de production, en milliers de FCFA, est donné par :
\[ C(x)=x^2-40x+2000. \]
Le coût moyen par sac est $CM(x)=\dfrac{C(x)}{x}$.
\begin{enumerate}
\item Montrer que $CM(x)=x-40+\dfrac{2000}{x}$.
\item Calculer $CM'(x)$ et montrer que $CM'(x)=\dfrac{x^2-2000}{x^2}$.
\item Étudier le signe de $CM'(x)$ sur $]0;+\infty[$ et dresser le tableau de variations de $CM$.
\item En déduire la production quotidienne qui minimise le coût moyen, ainsi que la valeur de ce coût moyen minimal (arrondie au franc près, en prenant $\sqrt{2000}\approx 44,72$).
\item Calculer le coût total correspondant à cette production.
\end{enumerate}
\end{exercice}

\begin{exercice}[Mouvement d'un mobile]
Un mobile se déplace sur un axe. Sa position, en mètres, à l'instant $t$ (en secondes, $t\geq 0$) est donnée par :
\[ x(t)=t^3-6t^2+9t. \]
\begin{enumerate}
\item Justifier que la vitesse du mobile à l'instant $t$ est $v(t)=x'(t)$, puis calculer $v(t)$.
\item Déterminer les instants où la vitesse s'annule.
\item Dresser le tableau de signes de $v(t)$ sur $[0;+\infty[$.
\item En déduire les intervalles de temps pendant lesquels le mobile avance (vitesse positive) et ceux où il recule (vitesse négative).
\item Calculer la position du mobile aux instants où sa vitesse s'annule, et interpréter.
\end{enumerate}
\end{exercice}

\begin{exercice}[Étude complète d'une fonction]
Soit $f$ la fonction définie sur $[0;4]$ par :
\[ f(x)=x^3-6x^2+9x+1, \]
et $\mathcal{C}_f$ sa courbe représentative dans un repère orthogonal.
\begin{enumerate}
\item Calculer $f'(x)$ et vérifier que $f'(x)=3(x-1)(x-3)$.
\item Étudier le signe de $f'(x)$ sur $[0;4]$.
\item Dresser le tableau de variations complet de $f$ sur $[0;4]$.
\item En déduire les extremums de $f$ sur $[0;4]$, en précisant les valeurs exactes.
\item Déterminer une équation de la tangente à $\mathcal{C}_f$ au point d'abscisse $0$.
\item Tracer $\mathcal{C}_f$ et cette tangente dans un repère orthogonal (unités : \SI{2}{cm} en abscisse, \SI{1}{cm} en ordonnée).
\end{enumerate}
\end{exercice}

\newpage

\coursec{Corrigés des exercices}

\coursec{Corrigés des exercices — Dérivation (Leçon 14)}

\begin{corrige}[Exercice 1 — QCM : nombre dérivé]
\begin{enumerate}
\item \textbf{Réponse b) } $5$. $f'(x)=2x+3$, donc $f'(1)=2\times 1+3=5$.
\item \textbf{Réponse b) } $f'(x)=-\dfrac{1}{x^2}$. C'est la dérivée de la fonction inverse (formulaire).
\item \textbf{Réponse b) } $\dfrac{1}{4}$. $f'(x)=\dfrac{1}{2\sqrt{x}}$, donc $f'(4)=\dfrac{1}{2\sqrt{4}}=\dfrac{1}{4}$.
\item \textbf{Réponse b) } $y=f'(a)(x-a)+f(a)$. C'est l'équation de la tangente au point d'abscisse $a$.
\end{enumerate}
\textbf{Points de vigilance :} ne pas confondre $f(a)$ (image) et $f'(a)$ (nombre dérivé) ; dans l'équation de la tangente, le coefficient directeur est $f'(a)$ et la tangente passe par le point $(a\,;\,f(a))$.
\end{corrige}

\begin{corrige}[Exercice 2 — Vrai ou faux, en justifiant]
\begin{enumerate}
\item \textbf{FAUX.} Contre-exemple : la fonction $x\mapsto |x|$ est continue en $0$ (car $\lim_{x\to 0}|x|=0=|0|$), mais elle n'est pas dérivable en $0$ (voir question 2). La continuité est une condition \emph{nécessaire} mais \emph{non suffisante} de dérivabilité.
\item \textbf{FAUX.} Pour $x>0$, $\dfrac{|x|-0}{x-0}=\dfrac{x}{x}=1\longrightarrow 1$ (dérivée à droite : $f'_d(0)=1$). Pour $x<0$, $\dfrac{|x|}{x}=\dfrac{-x}{x}=-1\longrightarrow -1$ (dérivée à gauche : $f'_g(0)=-1$). Comme $f'_g(0)\neq f'_d(0)$, $f$ n'est pas dérivable en $0$ (point anguleux).
\item \textbf{VRAI.} C'est le théorème du cours : si $f$ est dérivable sur un intervalle $I$ et si $f'(x)>0$ pour tout $x$ de $I$, alors $f$ est strictement croissante sur $I$.
\item \textbf{FAUX.} La bonne formule est $(f\times g)'=f'\times g+f\times g'$. Contre-exemple : $f(x)=g(x)=x$. Alors $(fg)(x)=x^2$, donc $(fg)'(x)=2x$, tandis que $f'(x)g'(x)=1\times 1=1$. Or $2x\neq 1$ en général.
\end{enumerate}
\textbf{Points de vigilance :} pour infirmer une propriété, un seul contre-exemple suffit ; pour la confirmer, il faut citer le théorème du cours avec ses hypothèses.
\end{corrige}

\begin{corrige}[Exercice 3 — Définitions et cours]
\begin{enumerate}
\item Soit $f$ une fonction définie sur un intervalle contenant $a$. On dit que $f$ est \cle{dérivable en $a$} si le taux d'accroissement $\dfrac{f(a+h)-f(a)}{h}$ admet une limite finie quand $h$ tend vers $0$. Cette limite est le \cle{nombre dérivé} de $f$ en $a$ :
\[ f'(a)=\lim_{h\to 0}\dfrac{f(a+h)-f(a)}{h}=\lim_{x\to a}\dfrac{f(x)-f(a)}{x-a}. \]
\item L'\cle{approximation affine} de $f$ au voisinage de $a$ est l'application affine $x\mapsto f(a)+f'(a)(x-a)$. Pour $x$ proche de $a$ (ou $h$ proche de $0$) :
\[ f(x)\approx f(a)+f'(a)(x-a) \qquad\text{ou}\qquad f(a+h)\approx f(a)+h\,f'(a). \]
\item \textbf{Critère de décroissance :} si $f$ est dérivable sur un intervalle $I$ et si $f'(x)\leq 0$ pour tout $x$ de $I$, alors $f$ est décroissante sur $I$ (strictement décroissante si $f'(x)<0$, sauf éventuellement en des points isolés).
\end{enumerate}
\end{corrige}

\begin{corrige}[Exercice 4 — Calculs directs de dérivées]
\begin{enumerate}
\item $f(x)=5x^3-2x^2+7x-4$ est une fonction polynôme, dérivable sur $\mathbb{R}$ :
\[ f'(x)=15x^2-4x+7. \]
\item $g(x)=\dfrac{3}{x}+\sqrt{x}=3\times\dfrac{1}{x}+\sqrt{x}$, dérivable sur $]0\,;\,+\infty[$ :
\[ g'(x)=-\dfrac{3}{x^2}+\dfrac{1}{2\sqrt{x}}. \]
\item $h(x)=(2x+1)^3$ : composée de $u(x)=2x+1$ (affine, $u'=2$) par $v\mapsto v^3$. $h$ est dérivable sur $\mathbb{R}$ et
\[ h'(x)=3\times 2\times (2x+1)^2=6(2x+1)^2. \]
\item $k(x)=\sqrt{5x-10}$ : définie si $5x-10\geq 0$, c'est-à-dire $x\geq 2$ ; dérivable si $5x-10>0$, c'est-à-dire sur $]2\,;\,+\infty[$. Avec $(\sqrt{u})'=\dfrac{u'}{2\sqrt{u}}$ :
\[ k'(x)=\dfrac{5}{2\sqrt{5x-10}}. \]
\end{enumerate}
\textbf{Points de vigilance :} la fonction racine carrée n'est jamais dérivable en la borne où son contenu s'annule ; pour $(ax+b)^n$ et $\sqrt{ax+b}$, ne pas oublier de multiplier par $a$ (dérivée de la fonction affine intérieure).
\end{corrige}

\begin{corrige}[Exercice 5 — Dérivabilité à gauche et à droite]
\begin{enumerate}
\item \textbf{Continuité en $1$.} $f(1)=1^2=1$. Limite à gauche : $\lim_{x\to 1^-}x^2=1$. Limite à droite : $\lim_{x\to 1^+}(2x-1)=1$. Les trois valeurs sont égales, donc $f$ est \textbf{continue en $1$}.
\item \textbf{Dérivabilité à gauche :} pour $x<1$,
\[ \dfrac{f(x)-f(1)}{x-1}=\dfrac{x^2-1}{x-1}=\dfrac{(x-1)(x+1)}{x-1}=x+1\xrightarrow[x\to 1^-]{}2. \]
Donc $f$ est dérivable à gauche en $1$ et $f'_g(1)=2$.

\textbf{Dérivabilité à droite :} pour $x>1$,
\[ \dfrac{f(x)-f(1)}{x-1}=\dfrac{(2x-1)-1}{x-1}=\dfrac{2(x-1)}{x-1}=2\xrightarrow[x\to 1^+]{}2. \]
Donc $f$ est dérivable à droite en $1$ et $f'_d(1)=2$.
\item Comme $f'_g(1)=f'_d(1)=2$, $f$ \textbf{est dérivable en $1$} et $f'(1)=2$. Graphiquement, la courbe admet au point $(1\,;\,1)$ une unique tangente, de coefficient directeur $2$ : le raccordement entre la parabole et la demi-droite se fait \emph{sans point anguleux} (raccordement « lisse »).
\end{enumerate}
\textbf{Points de vigilance :} la continuité ne suffit pas à conclure à la dérivabilité ; il faut comparer séparément les deux taux d'accroissement. Ici, les deux demi-tangentes ont le même coefficient directeur, donc il n'y a pas de point anguleux.
\end{corrige}

\begin{corrige}[Exercice 6 — Opérations sur les dérivées]
\begin{enumerate}
\item $f(x)=(3x^2-5x)(x+2)$ : produit de deux polynômes, définie et dérivable sur $\mathbb{R}$. Avec $(uv)'=u'v+uv'$, $u=3x^2-5x$, $u'=6x-5$, $v=x+2$, $v'=1$ :
\begin{align*}
f'(x)&=(6x-5)(x+2)+(3x^2-5x)\times 1\\
&=6x^2+12x-5x-10+3x^2-5x=9x^2+2x-10.
\end{align*}
(Vérification par développement : $f(x)=3x^3+x^2-10x$, d'où $f'(x)=9x^2+2x-10$.)
\item $g(x)=\dfrac{x-1}{2x+3}$ : définie si $2x+3\neq 0$, donc $D_g=\mathbb{R}\setminus\left\{-\dfrac{3}{2}\right\}$, et dérivable sur ce même ensemble. Avec $\left(\dfrac{u}{v}\right)'=\dfrac{u'v-uv'}{v^2}$ :
\[ g'(x)=\dfrac{1\times(2x+3)-(x-1)\times 2}{(2x+3)^2}=\dfrac{2x+3-2x+2}{(2x+3)^2}=\dfrac{5}{(2x+3)^2}. \]
\item $h(x)=\sqrt{4x-8}$ : définie si $4x-8\geq 0$, c'est-à-dire $D_h=[2\,;\,+\infty[$ ; dérivable sur $]2\,;\,+\infty[$ (la racine n'est pas dérivable en $2$). Avec $(\sqrt{u})'=\dfrac{u'}{2\sqrt{u}}$ et $u'=4$ :
\[ h'(x)=\dfrac{4}{2\sqrt{4x-8}}=\dfrac{2}{\sqrt{4x-8}}. \]
\end{enumerate}
\textbf{Points de vigilance :} toujours annoncer l'ensemble de définition \emph{avant} de dériver ; dans la formule du quotient, c'est $u'v-uv'$ au numérateur (attention à l'ordre) et $v^2$ au dénominateur.
\end{corrige}

\begin{corrige}[Exercice 7 — Tangente et position relative]
\begin{enumerate}
\item $f(-1)=(-1)^3-2\times(-1)=-1+2=1$. $f'(x)=3x^2-2$, donc $f'(-1)=3-2=1$.
\item Tangente au point d'abscisse $-1$ :
\[ y=f'(-1)(x-(-1))+f(-1)=1\times(x+1)+1=x+2. \]
Donc $T:\ y=x+2$.
\item On étudie la différence :
\[ f(x)-(x+2)=x^3-2x-x-2=x^3-3x-2. \]
Comme $-1$ est racine ($(-1)^3-3(-1)-2=-1+3-2=0$), on factorise par $(x+1)$ :
\[ x^3-3x-2=(x+1)(x^2-x-2)=(x+1)(x+1)(x-2)=(x+1)^2(x-2). \]
Au voisinage de $-1$, on a $x<2$, donc $x-2<0$, et $(x+1)^2\geq 0$. Ainsi $f(x)-(x+2)\leq 0$ au voisinage de $-1$, avec égalité seulement en $x=-1$.

\textbf{Conclusion :} au voisinage de $-1$, la courbe $\mathcal{C}_f$ est \textbf{en dessous} de sa tangente $T$ ; elle touche $T$ au point $(-1\,;\,1)$.
\end{enumerate}
\textbf{Points de vigilance :} pour étudier une position relative, on étudie le signe de la \emph{différence} $f(x)-y_T$ ; la factorisation par $(x-a)$ est souvent possible car $a$ est racine de la différence.
\end{corrige}

\begin{corrige}[Exercice 8 — Approximation affine sans calculatrice]
\begin{enumerate}
\item Soit $f(x)=\sqrt{x}$ ; $f'(x)=\dfrac{1}{2\sqrt{x}}$, donc $f(16)=4$ et $f'(16)=\dfrac{1}{8}=0{,}125$. Approximation affine au voisinage de $16$ :
\[ \sqrt{16{,}1}\approx f(16)+f'(16)\times 0{,}1=4+\dfrac{0{,}1}{8}=4{,}0125. \]
\item Soit $g(x)=\dfrac{1}{x^2}$ ; $g'(x)=-\dfrac{2}{x^3}$, donc $g(1)=1$ et $g'(1)=-2$. Approximation affine au voisinage de $1$ :
\[ \dfrac{1}{(1{,}02)^2}\approx g(1)+g'(1)\times 0{,}02=1-2\times 0{,}02=0{,}96. \]
\item Valeurs à la calculatrice : $\sqrt{16{,}1}\approx 4{,}01248$ et $\dfrac{1}{(1{,}02)^2}=\dfrac{1}{1{,}0404}\approx 0{,}96117$. Les valeurs approchées sont très proches des valeurs exactes (erreur inf\'erieure \`a $0{,}002$) : l'approximation affine est d'autant meilleure que $x$ est proche de $a$.
\end{enumerate}
\textbf{Points de vigilance :} bien identifier $a$ (valeur « ronde » proche) et $h=x-a$ ; l'approximation n'est valable que pour $h$ petit.
\end{corrige}

\begin{corrige}[Exercice 9 — Coût moyen dans une cimenterie]
\begin{enumerate}
\item Pour $x>0$ : $CM(x)=\dfrac{C(x)}{x}=\dfrac{x^2-40x+2000}{x}=\dfrac{x^2}{x}-\dfrac{40x}{x}+\dfrac{2000}{x}=x-40+\dfrac{2000}{x}$.
\item $CM$ est dérivable sur $]0\,;\,+\infty[$ et
\[ CM'(x)=1-\dfrac{2000}{x^2}=\dfrac{x^2-2000}{x^2}. \]
\item Sur $]0\,;\,+\infty[$, $x^2>0$, donc $CM'(x)$ a le signe de $x^2-2000=(x-\sqrt{2000})(x+\sqrt{2000})$. Comme $x>0$, $x+\sqrt{2000}>0$, donc $CM'(x)$ a le signe de $x-\sqrt{2000}$ : $CM'(x)<0$ sur $]0\,;\,\sqrt{2000}[$, $CM'(\sqrt{2000})=0$, $CM'(x)>0$ sur $]\sqrt{2000}\,;\,+\infty[$.

\begin{center}
\begin{tabular}{|c|ccccc|}\hline
$x$ & $0$ & & $\sqrt{2000}$ & & $+\infty$ \\ \hline
$CM'(x)$ & & $-$ & $0$ & $+$ & \\ \hline
$CM$ & & $\searrow$ & $2\sqrt{2000}-40$ & $\nearrow$ & \\ \hline
\end{tabular}
\end{center}
\item Le coût moyen est minimal pour $x=\sqrt{2000}\approx 44{,}72$, soit environ \textbf{45 sacs par jour} (production entière). Le coût moyen minimal est :
\[ CM(\sqrt{2000})=2\sqrt{2000}-40\approx 2\times 44{,}72-40=49{,}44 \text{ milliers de FCFA}, \]
soit environ $\mathbf{49\,440}$ \textbf{FCFA par sac}.
\item Coût total correspondant : $C(44{,}72)=44{,}72^2-40\times 44{,}72+2000=1999{,}88-1788{,}80+2000\approx 2211{,}08$ milliers de FCFA, soit environ $\mathbf{2\,211\,078}$ \textbf{FCFA} (on vérifie : $C(x)=x\times CM(x)\approx 44{,}72\times 49{,}44\approx 2211$ milliers).
\end{enumerate}
\textbf{Barème indicatif (sur 5) :} 0,5 pt (1) ; 1 pt (2) ; 1,5 pt (3) ; 1 pt (4) ; 1 pt (5).

\textbf{Points de vigilance :} justifier la dérivabilité ($x>0$) ; le signe du quotient se ramène au signe du numérateur car $x^2>0$ ; penser à convertir les milliers de FCFA en FCFA ; comme $x$ est un nombre de sacs, commenter l'arrondi à l'entier.
\end{corrige}

\begin{corrige}[Exercice 10 — Mouvement d'un mobile]
\begin{enumerate}
\item La vitesse est le taux de variation instantané de la position, c'est-à-dire la dérivée de $x$ par rapport au temps : $v(t)=x'(t)$. Comme $x$ est un polynôme, dérivable sur $[0\,;\,+\infty[$ :
\[ v(t)=x'(t)=3t^2-12t+9. \]
\item $v(t)=0\iff 3t^2-12t+9=0\iff t^2-4t+3=0\iff (t-1)(t-3)=0$. La vitesse s'annule aux instants $t=1$ s et $t=3$ s.
\item $v(t)=3(t-1)(t-3)$ est un trinôme du second degré ($a=3>0$), positif à l'extérieur des racines, négatif entre les racines :

\begin{center}
\begin{tabular}{|c|ccccccc|}\hline
$t$ & $0$ & & $1$ & & $3$ & & $+\infty$ \\ \hline
$v(t)$ & $9$ & $+$ & $0$ & $-$ & $0$ & $+$ & \\ \hline
\end{tabular}
\end{center}
\item Le mobile \textbf{avance} ($v>0$) pendant $[0\,;\,1[$ et $]3\,;\,+\infty[$ ; il \textbf{recule} ($v<0$) pendant $]1\,;\,3[$.
\item $x(1)=1-6+9=4$ m et $x(3)=27-54+27=0$ m. Interprétation : le mobile part de $x(0)=0$, avance jusqu'à la position $4$ m à $t=1$ s, puis fait demi-tour et revient à l'origine à $t=3$ s, avant de repartir dans le sens positif. Aux instants $t=1$ et $t=3$, la vitesse nulle marque les \textbf{changements de sens} du mouvement.
\end{enumerate}
\textbf{Barème indicatif (sur 5) :} 1 pt (1) ; 1 pt (2) ; 1 pt (3) ; 1 pt (4) ; 1 pt (5).

\textbf{Points de vigilance :} citer que la vitesse est la dérivée de la position (lien dérivation/grandeur physique) ; factoriser $v$ pour étudier le signe ; distinguer vitesse nulle et position nulle.
\end{corrige}

\begin{corrige}[Exercice 11 — Étude complète d'une fonction]
\begin{enumerate}
\item $f$ est un polynôme, dérivable sur $[0\,;\,4]$ : $f'(x)=3x^2-12x+9$. Vérification : $3(x-1)(x-3)=3(x^2-4x+3)=3x^2-12x+9=f'(x)$.
\item $f'(x)=3(x-1)(x-3)$ : trinôme avec $a=3>0$, racines $1$ et $3$. Donc $f'(x)>0$ sur $[0\,;\,1[$ et $]3\,;\,4]$ ; $f'(x)<0$ sur $]1\,;\,3[$ ; $f'(1)=f'(3)=0$.
\item Valeurs : $f(0)=1$ ; $f(1)=1-6+9+1=5$ ; $f(3)=27-54+27+1=1$ ; $f(4)=64-96+36+1=5$.

\begin{center}
\begin{tabular}{|c|ccccccc|}\hline
$x$ & $0$ & & $1$ & & $3$ & & $4$ \\ \hline
$f'(x)$ & $9$ & $+$ & $0$ & $-$ & $0$ & $+$ & $9$ \\ \hline
$f$ & $1$ & $\nearrow$ & $5$ & $\searrow$ & $1$ & $\nearrow$ & $5$ \\ \hline
\end{tabular}
\end{center}
\item $f$ admet un \textbf{maximum local} égal à $5$ en $x=1$ et un \textbf{minimum local} égal à $1$ en $x=3$. Sur $[0\,;\,4]$ tout entier, le maximum est $5$ (atteint en $x=1$ et en $x=4$) et le minimum est $1$ (atteint en $x=0$ et en $x=3$).
\item Tangente en $0$ : $f(0)=1$, $f'(0)=9$, donc $T:\ y=9(x-0)+1$, c'est-à-dire $\boxed{y=9x+1}$.
\item Tracé :

\begin{popfigure}
\begin{center}
\begin{tikzpicture}
\begin{axis}[width=0.85\linewidth, height=7cm, axis lines=middle, xlabel={$x$}, ylabel={$y$}, xmin=-0.4, xmax=4.4, ymin=-0.5, ymax=6.5, xtick={0,1,2,3,4}, ytick={1,2,3,4,5}, grid=both, grid style={popline!40}, samples=150]
\addplot[popcurve, domain=0:4]{x^3-6*x^2+9*x+1};
\addplot[poporange, thick, domain=-0.1:0.7]{9*x+1};
\addplot[only marks, mark=*, popdark] coordinates {(0,1) (1,5) (3,1) (4,5)};
\node[popblue, anchor=west] at (axis cs:3.1,4.2) {$\mathcal{C}_f$};
\node[poporange, anchor=south west] at (axis cs:0.35,4.4) {$T$};
\end{axis}
\end{tikzpicture}
\end{center}
\legende{Courbe de $f(x)=x^3-6x^2+9x+1$ sur $[0;4]$ et tangente au point d'abscisse $0$.}
\end{popfigure}
\end{enumerate}
\textbf{Barème indicatif (sur 6) :} 1 pt (1) ; 1 pt (2) ; 1 pt (3) ; 1 pt (4) ; 1 pt (5) ; 1 pt (6).

\textbf{Points de vigilance :} annoncer la dérivabilité avant de calculer $f'$ ; le tableau de variations doit comporter les images exactes des bornes et des points critiques ; distinguer extremums \emph{locaux} et extremums \emph{sur l'intervalle} $[0;4]$ ; pour le tracé, respecter les unités imposées et placer les tangentes horizontales en $x=1$ et $x=3$.
\end{corrige}

\newpage

\coursec{Fiche bilan}

\begin{popfigure}
\begin{center}
\begin{tikzpicture}[mindmap, grow cyclic, every node/.style={concept}, concept color=popblue!60, text=white, level 1 concept/.append style={sibling angle=60, font=\small}, level 2 concept/.append style={font=\scriptsize}]
\node[concept color=popdark, font=\bfseries] {Dérivation}
  child[concept color=popblue] { node {Nombre dérivé}
    child { node[concept color=popblue!70] {$f'(a)=\lim\frac{f(a+h)-f(a)}{h}$} }
  }
  child[concept color=popgreen] { node {Fonction dérivée}
    child { node[concept color=popgreen!70] {Dérivées usuelles} }
  }
  child[concept color=poporange] { node {Opérations}
    child { node[concept color=poporange!70] {$(u+v)'$, $(uv)'$, $\left(\frac{u}{v}\right)'$} }
    child { node[concept color=poporange!70] {$(f(ax+b))'=a f'(ax+b)$} }
  }
  child[concept color=poppurple] { node {Tangente}
    child { node[concept color=poppurple!70] {$y=f'(a)(x-a)+f(a)$} }
    child { node[concept color=poppurple!70] {$f(a+h)\approx f(a)+hf'(a)$} }
  }
  child[concept color=poppink] { node {Sens de variation}
    child { node[concept color=poppink!70] {$f'>0 \Rightarrow f$ croissante} }
    child { node[concept color=poppink!70] {$f'<0 \Rightarrow f$ décroissante} }
  }
  child[concept color=popgold] { node {Applications}
    child { node[concept color=popgold!70] {Coûts, vitesses, optimisation} }
  };
\end{tikzpicture}
\end{center}
\legende{Carte mentale de la leçon : les six idées clés de la dérivation.}
\end{popfigure}

\begin{bilan}
\textbf{1. Définitions clés}
\begin{itemize}
  \item \cle{Nombre dérivé} : $f$ est dérivable en $a$ si le taux d'accroissement admet une limite finie :
  \[ f'(a)=\lim_{h\to 0}\frac{f(a+h)-f(a)}{h}. \]
  \item \cle{Dérivabilité à gauche / à droite} : on note $f'_g(a)$ et $f'_d(a)$ les limites du taux pour $h<0$ et $h>0$. $f$ est dérivable en $a$ ssi $f'_g(a)=f'_d(a)$.
  \item \cle{Fonction dérivée} : si $f$ est dérivable en tout point d'un intervalle $I$, la fonction $f' : x\mapsto f'(x)$ est la fonction dérivée de $f$ sur $I$.
  \item \cle{Approximation affine} : pour $h$ proche de $0$, $f(a+h)\approx f(a)+h\,f'(a)$.
\end{itemize}

\textbf{2. Formules à connaître par cœur}
\begin{center}
\begin{tabularx}{\linewidth}{|l|X|X|}\hline
\rowcolor{popblueL} \textbf{Fonction $f$} & \textbf{Dérivée $f'$} & \textbf{Conditions} \\ \hline
$k$ (constante) & $0$ & sur $\mathbb{R}$ \\ \hline
$x$ & $1$ & sur $\mathbb{R}$ \\ \hline
$x^n$ ($n\in\mathbb{N}^*$) & $n x^{n-1}$ & sur $\mathbb{R}$ \\ \hline
$\dfrac{1}{x}$ & $-\dfrac{1}{x^2}$ & $x\neq 0$ \\ \hline
$\sqrt{x}$ & $\dfrac{1}{2\sqrt{x}}$ & $x>0$ \\ \hline
\end{tabularx}
\end{center}

\begin{center}
\begin{tabularx}{\linewidth}{|l|X|}\hline
\rowcolor{popgreenL} \textbf{Opération} & \textbf{Dérivée} \\ \hline
Somme & $(u+v)'=u'+v'$ \\ \hline
Produit par un réel & $(ku)'=k\,u'$ \\ \hline
Produit & $(uv)'=u'v+uv'$ \\ \hline
Quotient & $\left(\dfrac{u}{v}\right)'=\dfrac{u'v-uv'}{v^2}$, \quad $v\neq 0$ \\ \hline
Composée affine & $\bigl(f(ax+b)\bigr)'=a\,f'(ax+b)$ \\ \hline
\end{tabularx}
\end{center}

\textbf{3. Théorèmes fondamentaux}
\begin{itemize}
  \item \cle{Tangente} en $A\bigl(a\,;\,f(a)\bigr)$ : $\;y=f'(a)(x-a)+f(a)$.
  \item \cle{Sens de variation} sur un intervalle $I$ :
  \begin{itemize}
    \item si $f'(x)>0$ sur $I$, alors $f$ est strictement croissante sur $I$ ;
    \item si $f'(x)<0$ sur $I$, alors $f$ est strictement décroissante sur $I$ ;
    \item si $f'(x)=0$ sur $I$, alors $f$ est constante sur $I$.
  \end{itemize}
\end{itemize}

\textbf{4. Méthodes express}
\begin{itemize}
  \item \textbf{Calculer $f'(a)$} : former le taux $\frac{f(a+h)-f(a)}{h}$, simplifier, faire tendre $h$ vers $0$.
  \item \textbf{Dresser un tableau de variations} : (1) calculer $f'(x)$ ; (2) résoudre $f'(x)=0$ ; (3) étudier le signe de $f'$ ; (4) conclure sur les variations.
  \item \textbf{Équation de tangente} : calculer $f(a)$ et $f'(a)$, remplacer dans $y=f'(a)(x-a)+f(a)$.
  \item \textbf{Extremum local} : chercher où $f'$ s'annule \emph{en changeant de signe}.
\end{itemize}
\end{bilan}

\begin{aretenir}[Checklist avant le Probatoire]
\begin{itemize}
  \item[$\square$] Je sais écrire la définition de $f'(a)$ comme limite du taux d'accroissement.
  \item[$\square$] Je sais calculer un nombre dérivé par la définition (ex. $f(x)=x^2$ en $a=3$).
  \item[$\square$] Je sais vérifier la dérivabilité en un point en comparant $f'_g(a)$ et $f'_d(a)$.
  \item[$\square$] Je connais par cœur les dérivées de $x^n$, $\frac{1}{x}$ et $\sqrt{x}$.
  \item[$\square$] Je maîtrise $(uv)'=u'v+uv'$ et $\left(\frac{u}{v}\right)'=\frac{u'v-uv'}{v^2}$ sans les confondre.
  \item[$\square$] Je n'oublie pas le facteur $a$ dans $\bigl(f(ax+b)\bigr)'=a\,f'(ax+b)$.
  \item[$\square$] Je sais écrire l'équation de la tangente $y=f'(a)(x-a)+f(a)$ et la réduire.
  \item[$\square$] Je sais utiliser l'approximation affine $f(a+h)\approx f(a)+h\,f'(a)$ pour estimer une valeur.
  \item[$\square$] Je sais déduire le sens de variation du signe de $f'$ et dresser un tableau de variations complet.
  \item[$\square$] Je sais repérer un extremum : $f'$ s'annule \textbf{et} change de signe.
  \item[$\square$] Je vérifie toujours les conditions d'existence ($v\neq 0$, $x>0$ pour $\sqrt{x}$).
  \item[$\square$] Je rédige mes justifications en phrases complètes avec les hypothèses du théorème.
\end{itemize}
\end{aretenir}

\findecours

\end{document}
